\section{NN-A001: exact arithmetic label and a certified correction}
\label{sec:NN-A001}

\paragraph{Inputs.}
A supplies RH.  The analytic inputs are the Riemann--von Mangoldt identity
$N(t)=\theta(t)/\pi+1+S(t)$ with midpoint convention at a zero, and the
Hurwitz derivative identity DLMF 25.11.17.  The supplied consult
\path{OURS_CONSULT_ad855737_ZETA_LABELS.md}, Q1, Proposition 1, already
displays the correction term in its proof, but omits it from its claimed
equivalence.  We retain that proof's term and independently certify a
counterexample to the simplified assertion.

\begin{theorem}[NN-A001]
Let $\rho_n=1/2+i\gamma_n$ be a simple positive-ordinate zero, indexed with
multiplicity among all positive-ordinate zeros, and set
\[
 Z(t)=e^{i\theta(t)}\zeta(1/2+it),\quad
 \phi_n=\arg\zeta(3/2+i\gamma_n),\quad
 W_n=S(\gamma_n)-\phi_n/\pi.
\]
Choose the argument on the absolutely convergent Euler-product branch.
For the deformation $\zeta(s,1+t)$ put
$v_n=\rho_n\zeta(\rho_n+1)/\zeta'(\rho_n)$ and
$\delta_n=\arctan(1/(2\gamma_n))$.  Then
\[
 \operatorname{Re}v_n=
 \frac{|\zeta(\rho_n+1)|}{|Z'(\gamma_n)|}
 \left(\gamma_n\sin\pi W_n+\frac12\cos\pi W_n\right).       \tag{A001.1}
\]
Thus the exact right-mover criterion is
$\sin(\pi W_n+\delta_n)>0$, equivalently
$W_n+\delta_n/\pi\bmod2\in(0,1)$.  Tangency is exactly equality to an
integer modulo one after this correction.  The uncorrected criterion in
the prompt is false: the indexed zero $n=1169$ is a certified right-mover
although $W_{1169}<0$ and $\sin(\pi W_{1169})<0$.
\end{theorem}

\begin{proof}
Implicit differentiation using DLMF 25.11.17 gives the stated $v_n$.
Differentiating $Z(t)$ at a zero gives
$\zeta'(\rho_n)=-i e^{-i\theta(\gamma_n)}Z'(\gamma_n)$.
Since $Z(0)<0$, the sign of $Z'$ at this simple zero is
$(-1)^{n+1}$; preceding zeros of even multiplicity change neither this
sign nor the parity of the total zero count.  The midpoint convention
gives $\theta(\gamma_n)/\pi=n-3/2-S(\gamma_n)$.  Hence
\[
 \cos(\theta+\phi)=(-1)^n\sin\pi W_n,\qquad
 \sin(\theta+\phi)=(-1)^n\cos\pi W_n.
\]
Taking the real part of
$i\rho_n e^{i\theta}\zeta(\rho_n+1)/Z'(\gamma_n)$ proves (A001.1).
The positive factor $\sqrt{\gamma_n^2+1/4}$ converts its last parenthesis
to $\sqrt{\gamma_n^2+1/4}\sin(\pi W_n+\delta_n)$.

The included program \path{check_label_1169.py} calls the indexed
FLINT/Arb zero isolator and evaluates the derivative through a complex
ball power series, at 60 decimal digits.  It certifies
\[
 \gamma_{1169}=1613.0788554510274235\ldots,\quad
 W_{1169}=-0.0000606299724187599509\ldots,
\]
\[
 \delta_{1169}/\pi=0.0000986653178528222938\ldots,
 \qquad \operatorname{Re}v_{1169}=0.025925651967163814\ldots>0.
\]
The saved ball intervals in \path{label_certificate.json} exclude zero
for both sine signs and for the velocity's real part.  The program asserts
the relevant inequalities using ball comparisons, rather than midpoint
rounding.  This supplies the counterexample claimed in the theorem.
\end{proof}

\paragraph{Interpretation and geometry.}
The formula is a real normal component of a tangent vector to a moving
analytic Cartier divisor.  Its sign is a Boolean-valued statistic where
the normal component is nonzero; it is not a homomorphism of the compact
circle into a two-element group.  Indeed $\operatorname{sgn}\sin\pi x$ is
a discontinuous square wave, and even as a set map is not multiplicative
under addition: two small positive $x,y$ give $+,+,+$, while suitable
$x,y$ give $+,+,-$.  In particular the consult's description as a sign
character of a quotient circle is not a representation-theoretic identity.
The weight-zero Boolean skeleton in the split-zero construction contains
no value of $\zeta'(\rho)$ or of this normal vector.  This corrected label
therefore supplies no SAT, parity, or multiplicity-positivity algorithm.

\section{NN-A002: the exact statistical gate and why a CLT does not close it}
\label{sec:NN-A002}

\begin{theorem}[NN-A002]
A central limit theorem for $S(\gamma_n)$ after division by a standard
deviation tending to infinity, even together with boundedness of
$\phi_n$, does not imply that $W_n\bmod2$ equidistributes, nor that half
the labels in NN-A001 are positive.

A sufficient arithmetic gate for the half-proportion, on the chosen
simple-zero subsequence, is the following pair of assertions.  For every
odd integer $k$, the empirical means of $e^{i\pi kW_n}$ tend to zero;
and
\[
 \lim_{\epsilon\downarrow0}\limsup_{T\to\infty}
 \frac{\#\{n:\gamma_n\leq T,\ n\text{ simple},\
                 \operatorname{dist}(W_n,\mathbb Z)<\epsilon\}}
      {\#\{n:\gamma_n\leq T,\ n\text{ simple}\}}=0.       \tag{A002.1}
\]
All nonzero Fourier modes tending to zero is the stronger usual Weyl
equidistribution gate.  The required odd modes have the exact form
\[
 e^{i\pi kW_n}=e^{-3\pi i k/2}(-1)^n
          e^{-ik\theta(\gamma_n)}e^{-ik\phi_n}\qquad(k\text{ odd}).
                                                               \tag{A002.2}
\]
\end{theorem}

\begin{proof}
Let $G$ be standard normal, let $\sigma_j\to\infty$, and set
$X_j=2\lfloor\sigma_jG/2\rfloor+1/4$.
Then $|X_j-\sigma_jG|\leq7/4$, so $X_j/\sigma_j$ converges almost surely
and in every fixed $L^p$ to $G$.  In particular the normalized central
limit law and all its fixed normalized moment consequences hold.
Nevertheless $X_j\bmod2=1/4$ identically, so every square-wave label is
positive.  Equivalently, an exactly Gaussian variable $\sigma_jG$ can
be altered by a uniformly bounded, correlated correction to have this
completely nonuniform residue class.  This proves the failure of the
implication invoked in the supplied consult, Q2, Proposition 4.

For the sufficient gate, approximate the square wave
$f(x)=\operatorname{sgn}\sin\pi x$ by its Fejer trigonometric means.
It has only odd Fourier modes, because $f(x+1)=-f(x)$, and zero constant
term.  On any closed set a positive distance from $\mathbb Z$, these
means converge uniformly to $f$, and they are bounded in absolute value
by one.  For fixed truncation the empirical mean tends to zero by the
odd-mode hypothesis.  Equation (A002.1) controls the omitted boundary
neighborhood; then let the truncation tend to infinity and that
neighborhood shrink.  Thus the empirical mean of $f(W_n)$ tends to zero.
Since $\delta_n\to0$, (A002.1) also shows that replacing $f(W_n)$ by
$f(W_n+\delta_n/\pi)$ changes an asymptotically negligible proportion.
This proves the half-proportion for the exact labels.  Finally substitute
$W_n=n-3/2-\theta(\gamma_n)/\pi-\phi_n/\pi$ to get (A002.2).
\end{proof}

\paragraph{A primary CLT locator and its frequency range.}
The previously missing source can be made precise: Korolev,
\emph{On the Gram's Law in the Theory of Riemann Zeta Function},
arXiv:1011.3997, Theorem~14, printed pp.33--34, gives a Gaussian
law for $S(\gamma_n\pm0)/\sigma_N$, where
$\sigma_N=\sqrt{L/2}/\pi$ and $L=\log\log N$, on the stated
index blocks. Its proof invokes Theorem~10, pp.21--22. There the
characteristic-function argument is truncated at normalized frequency
$\lambda\asymp\sqrt{\log L}$. In the original $S$ coordinate
this reaches only $\lambda/\sigma_N
=O(\sqrt{\log L/L})\to0$, not the fixed nonzero frequencies
$\pi k$ required by NN-A002. The theorem also does not insert the
correlated exterior phase. Thus the concrete classical CLT, and not
merely an abstract CLT, leaves the identified fixed-frequency gate open.

\paragraph{Located progress beyond the invalid CLT route.}
Arias de Reyna, \emph{On the distribution (mod 1) of the normalized zeros
of the Riemann zeta-function}, arXiv:1403.7126, Theorem 3.1 and Corollary
3.2, prove cancellation of
$N(T)^{-1}\sum_{0<\gamma\leq T}e^{2i\kappa\theta(\gamma)}$ for
$0<\kappa<6/5$ once A identifies the ordinates with the real arguments
in that theorem.  The theorem itself is formulated without RH using
complex ordinates.  Its statement does not supply the alternating index
factor and exterior phase in (A002.2), does not select the simple-zero
subsequence, and does not establish all the required modes.  These are
specific additional gates.  The weighted first moment of velocities is
also insufficient: cancellation of a sum of real velocities controls
their sizes with signs, whereas the question here counts signs without
those weights.

\paragraph{Dictionary and complexity impact.}
The Fourier gate is a statement about characters $e^{i\pi kx}$ of a
compact real torus and a discontinuous observable on it.  The sample
points are analytic divisor points; no finite-type scheme or Frobenius
action on these phase values has been constructed.  This statement
neither establishes a nonstandard Weil purity theorem nor changes the
complexity of support, parity, or satisfiability.


\section{NN-A003: a CUE theorem with the exterior phase retained}
\label{sec:NN-A003}

\paragraph{Inputs and status.}
This result uses no instance of A, M, E, or a GCT hypothesis.  Its inputs are
the Haar eigenangle density for $U(N)$ (Weyl integration), the elementary
determinantal integral identity proved below, and elementary Fourier analysis.
The Haar character orthogonality used in the final phase estimate also appears
in the supplied paper 12, Theorem \texttt{thm:first-deriv-mean}.  This is a
proved form within the specified CUE model, not a transfer theorem for zeta
zeros.  In particular the arithmetic phase is not discarded on the grounds
that it is bounded.

\begin{theorem}[NN-A003: marked CUE labels and their correlations]
Let $U_N$ be Haar distributed in $U(N)$, let
$0\leq\theta_1<\cdots<\theta_N<2\pi$ be its eigenangles, and choose an index
$I$ uniformly from $\{1,\ldots,N\}$, independently after sampling $U_N$.
Fix $0\leq r<1$ and $a\in\mathbb R$.  Define the continuous exterior phase
\[
 \phi_{N,r}(x)=\operatorname{Im}\sum_{j=1}^N
       \log(1-r e^{i(x-\theta_j)}),
 \qquad
 W_{N,i}=i-\frac{N\theta_i}{2\pi}-a-
                 \frac{\phi_{N,r}(\theta_i)}{\pi}.
\]
The logarithm is its power-series branch in $r$, so there is no branch
ambiguity.  Set $\eta_{N,i}=\operatorname{sgn}\sin(\pi W_{N,i})$; a convention
at zero has no effect on the limits.  For fixed $j\geq1$, let
\[
 D_{j,N}=\frac{N}{2\pi}(\theta_{I+j}-\theta_I),
 \qquad \theta_{i+N}=\theta_i+2\pi.
\]
Then $(W_{N,I}\bmod2,D_{j,N})$ converges in distribution to $(V,D_j)$,
where $V$ is uniform on $\mathbb R/2\mathbb Z$, independent of $D_j$, and
$D_j$ is the distance to the $j$th point to the right of zero in the reduced
Palm sine process.  The latter process has kernel
\[
 K^{0}(x,y)=\frac{\sin\pi(x-y)}{\pi(x-y)}-
      \frac{\sin\pi x}{\pi x}\frac{\sin\pi y}{\pi y}.
\]
Consequently
\[
 \lim_{N\to\infty}\mathbb P(\eta_{N,I}=1)=\frac12,
 \qquad
 \lim_{N\to\infty}\mathbb E[\eta_{N,I}\eta_{N,I+j}]
       =\mathbb E[1-2d_2(j-D_j)],                 \tag{A003.1}
\]
where $d_2(x)=\min_{m\in\mathbb Z}|x-2m|\in[0,1]$.
In particular the adjacent agreement probability tends to
$1-\mathbb E d_2(1-D_1)$.  All assertions remain valid for $r=r_N$ in a fixed
compact subinterval of $[0,1)$, provided a subsequence is used if a limiting
phase law itself is requested; the displayed label and gap limits are
independent of that phase law.
\end{theorem}

\begin{proof}
Normalize circle measure to $d\theta/(2\pi)$ and let $P_N$ be the orthogonal
projection onto the span of $1,e^{i\theta},\ldots,e^{i(N-1)\theta}$.
The squared Vandermonde density is the determinant of the kernel of $P_N$.
Expanding its two determinants and integrating each coordinate shows directly
that, for a bounded multiplier $u$,
\[
 \mathbb E\prod_i u(\theta_i)
       =\det(P_NuP_N|_{\operatorname{ran}P_N}).       \tag{A003.2}
\]
Indeed each term in the determinant expansion is a product of the matrix
entries $\int u(\theta)e^{i(k-l)\theta}d\theta/(2\pi)$, and the permutation
signs are precisely the determinant signs.  Conditioning one uniformly
marked point to be $x$ gives the same formula with the rank $N-1$ projection
\[
 Q_{N,x}=P_N-|v_x\rangle\langle v_x|,
 \qquad v_x(\theta)=P_N(\theta,x)/\sqrt N.
\]
This conditional kernel follows either by the Schur complement of the
determinant with the marked row and column fixed, or by expanding that row
and column in (A003.2).  The marked point itself is uniform on the circle.

For any finite-rank projection $Q$ and multiplier $u$ with $|u|=1$, put
$A=QuQ|_{\operatorname{ran}Q}$.  Then
\[
 |\det A|^2=\det\bigl(I_Q-Qu^*(1-Q)uQ\bigr)
       \leq \exp\bigl(-\|(1-Q)uQ\|_{\rm HS}^2\bigr),       \tag{A003.3}
\]
because the eigenvalues of the positive contraction subtracted from $I_Q$
are in $[0,1]$, and $1-y\leq e^{-y}$.  For $Q=P_N$ its exponent equals
\[
 F_N(u):=\|(1-P_N)uP_N\|_{\rm HS}^2
       =\sum_{n\in\mathbb Z}\min(N,|n|)|\widehat u(n)|^2.   \tag{A003.4}
\]
This follows by applying $u$ to each Fourier basis vector and counting how
many translates by $n$ leave $\{0,\ldots,N-1\}$.  Removing one vector from
the projection changes this leakage by at most one.  To see this write
$P=Q+R$, where $R$ has rank one.  Orthogonality gives
\[
 F(P)-F(Q)=-\|RuQ\|_{\rm HS}^2+
                 \|(1-P)uR\|_{\rm HS}^2\in[-1,1].       \tag{A003.5}
\]

Consider a nonzero Fourier mode $k\in\mathbb Z$ of $W\bmod2$.
When $k$ is even, $e^{i\pi kI}=1$.  The exterior phase at the marked
eigenangle and all cyclic gaps are invariant under simultaneous rotation of
all eigenangles.  Averaging the factor $e^{-ikNx/2}$ over that rotation
therefore makes the Fourier mode exactly zero, also after multiplication
by any bounded function of cyclic gaps.

When $k$ is odd, condition on the marked point $x\in(0,2\pi)$.  Its index
is $1+$ the number of remaining points in $(0,x)$.  Up to a scalar of
modulus one, the conditional Fourier mode is (A003.2) for $Q_{N,x}$ and
\[
 u_x(\theta)=(-1)^{\mathbf1_{(0,x)}(\theta)}
        \exp\{-ik\,\operatorname{Im}\log(1-r e^{i(x-\theta)})\}.
\]
This function is smooth except for two jumps, at $0$ and $x$, each of
absolute size two.  Piecewise integration by parts twice gives
\[
 \widehat u_x(n)=\frac{J_0+J_x e^{-inx}}{2\pi i n}+O_{x,k,r}(n^{-2}),
 \qquad |J_0|=|J_x|=2.
\]
For fixed $x\notin2\pi\mathbb Z$ the cross sum
$\sum_{n\leq N}e^{inx}/n$ is bounded.  Thus (A003.4), including positive
and negative frequencies, gives
\[
 F_N(u_x)\geq\frac4{\pi^2}\log N+O_{x,k,r}(1).
\]
Equations (A003.3)--(A003.5) imply that the conditional odd mode tends to
zero.  Its absolute value is at most one, so dominated convergence over
the marked angle proves marginal equidistribution.  Derivative bounds of
the smooth factor are uniform when $r$ stays bounded away from one.

The same estimate proves independence from local gaps, which does not
follow from marginal equidistribution alone.  Fix $s>0$ and let $J$ be the
arc of length $2\pi s/N$ immediately after $x$.  For $|z|\leq1$ replace
$u_x$ by $w=u_x z^{\mathbf1_J}$.  This is a contraction multiplier, and
\[
 \|(w-u_x)Q_{N,x}\|_{\rm HS}^2
      \leq 4\operatorname{Tr}(Q_{N,x}\mathbf1_JQ_{N,x})\leq4s.
\]
For a contraction $w$, $I_Q-(QwQ)^*(QwQ)$ is the sum of
$Q(1-|w|^2)Q$ and $Qw^*(1-Q)wQ$, both positive.  Hence its determinant
bound has exponent at least $\|(1-Q)wQ\|_{\rm HS}^2$.  By the triangle
inequality this is at least
\[
 \bigl(\|(1-Q)u_xQ\|_{\rm HS}-2\sqrt s\bigr)_+^2,
\]
which still tends to infinity.  It follows uniformly on every circle
$|z|=b<1$ that the mixed conditional generating function of the odd
Fourier mode and $z^{\#J}$ tends to zero.  Cauchy coefficient extraction
then gives the same assertion with $\mathbf1_{\#J=l}$ for every fixed
$l$.  Since $D_{j,N}>s$ means $\#J\leq j-1$, this proves vanishing of
all mixed Fourier modes against these gap events.  The even modes already
vanish by rotation.  The gaps are tight, since their mean is exactly $j$
by cyclic summation of all gaps.

After scaling by $N/(2\pi)$ and removing a harmless phase, the conditional
kernel tends uniformly on each compact interval to $K^0$ displayed in the
theorem.  Determinantal generating functions on such an interval converge:
their Fredholm series converge term by term, and Hadamard's bound for
positive kernels with diagonal at most one bounds the $m$th term by
$ (2s)^m/m!$. The limiting $D_j$ has no atoms: a determinantal process
with continuous intensity has probability zero of a point at any
prescribed location, and $D_j=s$ requires such a point at $s$.
The local count formulas and tightness therefore give convergence of
the scaled $j$th gaps to $D_j$ at every survival threshold. The mixed Fourier
identities just proved characterize the joint limit as the product of
Haar measure modulo two and this gap law.

It remains to compare the phase at neighboring eigenangles.  The convergent
Fourier expansion gives
\[
 \sup_x|\phi_{N,r}'(x)|\leq
          \sum_{m\geq1}r^m|\operatorname{Tr}(U_N^m)|.
\]
Direct Fourier integration of the CUE two-point determinant gives
$\mathbb E|\operatorname{Tr}(U_N^m)|^2=\min(m,N)$; the same identity
follows by substituting $e^{im\theta}$ into (A003.4).  Cauchy--Schwarz
therefore bounds the expectation of this supremum by
$\sum_{m\geq1}r^m\sqrt m<\infty$, uniformly in $N$.
Tightness of $D_{j,N}$ shows that
$\phi_{N,r}(\theta_{I+j})-\phi_{N,r}(\theta_I)$ tends to zero in
probability.  With unwrapped indices and angles the exact difference is
\[
 W_{N,I+j}-W_{N,I}=j-D_{j,N}
       -\frac{\phi_{N,r}(\theta_{I+j})-\phi_{N,r}(\theta_I)}{\pi}.
\]
Finally, for uniform $V\bmod2$, the two square-wave signs at $V$ and
$V+v$ agree on a set of normalized measure $1-d_2(v)$.  Their correlation
is $1-2d_2(v)$.  The discontinuity boundaries have zero limiting Haar
measure, so the joint convergence proves (A003.1).
\end{proof}

\paragraph{What this changes.}
The CUE half-proportion and correlation formulas are justified within an
explicit model containing the correlated exterior phase.  The earlier
nonperiodic expression $1-\mathbb E|1-D_1|$ for adjacent agreement must be
replaced by $1-\mathbb E d_2(1-D_1)$; gaps larger than two cannot be silently
omitted.  To apply this to the arithmetic labels still requires a joint
arithmetic limit theorem for the label phase and local zero gaps.
Montgomery's restricted pair-correlation theorem supplies neither the odd
label modes nor the nearest-neighbor law.  This finding supplies no
algorithm for a Boolean counting problem and no separation of complexity
classes.

\paragraph{Geometric and representation-theoretic dictionary.}
At finite $N$, (A003.2) is an exterior-algebra determinant of a compressed
operator, derived from Weyl integration on the compact real form of
$\mathrm{GL}_N$.  The rank-one Palm operation is removal of the evaluation
line, and the leakage in (A003.3) is a Hilbert--Schmidt norm.  These are
precise Hilbert-space and representation-theoretic constructions; they do
not assert a scheme over $\mathbb F_q$, a finite Weil object, or Frobenius
purity.  The limiting sine projection has infinite rank.  No Grothendieck
archive cote is claimed for this new analytic argument.

\paragraph{Exact scope of the classical arithmetic pair theorem.}
The primary locator is Montgomery, \emph{The pair correlation of zeros
of the zeta function} (1973), Section~1, Theorem and equation~(2),
printed p.181; the Fourier identity (3)--(4) and explicit support
restriction are on p.182. Under RH, for each fixed $\epsilon>0$,
\[
 F(\alpha,T)=(1+o(1))T^{-2|\alpha|}\log T+|\alpha|+o(1)
 \quad (|\alpha|\le1-\epsilon),
\]
with the stated uniformity. The full continuation beyond this range
is presented as conjecture (11)--(12), pp.183--184. This controls
band-limited pair observables, while a neighboring-gap event requires
an empty interval and thus higher count information. The arithmetic
derivative and exterior phase are additional marks. Neither of these
missing pieces is supplied by that restricted pair theorem. This is
existing background, not a repeated consequence from the prior catalogue.


\section{Quantitative empirical balance of the CUE labels}
\label{sec:analytic-cue-empirical-balance}

NN-A003 concerns a uniformly marked eigenangle. The following
statement controls the actual fraction of labels in a single
sample, while retaining the same correlated exterior phase.

\begin{theorem}[NN-A031: empirical half proportions in mean square]
\label{thm:analytic-cue-empirical-balance}
Fix $r_0<1$. For each $N$, use the CUE labels of NN-A003 with
arbitrary $0\le r_N\le r_0$ and arbitrary offset $a_N\in\mathbb R$:
\[
 W_{N,i}=i-\frac{N\theta_i}{2\pi}-a_N
          -\frac1\pi\operatorname{Im}\sum_j
                   \log(1-r_Ne^{i(\theta_i-\theta_j)}),
 \qquad \eta_{N,i}=\operatorname{sgn}\sin(\pi W_{N,i}).
\]
Almost surely none of these labels is zero. There is a constant
$C_{r_0}$, independent of $N,r_N,a_N$, such that for $N\ge3$,
\begin{equation}\label{eq:analytic-cue-empirical-rate}
 \mathbb E\left[\left(\frac1N\sum_{i=1}^N\eta_{N,i}\right)^2\right]
                          \le\frac{C_{r_0}}{\log N}.
\end{equation}
In particular, the observed positive-label fraction $L_N$ satisfies
\[
 \mathbb E[(L_N-1/2)^2]\le\frac{C_{r_0}}{4\log N},\qquad
 \mathbb P(|L_N-1/2|>\varepsilon)
                   \le\frac{C_{r_0}}{4\varepsilon^2\log N}.
\]
Thus $L_N\to1/2$ in $L^2$ and in probability. No optimality of
the logarithmic error bound, almost-sure convergence along all
$N$, or arithmetic transfer is asserted.
\end{theorem}

\begin{proof}
Write $P=P_N$ for the CUE projection in
$L^2(S^1,d\theta/(2\pi))$, and let $I,J$ be two distinct ordered
uniform indices sampled after the matrix. All constants below
are independent of the offset and of $r\in[0,r_0]$.

\emph{An exact uniform marginal modulo one.}
Under simultaneous rotation by $\alpha$, the exterior phase at
a carried marked point is unchanged. Its coordinate changes by
$\alpha$ modulo $2\pi$, and its new index differs from its old
index by an integer. Consequently its $W$ coordinate modulo one
changes by $-N\alpha/(2\pi)$. The marked Haar ensemble is
rotation invariant, so for every nonzero integer $h$,
\[
 \mathbb E e^{2\pi ihW_{N,I}}
 =e^{-ihN\alpha}\mathbb E e^{2\pi ihW_{N,I}}
 \quad\hbox{for every }\alpha.
\]
These expectations vanish. Hence $W_{N,I}\bmod1$ is exactly
uniform for every $N$. In particular a uniformly marked label
is almost surely nonzero, and a finite union gives the same
assertion for all labels in a sample. For any integrable
one-periodic function $g$ we have the exact identity
\begin{equation}\label{eq:analytic-cue-exact-modone}
             \mathbb E g(W_{N,I})=\int_0^1g(u)\,du.
\end{equation}

\emph{Two-Palm odd Fourier modes.}
Condition on the marked eigenangles $x\ne y$. The remaining
points form a projection determinantal process with projection
$Q_{N,x,y}$ obtained by removing the two-dimensional span of
the evaluation vectors at $x,y$ from $\operatorname{ran}P$.
This follows by the Schur complement of their two rows and
columns in the CUE density; the two evaluation vectors are
linearly independent for $N\ge2$ and $x\ne y$.

Let $k,l$ be odd integers. Their index parities combine to the
parity of the number of remaining eigenangles in the arc
between $x$ and $y$ that does not cross the coordinate origin.
For example, when $x<y$, the two counts in $(0,x)$ and $(0,y)$
have sum congruent modulo two to the count in $(x,y)$; the
contribution from the two marked points is a constant sign.
Define
\[
 j_{x,y}(\theta)=(-1)^{\mathbf1_{(\min(x,y),\max(x,y))}(\theta)},
 \qquad g_x(\theta)=\operatorname{Im}\log(1-re^{i(x-\theta)}),
\]
and
\[
 v(\theta)=e^{-ikg_x(\theta)-ilg_y(\theta)},\qquad u=j_{x,y}v.
\]
The self- and cross-contributions of the marked points to the
exterior phases are fixed scalars of modulus one. The determinant
identity in NN-A003 therefore gives
\begin{equation}\label{eq:analytic-cue-two-palm-mode}
 \left|\mathbb E\left[e^{i\pi(kW_{N,I}+lW_{N,J})}
                     \mid x,y\right]\right|
   =\left|\det(Q_{N,x,y}uQ_{N,x,y})\right|.
\end{equation}
The determinant is on the range of $Q_{N,x,y}$.

For a unimodular multiplier $w$, retain the notation
\[
 F_N(w)=\|(1-P)wP\|_{\rm HS}^2
       =\sum_{n\in\mathbb Z}\min(N,|n|)|\widehat w(n)|^2.
\]
Expanding the two off-diagonal blocks of the commutator gives
$\|[P,w]\|_{\rm HS}^2=2F_N(w)$: both blocks have squared norm
$N-\|PwP\|_{\rm HS}^2$. Since
$[P,jv]=[P,j]v+j[P,v]$, the triangle inequality gives
\begin{equation}\label{eq:analytic-cue-smooth-multiplier}
       \sqrt{F_N(jv)}\ge\sqrt{F_N(j)}-\sqrt{F_N(v)}.
\end{equation}
The smooth multiplier has derivative bounded by
\[
 \|v'\|_2\le C_0(|k|+|l|),\qquad C_0=\frac{r_0}{1-r_0}.
\]
By Cauchy--Schwarz and Parseval,
\begin{equation}\label{eq:analytic-cue-smooth-leakage}
 F_N(v)\le\sum_n|n||\widehat v(n)|^2
     \le\|v'\|_2\|v\|_2
     \le C_0(|k|+|l|).
\end{equation}

Let $d(x,y)\in(0,\pi]$ be their circular separation and let
$\ell=|x-y|$. The arc multiplier has the exact coefficients
\[
 |\widehat j_{x,y}(n)|^2
       =\frac{2-2\cos(n\ell)}{\pi^2n^2}\quad(n\ne0).
\]
Consequently
\[
 F_N(j_{x,y})\ge\frac4{\pi^2}
       \left(H_N-\sum_{n=1}^N\frac{\cos(n\ell)}n\right).
\]
Summation by parts bounds the tail of
$\sum_{n\ge1}e^{in\ell}/n$ by
$2/(N|1-e^{i\ell}|)$. Its real infinite sum is
$-\log|1-e^{i\ell}|$. Since
$|1-e^{i\ell}|\ge2d(x,y)/\pi$, for
$d(x,y)\ge N^{-1/2}$ we obtain, uniformly,
\[
 F_N(j_{x,y})\ge\frac2{\pi^2}\log N-O(1)
                         \ge\frac1{\pi^2}\log N
\]
for all sufficiently large $N$.

Put
\[
 c_* =\frac1{64\pi^2(1+C_0)},\qquad
 M=\lfloor c_*\log N\rfloor,
\]
and assume $N$ is sufficiently large that $M\ge2$. If $k,l$ are
odd and $|k|,|l|\le2M$, then
\[
 F_N(v)\le4C_0M\le\frac{\log N}{16\pi^2}
                         \le\frac1{16}F_N(j_{x,y}).
\]
Equation \eqref{eq:analytic-cue-smooth-multiplier} yields
$F_N(u)\ge(9/16)F_N(j_{x,y})\ge(\log N)/(2\pi^2)$.
Removing a rank-one projection changes the leakage by at most
one, as proved in NN-A003; hence the two-Palm leakage is at
least $F_N(u)-2$. The determinant contraction from NN-A003 gives
\[
 |\det(Q_{N,x,y}uQ_{N,x,y})|
                 \le eN^{-\kappa},\qquad
                      \kappa=\frac1{4\pi^2},
\]
whenever $d(x,y)\ge N^{-1/2}$.

The joint density of the ordered marked pair, relative to
$dx\,dy/(2\pi)^2$, is
\[
 \frac{N^2-|P_N(x,y)|^2}{N(N-1)}\le\frac{N}{N-1}\le2.
\]
The excluded region $d(x,y)<N^{-1/2}$ has measure
$O(N^{-1/2})$, and a conditional characteristic function has
absolute value at most one there. Integration therefore proves
the uniform estimate
\begin{equation}\label{eq:analytic-cue-uniform-odd-pair}
 \left|\mathbb E e^{i\pi(kW_{N,I}+lW_{N,J})}\right|
                  \le C N^{-\kappa}
 \quad\left(\begin{array}{c}k,l\text{ odd},\\|k|,|l|\le2M\end{array}\right).
\end{equation}
No independence of the exterior phases has been used.

\emph{A positive finite Fourier approximation to the labels.}
On $\mathbb R/2\mathbb Z$, with normalized measure $du/2$, let
$f(u)=\operatorname{sgn}\sin(\pi u)$ and normalize the positive
Jackson kernel
\[
 J_M(u)=\frac1{Z_M}
       \left|\frac{\sin(M\pi u/2)}{\sin(\pi u/2)}\right|^4,
 \qquad \int_{-1}^1J_M(u)\frac{du}{2}=1.
\]
The quotient has its continuous value $M$ at zero. Elementary
sine bounds give $Z_M\asymp M^3$ and
\[
 J_M(u)\le CM\min\{1,(M|u|)^{-4}\},\qquad
 \int_{-1}^1|u|J_M(u)\frac{du}{2}\le\frac CM.
\]
For the lower bound on $Z_M$, restrict to $|u|\le1/M$ and use
$|\sin(M\pi u/2)|\ge M|u|$ and
$|\sin(\pi u/2)|\le\pi|u|/2$. The upper and moment bounds follow
by splitting at $1/M$.

Set $q_M=f*J_M$. This is a real trigonometric polynomial of
degree at most $2(M-1)$, with $|q_M|\le1$. Because
$f(u+1)=-f(u)$, it has only odd Fourier modes. Moreover,
\begin{equation}\label{eq:analytic-cue-jackson-error}
 \|f-q_M\|_{L^1(\mathbb R/2\mathbb Z)}\le\frac CM,
 \qquad \sum_k|\widehat q_M(k)|\le C\log(M+1).
\end{equation}
The first estimate follows by averaging
$\|f(\cdot-u)-f\|_1\le2|u|$ against $J_M$. For the second use
$|\widehat J_M(k)|\le1$ and
$|\widehat f(k)|=2/(\pi|k|)$ for odd $k$.

The function $|f-q_M|$ is one-periodic. Thus
\eqref{eq:analytic-cue-exact-modone} makes its expected value at
the marked $W$ exactly its circle integral, without an
asymptotic equidistribution error. Since all labels and $q_M$
values have absolute value at most one,
\[
 \left|\mathbb E[\eta_{N,I}\eta_{N,J}]
       -\mathbb E[q_M(W_{N,I})q_M(W_{N,J})]\right|
                                             \le\frac CM.
\]
Using the odd-mode estimate
\eqref{eq:analytic-cue-uniform-odd-pair} in the finite Fourier
expansion gives
\[
 |\mathbb E[\eta_{N,I}\eta_{N,J}]|
       \le\frac CM+C N^{-\kappa}\log^2(M+1).
\]
Finally the diagonal terms of the empirical square contribute
$1/N$, so
\[
 \mathbb E\left[\left(\frac1N\sum_i\eta_{N,i}\right)^2\right]
 \le\frac1N+\frac CM+C N^{-\kappa}\log^2(M+1)
                              =O_{r_0}(1/\log N).
\]
Increasing the constant covers the finitely many smaller $N$.
The assertions about $L_N=(1+N^{-1}\sum_i\eta_{N,i})/2$ follow
from this identity and Markov's inequality.
\end{proof}

\paragraph{Geometric translation.}
Identify $\operatorname{ran}P_N$ with the restrictions to the unit
circle of $H^0(\mathbb P^1,\mathcal O(N-1))$, carrying its Haar
Hermitian norm. At two distinct marked points $z_x,z_y$, polynomial
division gives the exact evaluation sequence
\[
 0\longrightarrow (z-z_x)(z-z_y)\mathbb C[z]_{\le N-3}
 \longrightarrow\mathbb C[z]_{\le N-1}
 \longrightarrow\mathbb C^2\longrightarrow0.
\]
The two-Palm projection is the orthogonal projection onto this
vanishing subspace. Its exterior determinant controls the marked
Fourier statistics; the rank-two correction and the commutator
norm make that control quantitative. The representation-ring
exterior-power formalism is SGA~6, Expos\'e~V, Example~2.2(b) and
Lemma~2.2.1. These actual evaluation spaces and Haar projections
give a probabilistic balance law. They do not provide an
arithmetic comparison map or a Weil purity argument.


\section{A local gap limit and distant-label decorrelation in the CUE limit}

\begin{theorem}[NN-A016: periodic gap displacement and pair decorrelation]
Let $D_j$ be the $j$th positive point of the reduced Palm sine process
at zero, with kernel $K^0$ as in NN-A003. It has a density $f_j$.
For $j\ge2$ put $\sigma_j=\sqrt{\log j}/\pi$ and extend $f_j(s)$
by zero for $s\le0$. Then
\[
 \int_{\mathbb R}\left|
 \sigma_j f_j(j+\sigma_j y)
       -\frac{e^{-y^2/2}}{\sqrt{2\pi}}\right|dy\longrightarrow0.
                                                        \tag{A016.1}
\]
In particular, for every fixed nonzero real $a$,
\[
 \mathbb E e^{ia(D_j-j)}\longrightarrow0,               \tag{A016.2}
\]
and $(D_j-j)\bmod2$ tends to Haar measure on $\mathbb R/2\mathbb Z$.
Consequently the limiting CUE label correlations in NN-A003 satisfy
\[
 \mathbb E[1-2d_2(j-D_j)]\longrightarrow0.
\]
Thus two labels far apart in the limiting indexed model become
independent fair signs. The order of limits is first $N\to\infty$
for fixed $j$ in NN-A003, then $j\to\infty$ here. No uniform
simultaneous-$j,N$ claim or arithmetic transfer is asserted.
\end{theorem}

\begin{proof}
We include the local-limit step, because an ordinary central limit
theorem alone would not imply (A016.2).

Write $k(x)=\sin(\pi x)/(\pi x)$. The sine kernel $K(x,y)=k(x-y)$
is the orthogonal projection in $L^2(\mathbb R)$ onto Fourier
frequencies $[-1/2,1/2]$. Its compression $A_s$ to $(0,s)$ is a
positive trace-class contraction with trace $s$. The determinant
generating function gives
\[
 \mathbb E z^{N(0,s)}=\det(I+(z-1)A_s)
       =\prod_\ell(1-\lambda_\ell+\lambda_\ell z).
\]
Thus this count has the distribution of a sum of independent Bernoulli
variables with parameters the eigenvalues $\lambda_\ell\in[0,1]$.
The product converges because $\sum\lambda_\ell=s<\infty$;
its probability generating function identifies the distribution.
Direct integration gives
\[
 \operatorname{Var}N(0,s)
 =s-2\int_0^s(s-u)k(u)^2du
 =\frac{\log s}{\pi^2}+O(1).                          \tag{A016.3}
\]
Indeed $\int_0^\infty k(u)^2du=1/2$; the remaining tail
$2s\int_s^\infty k(u)^2du$ is bounded, while
$2\int_0^s uk(u)^2du=(\log s)/\pi^2+O(1)$ by
$\sin^2(\pi u)=(1-\cos(2\pi u))/2$.

Conditioning on a fixed point removes its kernel evaluation line
from the projection. Conditioning on two distinct points $0,s$
removes their two-dimensional span. This follows either directly
from the Schur complement of the determinant correlation functions
or from the Gram determinant formula for those evaluation vectors.
Let $B_s$ be the restriction to $(0,s)$ of this two-point reduced
Palm projection. Then $A_s-B_s\ge0$ has trace at most two and rank
at most two. Hence
\[
 \operatorname{Tr}B_s=s+O(1),\qquad
 \operatorname{Tr}(B_s-B_s^2)
       =\frac{\log s}{\pi^2}+O(1).                    \tag{A016.4}
\]
For clarity, if $R=A_s-B_s$, cyclicity of trace gives the variance
difference $\operatorname{Tr}R-\operatorname{Tr}(R(A_s+B_s))$.
Since $0\le A_s+B_s\le2I$, its absolute value is at most
$\operatorname{Tr}R\le2$. The same assertions, with rank one,
hold for the Palm process at zero alone. All these bounds are
uniform as $s\to\infty$.

Here is the needed local central limit lemma. If $X$ is any finite
or countable sum of independent Bernoulli variables, has finite
mean $\mu$, and variance $v\to\infty$ in a sequence of such sums,
then
\[
 \sup_{k\in\mathbb Z}\left|
 \sqrt v\,\mathbb P(X=k)
 -\frac1{\sqrt{2\pi}}e^{-(k-\mu)^2/(2v)}
 \right|\longrightarrow0.                            \tag{A016.5}
\]
This statement is uniform over the Bernoulli parameters once $v$
is sufficiently large. To prove it, the centered characteristic
function $\psi$ satisfies, for $|\theta|\le\pi$,
\[
 |\psi(\theta)|\le
 \exp[-2v\sin^2(\theta/2)]
 \le\exp[-2v\theta^2/\pi^2].
\]
For fixed $u$, Taylor expansion of the Bernoulli logarithms gives
$\log\psi(u/\sqrt v)=-u^2/2+O(|u|^3/\sqrt v)$.
The remainder is uniform because the third derivative of each
centered Bernoulli logarithm, for small argument, is bounded by a
constant times $p(1-p)$. Fourier inversion, scaling by $\sqrt v$,
and dominated convergence on the interval
$[-\pi\sqrt v,\pi\sqrt v]$ prove (A016.5): the uniform error
in $k$ is at most the $L^1$ difference between the scaled
characteristic function, extended by zero, and $e^{-u^2/2}$.
The same argument applies to the countable product by finite
truncation. This proves the lemma without a local-limit assumption.

The Palm Campbell formula gives a density for the $j$th positive
point:
\[
 f_j(s)=K^0(s,s)\,
 \mathbb P^{0,s}\{N(0,s)=j-1\},\qquad s>0,            \tag{A016.6}
\]
where $\mathbb P^{0,s}$ is the two-point reduced Palm law.
Indeed summing a test function over points $s$ with exactly $j-1$
points between $0$ and $s$ selects precisely the $j$th point;
conditioning that summand at $s$ gives the right side.
The count in (A016.6) has the Bernoulli eigenvalue description for
$B_s$. Also $K^0(s,s)=1-k(s)^2\to1$.
For $s=j+y\sigma_j$ with $y$ in any fixed compact set, (A016.4)
and (A016.5) therefore give, uniformly in $y$,
\[
 \sigma_j f_j(j+y\sigma_j)
       \longrightarrow\frac{e^{-y^2/2}}{\sqrt{2\pi}}. \tag{A016.7}
\]
Only an $O(1)$ error in the Palm mean is needed, since
$\sigma_j\to\infty$.

To control tails, apply Chebyshev to the one-point Palm count,
whose mean is $s+O(1)$ and variance is (A016.3) up to $O(1)$.
For each fixed large $M$, the identities
$D_j>s\iff N(0,s)\le j-1$ and their lower-tail counterparts give
\[
 \limsup_{j\to\infty}
 \mathbb P(|D_j-j|>M\sigma_j)\le C/M^2.
\]
There are infinitely many positive points almost surely: the same
variance estimate and Chebyshev at dyadic $s$ give this by
Borel--Cantelli. Thus (A016.6) has total integral one.
Compact-uniform density convergence (A016.7), the displayed tail
bound, and Gaussian tail decay prove (A016.1).

Finally write $g_j(y)=\sigma_j f_j(j+\sigma_j y)$ and
$\varphi(y)=e^{-y^2/2}/\sqrt{2\pi}$. For $a\ne0$,
\[
 \left|\mathbb E e^{ia(D_j-j)}-e^{-a^2\sigma_j^2/2}\right|
 \le\|g_j-\varphi\|_{L^1}\longrightarrow0.
\]
This proves (A016.2), and taking $a=\pi k$ proves all nonconstant
Fourier coefficients of the residue-class law vanish. The continuous
periodic function $1-2d_2(x)$ has Haar mean zero, so its expectation
tends to zero. In NN-A003 both signs have zero mean and the joint
law is unchanged by flipping both, because the uniform phase can
be shifted by one. Their four joint probabilities are therefore
$(1\pm C_j)/4$, where $C_j$ is that correlation, and tend to $1/4$.
\end{proof}

The geometry in this result is again the sine projection, finite-rank
Palm modifications, and positive trace-class contractions on intervals.
The full Fredholm generating function supplies local count information
that a pair-correlation statement does not contain. The proof illustrates
the specific strengthening needed to pass from a Gaussian scaling law
to a fixed-frequency periodic statistic. It asserts neither the
arithmetic counterpart nor mixing of every finite label block.


\section{An exact series and certified value for adjacent CUE agreement}

\begin{theorem}[NN-A017: adjacent agreement from integer gap determinants]
Let $K_s$ be the sine-kernel operator on an interval of length $s$,
and put $E(s)=\det(I-K_s)$, the empty-interval probability for the
unit-intensity stationary sine process. The limiting adjacent-label
agreement probability in NN-A003 is
\[
 A_{\rm adj}=1-2\sum_{n=1}^\infty(-1)^{n+1}E(n).
                                                        \tag{A017.1}
\]
The series is absolutely convergent and its alternating partial sums
give certified two-sided bounds. In the specified model, for every
fixed exterior radius $0\le r<1$ and offset,
\[
 0.66654660183721072003
 < A_{\rm adj} <
 0.66654660183721072008.                                \tag{A017.2}
\]
The interval is certified using degree-$100$ Taylor features and
Arb ball arithmetic at $160$ decimal digits; no quadrature or
uncertified extrapolation is used. It is a CUE-model value, not a
proved proportion for arithmetic zero labels.
\end{theorem}

\begin{proof}
Let $D_1$ be the first positive point under the reduced Palm sine law.
The endpoint derivative of a void probability gives
\[
 -E'(s)=\mathbb P^0(D_1>s),\qquad E''(s)=f_1(s),
                                                        \tag{A017.3}
\]
where the second equality may first be read distributionally.
Indeed, enlarging $(0,s)$ by $ds$ loses precisely configurations
with a point in that added strip and no point in the old interval;
the Campbell formula gives the reduced Palm void probability at
the right endpoint. Translation and reflection symmetry identify
it with the probability on the right side of (A017.3).
The chance of two added points is $O(ds^2)$, since the two-point
intensity is bounded by one. Formula (A016.6) supplies its density;
the two-point Palm kernel varies continuously in trace norm after
rescaling each interval to a fixed unit interval, on every compact
set of $s>0$, so that formula also gives continuity. The distributional
form alone is enough below. In particular $E(0)=1$ and $E'(0)=-1$.

Write $g(s)=d_2(1-s)$ for $s\ge0$. It is the triangular function
with $g(0)=1$, right derivative $g'(0)=-1$, and distributional
second derivative
\[
 g''=2\sum_{n\ge1}(-1)^{n+1}\delta_n.
\]
All eigenvalues of $K_s$ lie in $[0,1]$ and their sum is $s$;
therefore $0\le E(s)\le e^{-s}$. Also $E'(s)\to0$ because
$D_1$ is finite almost surely. Twice integrating by parts, first on
a finite interval and then passing to infinity, gives
\[
 \mathbb E d_2(1-D_1)
 =\int_0^\infty g(s)E''(s)ds
 =2\sum_{n\ge1}(-1)^{n+1}E(n).
\]
The two terms at zero cancel: $-g(0)E'(0)=1$ and
$g'(0)E(0)=-1$. The boundary terms at infinity vanish; the series
converges absolutely by $E(n)\le e^{-n}$. NN-A003 now gives
(A017.1). Since void probabilities decrease with interval length,
the alternating remainder after $m$ terms has absolute value at
most $E(m+1)$ and the usual alternating sign.

Here are full error bounds for the supplied numerical certificate.
Center the spatial interval and write $x=sa$, with
$a,\xi\in[-1/2,1/2]$. The Fourier-feature map has kernel
\[
 T(a,\xi)=\sqrt s\,e^{2\pi is a\xi},\qquad K_s\simeq TT^*.
\]
Replace the exponential by its Taylor polynomial of degree $p$,
defining $T_p$. With $h=\pi s/2$ and
$R=e^h h^{p+1}/(p+1)!$, the uniform remainder gives
\[
 \|T-T_p\|_{\rm HS}\le\sqrt s R,\quad
 \|T\|_{\rm HS}=\sqrt s,\quad
 \|T_p\|_{\rm HS}\le\sqrt s(1+R).
\]
Thus, for $A=TT^*$ and $B=T_pT_p^*$,
\[
 \|A-B\|_1\le sR(2+R)=:\delta,
 \qquad \operatorname{Tr}B\le s(1+R)^2.               \tag{A017.4}
\]
For positive finite matrices $A,B$, differentiation along their
positive linear interpolation proves
\[
 |\det(I-A)-\det(I-B)|
 \le\|A-B\|_1e^{\max(\operatorname{Tr}A,
                                    \operatorname{Tr}B)}.
\]
To see this even when the interpolation passes through eigenvalue
one, use the polynomial adjugate derivative. If its positive
eigenvalues are $\lambda_i$, the adjugate operator norm is bounded
by $\prod_i(1+\lambda_i)\le e^{\sum_i\lambda_i}$, since
$|1-\lambda_i|\le1+\lambda_i$. Pairing with $A-B$ uses its
trace norm. For positive trace-class operators, compress by common
finite-rank projections and pass to the limit. The determinant
convergence follows from its absolutely convergent exterior-power
series, bounded by $\sum_k\|A\|_1^k/k!$; trace-norm convergence
of the compressions follows first for finite-rank approximants and
then by approximation. Consequently the certificate's more generous
error bound
\[
 \delta\exp(1+s+s(1+R)^2)                             \tag{A017.5}
\]
is valid.

The finite determinant is evaluated without quadrature. Let
\[
 G_{kl}=\int_{-1/2}^{1/2}u^{k+l}du
 =\begin{cases}
  [2^{k+l}(k+l+1)]^{-1},&k+l\ \text{even},\\
 0,&k+l\ \text{odd},
 \end{cases}
 \quad D_{kk}=\frac{(2\pi is)^k}{k!},\quad0\le k,l\le p.
\]
Sylvester's finite-rank determinant identity gives
\[
 \det(I-B)=\det(I-sDGD^*G).
\]
The even and odd index blocks are real and separate. In those
blocks $(DGD^*)_{kl}=(-1)^{(k-l)/2}
 (2\pi s)^k(2\pi s)^lG_{kl}/(k!l!)$.
The program uses exactly these rational-moment matrices with Arb
balls for $\pi$, and enlarges each determinant ball by (A017.5).
For $p=100$ it certifies the six integer gap probabilities, beginning
with
\begin{align*}
 E(1)&=0.17021742137918523073265305256\ldots,\\
 E(2)&=0.00349732514916909769841342811\ldots.
\end{align*}
and ending with
\[
 E(5)=1.55338899144701217136\ldots\,10^{-14},\qquad
 E(6)=1.89553739144083908383\ldots\,10^{-20}.
\]
Their complete intervals and analytic approximation errors are in
\path{cue_agreement_certificate.json}. The odd partial sum
through five gives a lower bound for $A_{\rm adj}$, and the even
partial sum through six an upper bound. Outward rounding gives
exactly (A017.2). The included program asserts both displayed
decimal inequalities as ball comparisons.
\end{proof}

The determinant is the exterior-algebra trace of the sine projection
restricted to an interval. The finite certificate uses an actual
polynomial feature space and its two Gram matrices, not a finite
field or a claimed arithmetic motive. This supplies a rigorously
computable model statistic, with an analytic approximation error;
it does not certify that arithmetic labels have this limit.


\section{Actual motion, ramified branches, and contour traces}

\begin{theorem}[NN-A004: all shifted-series jets and factorization obstruction]
Let $L(s)=\sum_{n\ge1}a_n n^{-s}$ converge absolutely for $\Re s>1$
and continue holomorphically near $\rho$, where $\Re\rho>0$ and
$L(\rho)=0$. In a sufficiently small neighborhood of $(\rho,0)$, set
\[
 L_t(s)=\sum_{k\ge0}\frac{(-1)^k(s)_k}{k!}L(s+k)t^k.
\]
This is holomorphic and agrees with $\sum a_n(n+t)^{-s}$ in the initial
half-plane. If $\rho$ is simple, its unique analytic zero branch has
\[
 v=\rho'(0)=\frac{\rho L(\rho+1)}{L'(\rho)},
\quad
 \rho''(0)=-\frac{L''(\rho)v^2
 -2\{L(\rho+1)+\rho L'(\rho+1)\}v
 +\rho(\rho+1)L(\rho+2)}{L'(\rho)}.
\]
All higher coefficients $\rho(t)=\rho+\sum_{n\ge1}c_nt^n$ are determined
without a new existence hypothesis by
\[
 c_n=-\frac{[t^n]L_t(\rho+\sum_{j<n}c_jt^j)}{L'(\rho)}.
\]
For real $t\downarrow0$, the first nonzero $\Re c_n$ determines the
normal side of the critical line. If every $\Re c_n=0$, the branch stays
on that line near zero. A nonzero complex velocity alone does not exclude
normal tangency.

If $L=FG$, $F(\rho)=0$, $F'(\rho)G(\rho)\ne0$, then the shifts of the
Dirichlet coefficients satisfy
\[
 v_L=v_F\,\frac{G(\rho+1)}{G(\rho)}.
\]
Consequently factorwise GRH does not give a factorwise rule for the motion
labels: the nonvanishing factor has a retained complex multiplier.
\end{theorem}
\begin{proof}
For $s$ in a fixed compact neighborhood of $\rho$, choose $k_0$ such that
$\Re(s+k_0)>1$. Absolute convergence gives a uniform bound for $L(s+k)$
for $k\ge k_0$, while $(s)_k/k!=O(k^C)$ uniformly on this compact set.
The displayed series and its local holomorphic derivatives converge for
$|t|<1$. There are only finitely many lower $k$, and their neighborhoods
can avoid any poles. Expanding $(n+t)^{-s}$ by the binomial theorem and
using absolute convergence proves equality in the initial half-plane.
The implicit function theorem applies because $L'(\rho)\ne0$.
The first two derivatives of $L_t(\rho(t))=0$ use
$\partial_tL_t|_0=-sL(s+1)$ and
$\partial_t^2L_t|_0=s(s+1)L(s+2)$, and give the stated formulas.
Adding $c_nt^n$ to the already constructed truncation changes its degree
$n$ coefficient by exactly $L'(\rho)c_n$, proving the recursion. A
real analytic function is either identically zero locally or has the sign
of its first nonzero coefficient for small positive argument. Finally
$L'(\rho)=F'(\rho)G(\rho)$ and
$L(\rho+1)=F(\rho+1)G(\rho+1)$ prove the product formula.
\end{proof}

This applies to $\zeta$, Dirichlet $L$-functions, and Dedekind $\zeta_K$.
A supplies the critical line for $\zeta$ and Dirichlet factors; for abelian
$K/\mathbb Q$ the standard Dirichlet factorization transfers that location
to $\zeta_K$. General $K$ requires the stipulated Hecke input M, or a
separate GRH hypothesis. None of these inputs supplies simplicity. At a
nontrivial critical zero, the Euler product at $\rho+1$ shows $v\ne0$
when the zero is simple; it gives no sign for $\Re v$.
The geometric object here is the analytic zero divisor in the two
coordinates $(s,t)$ and its sections over the parameter disc. It is not
an algebraic family over a finite field, and its complex multipliers have
not been identified with Frobenius eigenvalues.

\begin{theorem}[NN-A005: multiple-zero branches and the residue trace]
Suppose instead that $L$ has multiplicity $m\ge2$ at $\rho$ and
$\rho L(\rho+1)\ne0$. Put
\[
 C=\frac{m!\rho L(\rho+1)}{L^{(m)}(\rho)}.
\]
The $m$ local branches on the ramified parameter $u=t^{1/m}$ satisfy
\[
 s_j(u)=\rho+c\omega_m^j u+O(u^2),\qquad c^m=C.
\]
In particular they are not individually differentiable as holomorphic
functions of $t$ at zero. Their symmetric trace is analytic. More
generally, let $\Gamma$ be a positively oriented finite union of simple
contours, bounding a region where $L_t$ is holomorphic and has no boundary
zeros for small $t$. If $P$ is holomorphic on that region, then
\[
 \left.\frac d{dt}\sum_{L_t(z)=0}P(z)\right|_{t=0}
 =\frac1{2\pi i}\int_\Gamma
 P'(s)\frac{sL(s+1)}{L(s)}\,ds.
\]
Zeros in this formula are counted with multiplicity. For $P(s)=s$, the
contribution of a multiplicity-$m$ zero is
\[
 \operatorname{Res}_{s=\rho}\frac{sL(s+1)}{L(s)}
 =\frac1{(m-1)!}
 \left.\left(\frac d{ds}\right)^{m-1}
 \frac{sL(s+1)}{L(s)/(s-\rho)^m}\right|_{s=\rho}.
\]
Thus a contour mean remains meaningful without simple zeros, whereas a
sum of the simple-zero velocities does not.
\end{theorem}
\begin{proof}
With $a=L^{(m)}(\rho)/m!$, the quotient
$u^{-m}L_{u^m}(\rho+uy)$ extends holomorphically to $u=0$ and there
equals $ay^m-\rho L(\rho+1)$. Its $m$ roots are nonzero and simple.
The implicit function theorem gives the $m$ analytic $y_j(u)$, and
Rouche's theorem counts all roots in the local cluster. The nonzero
leading terms $c\omega_m^j t^{1/m}$ exclude holomorphic individual
branches in $t$.
The argument principle writes the symmetric sum as
$(2\pi i)^{-1}\int_\Gamma P(s)\partial_s\log L_t(s)\,ds$.
The quotient $\partial_tL_t/L_t$ is analytic in $t$ on the compact
contour, so differentiation under the integral and integration by parts
give $-(2\pi i)^{-1}\int P'(s)\partial_tL_t/L_t$. Insert
$\partial_tL_t|_0=-sL(s+1)$. The residue expression is the ordinary
formula for a pole of order at most $m$.
\end{proof}

For $L=\zeta$, the right-hand contour integrand has the fully explicit
Dirichlet series, initially on $\Re s>1$,
\[
 \frac{s\zeta(s+1)}{\zeta(s)}
 =s\sum_{n\ge1}\frac{b(n)}{n^s},\qquad
 b(n)=\frac1n\prod_{p\mid n}(1-p).
\]
Indeed convolution with $1/\zeta(s)=\sum\mu(n)n^{-s}$ gives
$b(n)=n^{-1}\sum_{d\mid n}d\mu(d)$. For general $L$, replace
$b(n)$ by $\sum_{d\mid n}a_d d^{-1}a^{-1}_{n/d}$ whenever the
Dirichlet inverse converges absolutely. This is an exact residue
formula, not an unproved asymptotic on contours passing through zeros.
In the zero-divisor geometry, symmetric sums are functions on the relative
degree-$m$ divisor (or its symmetric power); individual sections require
the stated ramified base change. No complexity bound follows from the
existence of the contour integral.

\begin{proposition}[NN-A010: the completed Hurwitz deformation has extra poles]
Let $\Lambda_t(s)=\pi^{-s/2}\Gamma(s/2)\zeta(s,1+t)$. For $m\ge1$,
\[
 \operatorname{Res}_{s=-2m}\Lambda_t(s)
 =-\frac{2(-1)^m\pi^m}{m!(2m+1)}B_{2m+1}(1+t).
\]
At $t=0$ these residues vanish; their first $t$ derivatives are
$-2(-1)^m\pi^mB_{2m}/m!\ne0$. At $t=1$, every such residue is
$-2(-1)^m\pi^m/m!$. Consequently the rational principal-part
bookkeeping using the poles at zero and one is not the complete
meromorphic divisor of $\Lambda_t$. It does not provide a meromorphic
function on $\mathbb P^1$ for the full completion.
\end{proposition}
\begin{proof}
The residue of $\Gamma(s/2)$ at $s=-2m$ is $2(-1)^m/m!$.
Use the exact Hurwitz value
$\zeta(-2m,a)=-B_{2m+1}(a)/(2m+1)$ (DLMF 25.11.14).
Bernoulli differentiation and $B_{2m}(1)=B_{2m}$ give the derivative.
For $t=1$, $\zeta(s,2)=\zeta(s)-1$, so its value at every negative
even integer is $-1$. A meromorphic function on compact
$\mathbb P^1$ has finitely many poles, whereas this example has
infinitely many. This proves the asserted scope restriction.
\end{proof}


\section{Central Dedekind branches and an unconditional arithmetic imbalance}
\label{sec:analytic-central-puiseux}

\begin{theorem}[NN-A018: the second central Puiseux coefficient]
Let $K$ have degree $n=r_1+2r_2$ and absolute discriminant $D_K$.
Write $L=\zeta_K$ and suppose its zero at $\rho=1/2$ has order
$m\ge2$. Put
\[
 Q_K(s)=s(s-1)D_K^{s/2}\Gamma_{\mathbb R}(s)^{r_1}
           \Gamma_{\mathbb C}(s)^{r_2},\qquad
 \Gamma_{\mathbb R}(s)=\pi^{-s/2}\Gamma(s/2),\quad
 \Gamma_{\mathbb C}(s)=(2\pi)^{-s}\Gamma(s),
\]
so that $\xi_K=Q_KL$ is entire and satisfies
$\xi_K(s)=\xi_K(1-s)$. The harmless constant often inserted in
$\Gamma_{\mathbb C}$ changes none of the formulas below. Set
\[
 a_m=\frac{L^{(m)}(1/2)}{m!},\qquad
 C_K=2+\frac{\zeta_K'}{\zeta_K}(3/2)
                 +\frac{Q_K'}{Q_K}(1/2).
\]
The shifted Dirichlet-series deformation from NN-A004 has exactly $m$
local branches, counted with multiplicity. For positive real $t$ they
are distinct for all sufficiently small $t$ and have expansions
\begin{equation}\label{eq:central-second-jet}
 s_j(t)=\frac12+c_jt^{1/m}
              +\frac{C_K}{m}c_j^2t^{2/m}+O(t^{3/m}),\qquad
 c_j^m=\frac{\zeta_K(3/2)}{2a_m}.
\end{equation}
The central order $m$ is even. If $a_m<0$, let
$a=(\zeta_K(3/2)/(2|a_m|))^{1/m}$ and choose
$c_j=a\exp((2j+1)\pi i/m)$. If $m\equiv0\pmod4$, exactly $m/2$
branches lie on each side of the critical line for small positive $t$.
If $m\equiv2\pmod4$, precisely two leading directions are vertical.
When $C_K>0$ these two branches bend left, and the left/right counts
are $(m+2)/2,(m-2)/2$; when $C_K<0$ both bend right and these counts
are reversed. No conclusion about their normal direction is made
from this coefficient when $C_K=0$.

The hypothesis $a_m<0$ follows from GRH for $\zeta_K$, hence from A
for abelian $K$ and conditionally from M in general. It can also be
established without GRH by a local factorization, as in the example
below.
\end{theorem}

\begin{proof}
Evenness of $\xi_K(1/2+z)$ and the nonvanishing of $Q_K(1/2)$ imply
that $m$ is even. Write
\[
 L(1/2+z)=a_mz^m+a_{m+1}z^{m+1}+O(z^{m+2}).
\]
The coefficient of $z^{m+1}$ in $Q_K(1/2+z)L(1/2+z)$ vanishes, so
\begin{equation}\label{eq:central-unit-jet}
 \frac{a_{m+1}}{a_m}=-\frac{Q_K'}{Q_K}(1/2).
\end{equation}
The deformation has the local expansion
\[
 L_t(1/2+z)=a_mz^m+a_{m+1}z^{m+1}
       -t(B_0+B_1z)+O(z^{m+2})+O(tz^2)+O(t^2),
\]
where $B_0=L(3/2)/2>0$ and
$B_1=L(3/2)+L'(3/2)/2$. By NN-A005, the base change $t=u^m$
produces $m$ holomorphic branches $z=c_ju+d_ju^2+O(u^3)$,
with distinct $c_j$ and $a_mc_j^m=B_0$. The coefficient of $u^{m+1}$
is
\[
 ma_mc_j^{m-1}d_j+a_{m+1}c_j^{m+1}-B_1c_j=0.
\]
Because $2m\ge m+2$, the $O(t^2)$ term does not contribute.
Using $a_mc_j^m=B_0$ and \eqref{eq:central-unit-jet} gives
\[
 d_j=\frac{c_j^2}{m}
       \left(\frac{B_1}{B_0}-\frac{a_{m+1}}{a_m}\right)
     =\frac{C_K}{m}c_j^2.
\]
This proves \eqref{eq:central-second-jet} and the simplicity of
the branches at small nonzero $t$.

For $a_m<0$ the displayed $m$ directions are the roots of a negative
real number. A direction is vertical exactly when $m\equiv2\pmod4$,
and then the two vertical directions are $\pm ia$. Their second
normal coefficient is $-a^2C_K/m$, the same for both. The other
directions pair across the imaginary axis. This gives all the
claimed counts. The $O(t^{3/m})$ error is smaller than the nonzero
first or second normal term, so the statements hold on a common
positive interval for all $m$ branches.

Under GRH, $\xi_K$ has no real zero in $(1/2,\infty)$ and is positive
on $(1,\infty)$. Its value at $1$ is the positive Dedekind residue
times positive factors; thus it remains positive on $(1/2,\infty)$.
Its first nonzero coefficient at $1/2$ is consequently positive.
Since $Q_K(1/2)<0$, this gives $a_m<0$.
\end{proof}

\begin{proposition}[An effective conductor criterion]
Let
\[
 a_{\mathbb R}=\frac12\psi(1/4)-\frac12\log\pi,
 \qquad C_*=-a_{\mathbb R}-\frac{\zeta'}{\zeta}(3/2).
\]
Then
\begin{equation}\label{eq:central-C-bound}
 C_K\ge2+\frac12\log D_K-nC_*+\frac\pi2r_2.
\end{equation}
In particular, root discriminant $D_K^{1/n}\ge4371$ implies $C_K>0$.
Consequently, whenever $a_m<0$ and $m\equiv2\pmod4$, the two
initially vertical branches bend left for such a field.
\end{proposition}

\begin{proof}
The logarithmic derivative of $s(s-1)$ vanishes at $1/2$.
The gamma reflection and duplication formulas give
\[
 \frac{Q_K'}{Q_K}(1/2)=\frac12\log D_K
                   +na_{\mathbb R}+\frac\pi2r_2.
\]
For a rational prime $p$, write $f_{\mathfrak p}$ and
$e_{\mathfrak p}$ for the residue and ramification degrees.
For $\sigma>1$,
\[
 -\frac{\zeta_K'}{\zeta_K}(\sigma)
 =\sum_p\log p\sum_{\mathfrak p\mid p}
       \frac{f_{\mathfrak p}}{p^{f_{\mathfrak p}\sigma}-1}
 \le n\sum_p\frac{\log p}{p^\sigma-1}
 =-n\frac{\zeta'}{\zeta}(\sigma).
\]
Here $p^{f\sigma}-1\ge p^\sigma-1$ and
$\sum f_{\mathfrak p}\le\sum e_{\mathfrak p}f_{\mathfrak p}=n$.
This proves \eqref{eq:central-C-bound}. The saved Arb certificate
gives $e^{2C_*}<4371$ by a strict interval comparison; hence the
right side of \eqref{eq:central-C-bound} is positive at the stated
root discriminant. One may instead use the exact threshold
$\log D_K>2nC_*-\pi r_2-4$ without decimal constants.
\end{proof}

\begin{corollary}[An unconditional arithmetic example with two excess right branches]
Let
\[
 F=\mathbb Q(\sqrt5,\sqrt{41}),\qquad
 N=F\left(\sqrt{\frac{(5+\sqrt5)(41+\sqrt{205})}{4}}\right).
\]
The central order $m=\operatorname{ord}_{s=1/2}\zeta_N(s)$ is
positive and congruent to $2$ modulo $4$. For the shifted
Dirichlet-series deformation of $\zeta_N$, and every sufficiently
small positive $t$, its $m$ branches from the central zero consist
of exactly $(m+2)/2$ right branches and $(m-2)/2$ left branches.
The two initially vertical branches both bend right. This assertion
uses neither A nor M and does not assume $m=2$.
\end{corollary}

\begin{proof}
The field, its totally real tame $Q_8$ Galois group, ramification
only at $5,41$, and the nonfreeness of $\mathcal O_N$ over
$\mathbb Z[Q_8]$ are Martinet's Example~3, printed p.~407, with
Theorem~IV.1 on p.~406. They are classical inputs, already used for
central existence in prior MT-ART10; central existence is not the
new claim here. Fr\"ohlich's Theorem~1, printed p.~82 of his 1974
paper, identifies this nontrivial module class with root number
$-1$ for the unique irreducible character $\chi$ of degree two.

Choose a cyclic subgroup of order four. Its faithful linear
character induces $\chi$, so Artin induction identifies $L(s,\chi)$
with a nontrivial finite-order Hecke $L$-function over the quadratic
fixed field. It is entire. Its selfdual functional equation with
sign $-1$ makes its central order $\ell$ positive and odd. Its
Dirichlet coefficients are real, so its leading central coefficient
$b_\ell$ is a nonzero real number. The regular representation gives
the exact all-prime factorization
\[
 \zeta_N(s)=\zeta(s)L(s,\chi_5)L(s,\chi_{41})L(s,\chi_{205})
                     L(s,\chi)^2.
\]
The included 70-digit Arb certificate, using the finite Hurwitz
formula, proves
\[
 0.23<L(1/2,\chi_5)<0.24,\quad
 1.10<L(1/2,\chi_{41})<1.11,\quad
 1.17<L(1/2,\chi_{205})<1.18.
\]
Also $\zeta(1/2)<0$: the alternating eta series is positive by
pairing successive terms and $1-2^{1-s}<0$ for $0<s<1$.
It follows that $m=2\ell\equiv2\pmod4$ and the leading coefficient
of $\zeta_N$ is
$a_m=\zeta_F(1/2)b_\ell^2<0$, without a zero-location hypothesis.

The discriminant is $D_N=205^6$. For completeness, at either
$p=5$ or $41$ the inertia image in
$\operatorname{Gal}(F/\mathbb Q)=Q_8/\{\pm1\}$ is nontrivial.
Tame inertia in $Q_8$ is cyclic; any element with nontrivial image
in that quotient has order four. Thus the inertia order is four
and the tame discriminant exponent is $8(1-1/4)=6$. There are no
other ramified primes. This also agrees with Martinet's
Propositions~III.1--III.2, printed pp.~404--405.

The Euler product gives $\zeta_N'/\zeta_N(3/2)<0$ strictly.
Since $N$ is totally real of degree eight, the same certificate gives
\[
 C_N<2+3\log205+8a_{\mathbb R}
       <-3.5197<0.
\]
The theorem now gives the stated right/left branch counts.
\end{proof}

\paragraph{Source and verification locators.}
Martinet, \emph{Modules sur l'alg\`ebre du groupe quaternionien},
Ann. Sci. ENS~4 (1971), 399--408,
DOI~10.24033/asens.1216, Propositions~III.1--III.2,
Theorem~IV.1 and Example~3. The original p.~407 image was checked
again here: its second radical is $\sqrt{5\cdot41}$.
Fr\"ohlich, \emph{The Galois module structure of algebraic integer
rings in fields with generalised quaternion group},
M\'em. SMF~37 (1974), 81--86, DOI~10.24033/msmf.132,
Theorem~1 p.~82. This is an imported classical theorem; its
referenced 1972 proof is not claimed to have been re-proved here.
The short local deduction above is complete.
The script \path{check_central_jets.py} uses exact Kronecker
characters and
\[
 L(s,\chi_d)=d^{-s}\sum_{a=1}^d\chi_d(a)\zeta(s,a/d)
\]
with complex-ball Hurwitz evaluation, and writes every sign interval
to \path{central_jets_certificate.json}.

\paragraph{Geometric content and limitation.}
The central zero has local divisor length $m$; after the ramified
base change $t=u^m$, the normalization has $m$ sections. The
coefficient of $u^2$ depends on the completion unit and the shifted
Euler value, data not retained by the reduced support. On the chosen
positive real parameter ray the theorem assigns a canonical
multiset of sides. It does not assign a single unramified Boolean
label to the central support point, nor a colouring invariant under
complex monodromy. In the example, a nontrivial integral
$\mathbb Z[Q_8]$-module class supplies the odd Artin order, while
the analytic second jet supplies the excess of two right branches.
Those are precise distinct steps; no finite pure motive or
complexity lower bound is inferred.


\section{A uniform central Dirichlet second-jet sign under A}
\label{sec:analytic-central-dirichlet-sign}

Let $\chi$ be a nonprincipal real primitive Dirichlet character of
conductor $q$, and put $a\in\{0,1\}$ for its parity,
$\chi(-1)=(-1)^a$. Throughout this section the completion is
\begin{equation}\label{eq:analytic-dirichlet-completion}
 Q_\chi(s)=\left(\frac q\pi\right)^{(s+a)/2}
                \Gamma\left(\frac{s+a}{2}\right),\qquad
 \Lambda_\chi(s)=Q_\chi(s)L(s,\chi).
\end{equation}
This fixes every gamma and conductor factor used in the local
coefficient. The principal primitive character, whose completion
has poles, is excluded.

\begin{theorem}[NN-A032: positive central coefficient and exact side counts]
\label{thm:analytic-central-dirichlet-sign}
Let $m=\operatorname{ord}_{s=1/2}L(s,\chi)\ge0$. It is even.
The coefficient $C_\chi$ defined below satisfies, unconditionally,
$C_\chi>m+2-\pi/2$ for $a=0$ and
$C_\chi>m+\pi/2$ for $a=1$.
Under A, list the positive ordinates of all noncentral zeros
with their multiplicities. Then the exact identity is
\begin{equation}\label{eq:analytic-dirichlet-C-positive}
 \begin{split}
 C_\chi&:=2+\frac{L'}L(3/2,\chi)
                         +\frac{Q_\chi'}{Q_\chi}(1/2)\\
 &=m+\sum_{\gamma>0}\frac2{1+\gamma^2}
     +\begin{cases}2-\pi/2,&a=0,\\ \pi/2,&a=1.
       \end{cases}
 \end{split}
\end{equation}
The sum converges absolutely. In particular,
\[
 C_\chi\ge m+2-\pi/2>0\quad(a=0),\qquad
 C_\chi\ge m+\pi/2>0\quad(a=1).
\]

Assume A for the rest of the theorem. If $m>0$, write
$A_m=L^{(m)}(1/2,\chi)/m!$. Then $A_m>0$. For the shifted
Dirichlet deformation of NN-A004,
\[
 L_t(s,\chi)=\sum_{n\ge1}\frac{\chi(n)}{(n+t)^s},
\]
continued analytically in $s$ and locally in $t$, its $m$ zero
branches near the center have, for positive $t$, the expansions
\begin{equation}\label{eq:analytic-dirichlet-central-branches}
 s_j(t)=\frac12+c_jt^{1/m}
                  +\frac{C_\chi}{m}c_j^2t^{2/m}
                  +O(t^{3/m}),\qquad
 c_j^m=\frac{L(3/2,\chi)}{2A_m}>0.
\end{equation}
For all sufficiently small $t>0$, the exact counts are
\[
 \begin{array}{c|cc}
 &\Re s_j(t)<1/2&\Re s_j(t)>1/2\\ \hline
 m\equiv2\pmod4&m/2&m/2\\
 m\equiv0\pmod4,\ m\ge4&m/2+1&m/2-1.
 \end{array}
\]
In the second row the two branches with purely imaginary first
coefficient both bend strictly left at second order. No central
zero is asserted to exist; the side-count statement is
conditional on $m>0$. Neither quadratic central nonvanishing nor
simplicity follows from this theorem.
\end{theorem}

\begin{proof}
\emph{The exact functional-equation sign and entire completion.}
A nonprincipal real primitive character is quadratic. Its
quadratic Gauss sum has the sign
$\tau(\chi)=i^a\sqrt q$, so the root number in the convention
\eqref{eq:analytic-dirichlet-completion} is $+1$. A primary
locator, including the sign rather than just its square, is
B.~H.~Gross, \emph{Some remarks on signs in functional equations},
section~1, equations~(1.3)--(1.6), p.~2. There the sign is proved
by writing the completed quadratic $L$-function as the quotient
of the completed zeta functions of the quadratic field and
$\mathbb Q$. DLMF~25.15.5 gives the primitive functional equation
with the Gauss sum explicitly.

For completeness, the usual order-one analytic input can be
seen directly with these normalizations. Put
\[
 \theta_\chi(x)=\sum_{n\in\mathbb Z}n^a\chi(n)e^{-\pi n^2x/q},
 \qquad x>0,
\]
where the $n=0$ term is zero. Poisson summation in the residue
classes modulo $q$, using
\[
 \sum_{b\bmod q}\chi(b)e^{2\pi inb/q}=\tau(\chi)\chi(n),
\]
and the Gaussian Fourier transform gives
\[
 \theta_\chi(x)=\frac{(-i)^a\tau(\chi)}{\sqrt q}
        x^{-a-1/2}\theta_\chi(1/x)
              =x^{-a-1/2}\theta_\chi(1/x).
\]
For $a=1$ the needed transform follows by differentiating the
Gaussian transform once; this accounts for the factor $-i$.
Termwise Mellin integration for $\Re s>1$, followed by the
substitution $x\mapsto1/x$ on $(0,1)$, yields
\begin{equation}\label{eq:analytic-dirichlet-theta-entire}
 \Lambda_\chi(s)=\frac12\int_1^\infty\theta_\chi(x)
       \left(x^{(s+a)/2}+x^{(1-s+a)/2}\right)\frac{dx}{x}.
\end{equation}
The theta series decays exponentially for $x\ge1$; hence this
formula is entire, is invariant under $s\mapsto1-s$, and gives
$\log\max_{|s|\le R}|\Lambda_\chi(s)|=O_q(R\log(R+2))$ by
comparison with a gamma integral. Thus the entire function has
order at most one. Its zeros are precisely the nontrivial
$L$-zeros: the gamma poles are canceled by the trivial zeros,
with nonzero resulting values by the functional equation and
the Euler product on the opposite half-plane. At the endpoint
$s=0$ for an even character, the reflected point is $s=1$;
here one also uses Dirichlet's nonvanishing theorem
$L(1,\chi)\ne0$ (DLMF~25.15.9, with source Apostol, 1976,
p.~149), rather than the Euler product in its open half-plane.

\emph{Unconditional sign of the second-jet coefficient.}
Write $F(z)=\Lambda_\chi(1/2+z)$. All its zeros
$\alpha=\rho-1/2$ lie in $|\Re\alpha|\le1/2$ by the classical
critical-strip theorem, independently of A. The paired
Hadamard argument, exactly as in NN-A033, gives
\[
 \frac{F'(1)}{F(1)}
       =m+\sum_{\{\alpha,-\alpha\},\ \alpha\ne0}
                         \frac2{1-\alpha^2}.
\]
The series is absolutely convergent and is grouped under
conjugation. If $\alpha=b+ic$, then
$\Re(1-\alpha^2)=1-b^2+c^2>0$, whence
$\Re(1-\alpha^2)^{-1}>0$. Purely real and purely imaginary
pairs also contribute positive terms. There are noncentral
zeros, since the primitive $L$-function has infinitely many
zeros in the strip (DLMF~25.15(ii)). Hence $F'(1)/F(1)>m$.
The completion identity and the gamma differences calculated
below give, with no zero-location assumption,
\[
 C_\chi=\frac{F'(1)}{F(1)}+
       \begin{cases}2-\pi/2,&a=0,\\ \pi/2,&a=1.
       \end{cases}
\]
This proves the unconditional strict inequalities. It does not
fix the sign of the central leading coefficient without A.

\emph{Real positivity and the paired zero product under A.}
For real $\sigma>1$, every factor of the Euler product
\[
            L(\sigma,\chi)=\prod_p(1-\chi(p)p^{-\sigma})^{-1}
\]
is positive; this is DLMF~25.15.2. The gamma and conductor
factors are also positive. Assumption A excludes real zeros in
$(1/2,\infty)$, so continuity gives
$\Lambda_\chi(\sigma)>0$ throughout this interval. The
sign-$+1$ functional equation
makes $F$ even, so its central order $m$ is even. Its central
leading coefficient is positive, by its positivity for $z>0$.
Dividing by the positive number $Q_\chi(1/2)$ gives $A_m>0$
when $m>0$.

Under A all noncentral zeros of $F$ are $\pm i\gamma$. The
paired genus-one Hadamard product is
\begin{equation}\label{eq:analytic-dirichlet-positive-product}
 F(z)=Bz^m\prod_{\gamma>0}\left(1+\frac{z^2}{\gamma^2}\right),
                         \qquad B>0.
\end{equation}
Indeed the order-one factorization initially permits an
exponential $e^{b z}$; evenness forces $b=0$, and the paired
genus-one exponentials cancel. The sum of reciprocal squared
ordinates converges, so the product and its logarithmic
derivative converge normally on compact sets avoiding zeros.
At $z=1$ this gives the exact positive formula
\begin{equation}\label{eq:analytic-dirichlet-log-derivative}
       \frac{\Lambda_\chi'}{\Lambda_\chi}(3/2)
                    =m+\sum_{\gamma>0}\frac2{1+\gamma^2}.
\end{equation}

\emph{Cancellation of the conductor and the gamma differences.}
Since $L'/L=\Lambda_\chi'/\Lambda_\chi-Q_\chi'/Q_\chi$,
\[
 C_\chi=\frac{\Lambda_\chi'}{\Lambda_\chi}(3/2)
          +2+\frac12\left[
              \psi\left(\frac{1/2+a}{2}\right)
             -\psi\left(\frac{3/2+a}{2}\right)\right].
\]
The conductor terms have canceled exactly. The digamma
reflection and recurrence identities, DLMF~5.5.4 and~5.5.2,
give
\[
 \psi(1/4)-\psi(3/4)=-\pi,\qquad
 \psi(3/4)-\psi(5/4)=\pi-4.
\]
Substitution in \eqref{eq:analytic-dirichlet-log-derivative}
proves \eqref{eq:analytic-dirichlet-C-positive} and both strict
sign bounds. This uses the actual completion unit, without a
conductor-size estimate or a numerical verification.

\emph{The shifted Puiseux coefficient.}
Assume $m>0$, hence $m\ge2$. Write $x=s-1/2$ and
\[
 L(1/2+x,\chi)=A_mx^m+A_{m+1}x^{m+1}+O(x^{m+2}).
\]
Evenness of $\Lambda_\chi$ gives
$A_{m+1}/A_m=-Q_\chi'/Q_\chi(1/2)$. The shifted series of
NN-A004 is locally
\[
 L_t(s,\chi)=L(s,\chi)-sL(s+1,\chi)t+O(t^2).
\]
Put $t=u^m$ and $x=cu+du^2+O(u^3)$. The coefficient of $u^m$
gives $A_mc^m=L(3/2,\chi)/2$. Its $m$ roots are distinct.
The coefficient of $u^{m+1}$ then gives
\[
 mA_mc^{m-1}d+A_{m+1}c^{m+1}
       -\left(L(3/2,\chi)+\tfrac12L'(3/2,\chi)\right)c=0,
\]
or $d=(C_\chi/m)c^2$. Since $2m\ge m+2$, the $O(t^2)$ term
does not affect these two equations. The analytic implicit
function theorem after dividing by $u^m$ at each distinct
leading root gives all $m$ analytic branches in $u$, with the
stated error. This proves
\eqref{eq:analytic-dirichlet-central-branches}.

Finally put $R=(L(3/2,\chi)/(2A_m))^{1/m}>0$. The leading
coefficients are $c_j=Re^{2\pi ij/m}$. If $m\equiv2\pmod4$,
none is purely imaginary and exactly half have each sign of
their real part. If $m\equiv0\pmod4$, precisely two are
$\pm iR$; each has second coefficient
$-C_\chi R^2/m<0$. All other branches keep the sign of their
nonzero first real coefficient. Counting them gives the two
rows of the table, with no branch left on the line for
sufficiently small positive $t$.
\end{proof}

\paragraph{Geometric translation.}
The quadratic character is the rank-one Artin representation of
its quadratic finite etale cover away from the conductor; its
root number is the actual $+1$ in Gross's equations~(1.3)--(1.6).
At a hypothetical central zero of length $m$, the shifted
analytic divisor is finite flat of rank $m$ near the parameter
origin. The ramified base change $t=u^m$ separates its leading
sections. Complex conjugation pairs the two sections with
leading coefficients $\pm iR$, and the strictly negative real
second coefficient puts both on the left of the chosen real
normal line. This is a statement about the full local analytic
algebra and its completion unit, rather than a single label on
its reduced support. The finite-free discriminant/etale
criterion is SGA~1, Expos\'e~I, section~4, Proposition~10; no
arithmetic family realizing the shifted parameter, central
nonvanishing theorem, or positive integral cohomology is inferred.


\section{An unconditional central-order certificate and exactly two right branches}
\label{sec:analytic-central-order-certificate}

The next argument uses only the classical critical strip, not a
Riemann hypothesis. It complements NN-A018 by determining the
central multiplicity of its explicit quaternionic example.
The central order itself is treated as a classical type of
calculation, without a historical novelty claim. The new
catalogue consequence is the complete two-branch conclusion for
the specified shifted deformation, together with the fully
written rational certificate below.

For a number field $K$ of signature $(r_1,r_2)$, use the exact
pole-removed completion
\begin{equation}\label{eq:analytic-order-completion}
 \xi_K(s)=s(s-1)D_K^{s/2}\pi^{-r_1s/2}(2\pi)^{-r_2s}
        \Gamma(s/2)^{r_1}\Gamma(s)^{r_2}\zeta_K(s).
\end{equation}
It is real entire of order one, is invariant under $s\mapsto1-s$,
and its zeros are the nontrivial zeros in $0\le\Re s\le1$,
with multiplicity. A primary locator for these analytic inputs
and the completion is Hasanalizade--Shen--Wong,
\emph{Counting zeros of Dedekind zeta functions},
\emph{Math. Comp.}~91 (2022), section~2.1, equations~(2.1)--(2.2),
p.~279. Its complex-place gamma convention differs from
\eqref{eq:analytic-order-completion} by the positive constant
$2^{r_2}$, by the duplication formula; every logarithmic
derivative here is therefore unchanged. The Euler product gives
nonvanishing in $\Re s>1$, and the functional equation and the
nonzero residues at $0,1$ remove any additional completed zeros.

\begin{theorem}[NN-A033: a positive rational test for central order]
\label{thm:analytic-central-order-test}
Put $F_K(z)=\xi_K(1/2+z)$,
$G_K(z)=F_K'(z)/F_K(z)$, and
$m_K=\operatorname{ord}_{s=1/2}\zeta_K(s)$. For every real
$\delta>1/\sqrt2$,
\begin{equation}\label{eq:analytic-T2-order}
 m_K\le T_{2,K}(\delta):=
 \frac{\delta G_K(\delta)-\delta^2G_K'(\delta)}2
 =m_K+\sum_{\{\alpha,-\alpha\}}
                 \frac{2\delta^4}{(\delta^2-\alpha^2)^2}.
\end{equation}
Here $\alpha=\rho-1/2\ne0$, one representative is chosen for
each pair $\{\alpha,-\alpha\}$, multiplicities are retained,
and conjugate terms are grouped. Each grouped contribution is
strictly positive. The series is absolutely convergent.

Set $\sigma=1/2+\delta$ and
\[
 B_K(s)=\frac1s+\frac1{s-1}+\frac12\log D_K
       -\frac{r_1}2\log\pi-r_2\log(2\pi)
       +\frac{r_1}2\psi(s/2)+r_2\psi(s).
\]
Then the entirely archimedean upper bound is
\begin{equation}\label{eq:analytic-T2-arch}
 T_{2,K}(\delta)<
 U_K(\delta):=
       \frac{\delta B_K(\sigma)-\delta^2B_K'(\sigma)}2.
\end{equation}
Both assertions are unconditional.
\end{theorem}

\begin{proof}
Evenness gives $m_K$ even. Hadamard factorization and pairing
opposite zeros give
\[
 F_K(z)=c z^{m_K}\prod_{\{\alpha,-\alpha\}}
                      (1-z^2/\alpha^2),\qquad c\ne0.
\]
The product converges normally because the order-one zero
divisor satisfies $\sum_{\alpha\ne0}|\alpha|^{-2}<\infty$.
Pairing the genus-one factors cancels their exponential parts;
any remaining factor $e^{az+b}$ has $a=0$ by evenness.
The point $\delta$ is not a zero, since
$\sigma>1$. Logarithmic differentiation on a neighbourhood of
that point is consequently legitimate and gives
\[
 G_K(z)=\frac{m_K}{z}
       +\sum_{\{\alpha,-\alpha\}}
                        \frac{2z}{z^2-\alpha^2}.
\]
Differentiating once more, the contribution of a noncentral
pair to $(\delta G_K-\delta^2G_K')/2$ is
$2\delta^4/(\delta^2-\alpha^2)^2$, and the central contribution
is $m_K$. This also proves absolute convergence of the series
in \eqref{eq:analytic-T2-order}, whose tail is
$O_\delta(\sum |\alpha|^{-4})$.

It remains to check positivity without moving zeros to the
critical line. Write $\alpha=a+ib$, where $|a|\le1/2$. Put
\[
 A=\delta^2-a^2+b^2,\qquad B=2ab.
\]
Then
\[
 A-|B|=\delta^2-2a^2+(|b|-|a|)^2
                    >0.
\]
For a genuine quartet $\alpha,-\alpha,\bar\alpha,-\bar\alpha$,
the two paired contributions therefore have positive real sum,
because
\[
 \Re\frac1{(\delta^2-\alpha^2)^2}
             =\frac{A^2-B^2}{(A^2+B^2)^2}>0.
\]
If the pair is real or purely imaginary, its single contribution
is itself positive. This proves the first inequality with all
possible off-line zeros included.

For real $\sigma>1$, the ideal Euler product gives the absolutely
convergent positive sums
\[
 E_1(\sigma)=\sum_{\mathfrak p}\sum_{k\ge1}
           \frac{\log N\mathfrak p}{(N\mathfrak p)^{k\sigma}}>0,
 \qquad
 E_2(\sigma)=\sum_{\mathfrak p}\sum_{k\ge1}
           \frac{k(\log N\mathfrak p)^2}
                {(N\mathfrak p)^{k\sigma}}>0.
\]
Thus $\zeta_K'/\zeta_K=-E_1$ and
$(\zeta_K'/\zeta_K)'=E_2$. Taking the logarithmic derivatives of
\eqref{eq:analytic-order-completion} now yields exactly
\[
 T_{2,K}(\delta)=U_K(\delta)
                  -\frac\delta2E_1(\sigma)
                  -\frac{\delta^2}2E_2(\sigma)
                  <U_K(\delta),
\]
which proves \eqref{eq:analytic-T2-arch}.
\end{proof}

\begin{corollary}[The quaternionic example has central order exactly two]
\label{cor:analytic-martinet-exact-two}
Let
\[
 F=\mathbb Q(\sqrt5,\sqrt{41}),\qquad
 N=F\left(\sqrt{\frac{(5+\sqrt5)(41+\sqrt{205})}{4}}\right).
\]
For this totally real tame quaternionic field of discriminant
$205^6$ from NN-A018,
\[
               \operatorname{ord}_{s=1/2}\zeta_N(s)=2.
\]
Its unique irreducible degree-two Artin factor has central order
exactly one. These statements require neither A nor M.

For the specific shifted ideal-norm series
$\zeta_{N,t}(s)=\sum_{n\ge1}a_N(n)(n+t)^{-s}$, continued as in
NN-A004, there are exactly two zero branches near $1/2$. Both
are strictly to the right for all sufficiently small real
$t>0$. More precisely, put
\[
 a_2=\frac{\zeta_N''(1/2)}2<0,\qquad
 A_N=\left(\frac{\zeta_N(3/2)}{2|a_2|}\right)^{1/2},\qquad
 d_N=-\frac{C_N A_N^2}{2}>0,
\]
where $C_N$ is the coefficient of NN-A018. Then
\begin{equation}\label{eq:analytic-martinet-two-right}
 s_\pm(t)=\frac12\pm i A_N\sqrt t+d_Nt+O(t^{3/2}),
 \qquad
 \Re s_\pm(t)-\frac12=d_Nt+O(t^2).
\end{equation}
In particular, their real velocities are
$\frac{d}{dt}\Re s_\pm(t)=d_N+O(t)>0$ on a sufficiently small
positive interval.
\end{corollary}

\begin{proof}
The field, discriminant and root-number input are those already
verified in NN-A018: Martinet, 1971, Propositions~III.1--III.2,
Theorem~IV.1 and Example~3, pp.~404--407, together with
Fr\"ohlich, 1974, Theorem~1, p.~82. The three quadratic central
values for discriminants $5,41,205$ are strictly positive by
the saved interval certificate of NN-A018. Hence the
factorization
\[
 \zeta_N(s)=\zeta(s)L(s,\chi_5)L(s,\chi_{41})
                         L(s,\chi_{205})L(s,\psi)^2
\]
and the negative symplectic root number imply
$m_N=2\operatorname{ord}_{1/2}L(s,\psi)\equiv2\pmod4$.
In particular $m_N\ge2$ and $a_{m_N}<0$; none of these steps
uses a Riemann hypothesis.

Apply the theorem at $\delta=1$. Since $r_1=8$, $r_2=0$ and
$\log D_N=6\log205$,
\[
 B_N(3/2)=\frac83+3\log205-4\log\pi+4\psi(3/4),
 \qquad
 B_N'(3/2)=-\frac{40}9+2\psi'(3/4).
\]
The reflection and duplication formulas for $\psi$, with
$\psi(1)=-\gamma_{\!E}$, give
\[
 \psi(3/4)=-\gamma_{\!E}+\frac\pi2-3\log2.
\]
For clarity, the other special value follows directly from
DLMF~5.15.1:
\[
 \psi'(3/4)=16\sum_{n\ge0}\frac1{(4n+3)^2}
                  =\pi^2-8\mathcal G,
 \qquad
 \mathcal G=\sum_{n\ge0}\frac{(-1)^n}{(2n+1)^2}.
\]
Indeed the sum of the two odd residue-class sums is
$\pi^2/8$ and their difference is $\mathcal G$.
Consequently
\begin{equation}\label{eq:analytic-martinet-U-exact}
 U_N(1)=\frac{32}9+\frac32\log205-2\log\pi
       -2\gamma_{\!E}+\pi-6\log2-\pi^2+8\mathcal G.
\end{equation}

Only coarse rational estimates are needed:
\begin{equation}\label{eq:analytic-martinet-coarse}
 \frac{157}{50}<\pi<\frac{22}7,\quad
 \log205<\frac{16}3,\quad \log\pi>\frac{57}{50},\quad
 \gamma_{\!E}>\frac12,\quad \log2>\frac{69}{100},\quad
 \mathcal G<\frac{23}{25}.
\end{equation}
Here are elementary finite checks, included so that no
floating-point constant is required. The Machin identity
$\pi=16\arctan(1/5)-4\arctan(1/239)$ follows by the tangent
addition law and the angle ranges. Alternating-series bounds
give
\[
 \begin{split}
 16\left(\frac15-\frac1{3\cdot5^3}\right)-\frac4{239}
                         &>\frac{157}{50},\\
 16\left(\frac15-\frac1{3\cdot5^3}
                         +\frac1{5\cdot5^5}\right)
 -4\left(\frac1{239}-\frac1{3\cdot239^3}\right)
                         &<\frac{22}7.
 \end{split}
\]
The positive series $\log x=2\sum_{j\ge0}y^{2j+1}/(2j+1)$,
$y=(x-1)/(x+1)$ for $x>1$, gives
\[
 \log\pi>\log(157/50)
       >2\sum_{j=0}^3\frac{(107/207)^{2j+1}}{2j+1}
       >\frac{57}{50},\qquad
 \log2>2\left(\frac13+\frac1{81}\right)
       =\frac{56}{81}>\frac{69}{100}.
\]
For the two upper logarithm bounds used here, the positive
exponential series gives the exact rational comparisons
\[
 \begin{split}
 \sum_{j=0}^{16}\frac{(16/3)^j}{j!}-205
       &=\frac{58220464777647308}{27485885585032875}>0,\\
 \sum_{j=0}^{8}\frac{(39/20)^j}{j!}-7
       &=\frac{3130848254169}{114688000000000}>0.
 \end{split}
\]
The first proves $\log205<16/3$. The elementary integral
comparison $\gamma_{\!E}>H_6-\log7$, the second comparison and
$H_6=49/20$ give $\gamma_{\!E}>1/2$. Finally, the alternating
Catalan series gives
\[
 \mathcal G<\sum_{j=0}^6\frac{(-1)^j}{(2j+1)^2}
       =\frac{23}{25}-\frac{3085982}{2029052025}
       <\frac{23}{25}.
\]
This proves every inequality in
\eqref{eq:analytic-martinet-coarse}. Substitution in
\eqref{eq:analytic-martinet-U-exact} now gives
\[
 \begin{split}
 U_N(1)&<\frac{32}9+8-\frac{57}{25}-1+\frac{22}7
            -\frac{207}{50}-\frac{24649}{2500}
            +\frac{184}{25}\\
       &=\frac{752663}{157500}
         =5-\frac{34837}{157500}<5.
 \end{split}
\]
The theorem yields $m_N<5$. Together with
$m_N\equiv2\pmod4$, this forces $m_N=2$ and the Artin order
one.

NN-A018 already proves $C_N<-3.5197<0$ using its saved
certificate. Its full local formula, now with $m=2$, gives
the first expansion in \eqref{eq:analytic-martinet-two-right}.
There are exactly two holomorphic sections after $t=u^2$.
Real coefficients and uniqueness identify these sections both
by conjugation on real $u$ and by $u\mapsto-u$; thus the odd
coefficients of either section are imaginary and its even
coefficients are real. The real part consequently has a
convergent power series in $t$, with first coefficient
$d_N>0$. This proves the improved $O(t^2)$ remainder, permits
termwise differentiation, and establishes the strict side and
velocity assertions.
\end{proof}

\paragraph{Classification and source boundary.}
The quaternionic root-number and central-existence results are
classical. Sami Omar, \emph{On Artin $L$-Functions for Octic
Quaternion Fields}, \emph{Experimental Mathematics}~10 (2001),
237--245, DOI~10.1080/10586458.2001.10504445, studies exact
central orders in these fields. The present proof does not
claim to be the first determination of the order of this field,
nor does it import a table entry that has not been inspected.
Its order bound is fully proved above and uses the existing
NN-A018 certificate only for the quadratic nonvanishing and
the second-jet sign. The comparison with an exact Omar table
entry remains a source-location audit, not a missing hypothesis
of the deduction.

\paragraph{AG and representation translation.}
The symplectic irreducible representation of the actual
$Q_8$-extension has negative functional-equation sign and,
by the proved upper bound, a length-one central analytic zero.
Its multiplicity two in the regular representation produces
the local Dedekind Artin algebra $\mathbb C[z]/(z^2)$. For the
specified shifted deformation, the base change $t=u^2$ gives
two sections with a common positive second real normal
coefficient $d_N$. Their discriminant is
$-4A_N^2t+O(t^2)$, so the cover has simple ramification at
$t=0$, is etale off that point on a sufficiently small disc,
and has the transposition monodromy. This is the concrete
finite-free discriminant criterion of SGA~1, Expos\'e~I,
section~4, Proposition~10. Positivity of the rational test is
a statement about conjugate analytic zero orbits, not
Frobenius positivity or a claimed motive for the shifted
series. The sign of the second normal jet retains the analytic
unit data that the length and the reduced support alone do not
record.


\section{An explicit cubic compositum with two excess left branches}
\label{sec:analytic-cubic-left-imbalance}

Let $N$ be the quaternionic field of NN-A018 and NN-A033.
Put $p=751$, and let $\vartheta$ be the primitive cubic
Dirichlet character modulo $p$ specified by
\[
  \vartheta(3^j)=e^{2\pi ij/3},\qquad
  \vartheta(n)=0\quad(p\mid n).
\]
Let $E$ be the fixed field of $\ker\vartheta$ in
$\mathbb Q(\zeta_{751})$, and define $K=NE$.
This explicitly specifies the field; no search over number
fields is part of the statement.

\begin{theorem}[NN-A034: an unconditional arithmetic left imbalance]
\label{thm:analytic-cubic-left-imbalance}
The extension $K/\mathbb Q$ is totally real, has Galois group
$Q_8\times C_3$, degree $24$, and absolute discriminant
\[
                   D_K=205^{18}751^{16}.
\]
There is an integer $k\ge0$ such that
\[
 \operatorname{ord}_{s=1/2}\zeta_K(s)=m=2+4k,
 \qquad \frac{\zeta_K^{(m)}(1/2)}{m!}<0,
 \qquad C_K>0.
\]
For all sufficiently small positive $t$, the shifted ideal-norm
deformation of NN-A004 has exactly $2+2k$ central branches on
the left of $\Re s=1/2$ and exactly $2k$ on the right. Thus it
has exactly two excess left branches. Its two initially vertical
branches both bend strictly left at the second Puiseux order.
All these conclusions are unconditional; neither A nor M is
used. The integer $k$, a central order of a complex Artin
twist, is not determined by this theorem.
\end{theorem}

\begin{lemma}[Four certified nonzero central values]
\label{lem:analytic-cubic-four-values}
Write $\chi_d(n)=(d/n)$ for $d=1,5,41,205$, with
$\chi_1=1$. Each $\eta_d=\chi_d\vartheta$ is primitive and
even, of conductor $q_d=751d$. Its central value lies in the
following open rectangle. Every decimal endpoint in this table
is an exact rational number.
\[
\begin{array}{c|r|c|c}
d&q_d&\Re L(1/2,\eta_d)&\Im L(1/2,\eta_d)\\\hline
1&751&(-0.0265980,-0.0265979)&(0.2188449,0.2188450)\\
5&3755&(1.2071682,1.2071683)&(-0.5141633,-0.5141632)\\
41&30791&(0.3418888,0.3418889)&(0.0415524,0.0415525)\\
205&153955&(0.8806422,0.8806423)&(2.0675988,2.0675989)
\end{array}
\]
In particular all four values are nonzero.
\end{lemma}

\begin{proof}
Primality of $751$ follows by trial division by the primes at
most $\sqrt{751}<28$. Since $750=2\cdot3\cdot5^3$ and
\[
 3^{375}\equiv750,\qquad3^{250}\equiv678,\qquad
 3^{150}\equiv80\pmod{751},
\]
the residue class of $3$ is primitive. This defines the character
exactly by its finite discrete-log table. Its parity is even
because $\vartheta(-1)=e^{2\pi i375/3}=1$. The quadratic
characters have positive fundamental discriminants $d$, and
their conductors are coprime to $751$; the product characters
are consequently primitive of conductor $q_d$ and even.

Here is the exact analytic formula and rigorous tail used for
the finite check. For any primitive even character $\eta$ of
conductor $q$, put
\[
 \epsilon_\eta=\frac1{\sqrt q}
          \sum_{a=1}^{q-1}\eta(a)e^{2\pi ia/q},\qquad
 w_q(n)=\frac{\Gamma(1/4,\pi n^2/q)}{\Gamma(1/4)\sqrt n}.
\]
The primitive Gauss sum has modulus $\sqrt q$, so
$|\epsilon_\eta|=1$. The completed functional equation has
this sign in the convention
$\Lambda(s,\eta)=(q/\pi)^{s/2}\Gamma(s/2)L(s,\eta)$;
compare DLMF~25.15.5--25.15.6. To derive the central formula,
apply Poisson summation in the residue classes to
$\Theta_\eta(x)=\sum_{n\in\mathbb Z}\eta(n)e^{-\pi n^2x/q}$.
The finite Fourier transform is
\[
 \sum_{a\bmod q}\eta(a)e^{2\pi ian/q}
              =\bar\eta(n)\tau(\eta),
\]
so
$\Theta_\eta(x)=\epsilon_\eta x^{-1/2}
\Theta_{\bar\eta}(1/x)$. Splitting its Mellin integral at $1$
gives the absolutely convergent identity
\begin{equation}\label{eq:analytic-cubic-central-sum}
 L(1/2,\eta)=\sum_{n\ge1}
            \bigl(\eta(n)+\epsilon_\eta\bar\eta(n)\bigr)w_q(n).
\end{equation}
The change of variable $u=\pi n^2x/q$ evaluates each integral
as the upper incomplete gamma function displayed in $w_q$;
dividing by $(q/\pi)^{1/4}\Gamma(1/4)$ leaves precisely
$n^{-1/2}$. This verifies every scale and phase in
\eqref{eq:analytic-cubic-central-sum}.

For $x>0$,
\[
 \Gamma(1/4,x)=\int_x^\infty u^{-3/4}e^{-u}\,du
                  \le x^{-3/4}e^{-x}.
\]
Also
$\Gamma(1/4)>\int_0^1u^{-3/4}(1-u)\,du=16/5>1$.
Consequently the omitted tail after an integer $B\ge1$ is
bounded by
\begin{equation}\label{eq:analytic-cubic-tail}
 \begin{split}
 \left|\sum_{n>B}
       (\eta(n)+\epsilon_\eta\bar\eta(n))w_q(n)\right|
 &\le 2(q/\pi)^{3/4}
       \int_B^\infty x^{-2}e^{-\pi x^2/q}\,dx\\
 &\le\frac{(q/\pi)^{7/4}}{B^3}
                         e^{-\pi B^2/q}.
 \end{split}
\end{equation}
The first comparison uses the decreasing positive summand;
the second uses $x^{-2}\le x/B^3$ for $x\ge B$ and integrates
the resulting Gaussian derivative exactly.

Only the Gauss sum of conductor $751$ is needed. The Chinese remainder
theorem, writing a residue modulo $751d$ as $751a+db$, gives
\[
 \tau(\chi_d\vartheta)
   =\chi_d(751)\vartheta(d)\tau(\chi_d)\tau(\vartheta),
 \qquad
 \epsilon_{\eta_d}
   =\chi_d(751)\vartheta(d)\epsilon_\vartheta.
\]
Here $\tau(\chi_d)=\sqrt d$, with the positive quadratic Gauss
sign verified in Gross, section~1, equations~(1.3)--(1.6),
as used in NN-A032; for $d=1$ this is the trivial identity.
This exact formula avoids any ambiguity about conjugating a
root number.

The saved script \path{work/analytic/check_cubic_compositum.py}
uses these formulas with FLINT/Arb complex balls at $60$
decimal digits. It takes
\[
 B=5(\lfloor\sqrt q\rfloor+1),\qquad
 B\in\{140,310,880,1965\},
\]
computes the full finite sum in
\eqref{eq:analytic-cubic-central-sum}, and adds the upper bound
\eqref{eq:analytic-cubic-tail} to the error in both real and
imaginary components. The respective tail bounds are strictly
less than
\[
 1.31\cdot10^{-38},\quad9.86\cdot10^{-38},\quad
 6.87\cdot10^{-37},\quad1.29\cdot10^{-36}.
\]
It checks all eight real and imaginary interval inclusions in
the table, as well as nonvanishing, and writes the complete
balls and tail bounds to
\path{work/analytic/cubic_compositum_certificate.json}.
The stated rectangles are certified enclosures, not numerical
approximations substituted for a proof.
\end{proof}

\begin{proof}[Proof of Theorem~\ref{thm:analytic-cubic-left-imbalance}]
The cyclic field $E$ exists by the cyclotomic Galois
correspondence. It has degree three and is real because
$\vartheta(-1)=1$. The only ramified prime is $751$, where the
ramification is total and tame of index three, so
$D_E=751^2$. This also follows directly from the
conductor--discriminant formula for its two nontrivial cubic
characters. The degrees $8$ and $3$ are coprime, so
$N\cap E=\mathbb Q$. Both fields are Galois, which gives
$\operatorname{Gal}(K/\mathbb Q)=Q_8\times C_3$ and degree
$24$. The discriminants $205^6$ and $751^2$ are coprime.
The integral tensor-product discriminant formula therefore
gives
\[
 D_K=D_N^3D_E^8=205^{18}751^{16}.
\]
A precise primary exposition is B.~Conrad,
\emph{Discriminant of composite fields}, Math.~154 handout,
p.~1: its displayed formula and trace-tensor determinant
calculation apply with exactly these degree and coprimality
hypotheses. Equivalently, the tame inertia indices $4,4,3$ at
$5,41,751$ give exponents $24(1-1/4)=18$ and
$24(1-1/3)=16$ directly.

Let $\psi$ denote the unique degree-two irreducible character
of $Q_8$, inflated where appropriate. The irreducible
characters of the product group are
\[
 \chi_d\vartheta^j\quad(d=1,5,41,205;\ j=0,1,2),
 \qquad
 \psi\vartheta^j\quad(j=0,1,2).
\]
The former have degree one and the latter degree two.
Decomposing the regular representation, Artin formalism gives
the exact factorization, including ramified Euler factors,
\begin{equation}\label{eq:analytic-cubic-factorization}
 \zeta_K(s)=
  \prod_{d\in\{1,5,41,205\}}\prod_{j=0}^2
                L(s,\chi_d\vartheta^j)
  \prod_{j=0}^2 L(s,\psi\vartheta^j)^2.
\end{equation}
For clarity, the regular-representation Euler identity does
not discard bad primes: on inertia invariants, Frobenius acts
by cycles corresponding to the prime ideals over the given
rational prime, with their residue degrees. Its determinant
is therefore the product of the corresponding Dedekind Euler
factors. Direct-sum decomposition of that same invariant
space gives \eqref{eq:analytic-cubic-factorization} at every
prime.

There is no unproved Artin holomorphy input. Choose $C_4\subset
Q_8$ and a faithful linear character $\lambda$ of $C_4$.
Since its nontrivial conjugate is $\lambda^{-1}$,
$\operatorname{Ind}_{C_4}^{Q_8}\lambda=\psi$. Thus
\[
 \psi\vartheta^j=
 \operatorname{Ind}_{C_4\times C_3}^{Q_8\times C_3}
                       (\lambda\boxtimes\vartheta^j).
\]
Artin induction invariance identifies each $L$-function with
the nontrivial finite-order Hecke $L$-function over the
quadratic fixed field of $C_4\times C_3$. Such a function is
entire: the classical analytic input is, for example,
B.~Poonen, \emph{Tate's thesis}, section~5.11,
Theorem~5.22(b), p.~38, together with the Galois-character
construction in Example~5.14, p.~32. A nontrivial finite-order
character cannot be a norm power, the only exceptional
component in that theorem. The local induction identity is
the determinant identity for induced cosets, and holds at the
ramified places as well; it is the same unconditional
monomial argument already used for $\psi$ in NN-A018.

NN-A033 proves
$\operatorname{ord}_{1/2}L(s,\psi)=1$. Put
$k=\operatorname{ord}_{1/2}L(s,\psi\vartheta)\ge0$.
Conjugating the Artin Euler factors gives
\[
 L(s,\psi\bar\vartheta)
       =\overline{L(\bar s,\psi\vartheta)},
\]
because $\psi$ is real valued. This identity extends from
$\Re s>1$ to the entire continuations and shows that the
two cubic Artin twists have the same order $k$ at the center.
The four nonzero values in Lemma~\ref{lem:analytic-cubic-four-values}
and their conjugates dispose of the eight nonreal
degree-one factors. The remaining factors are $\zeta$ and
the three positive quadratic central values already certified
in NN-A018; $\zeta(1/2)<0$. Hence
\eqref{eq:analytic-cubic-factorization} gives exactly
\[
                  m=2+4k.
\]
It also determines the leading sign. On the real axis the
nonreal character factors occur as absolute squares. If
$b_1$ and $c_k$ are the first nonzero coefficients of
$L(s,\psi)$ and $L(s,\psi\vartheta)$, respectively, their
contribution to the leading coefficient is
$b_1^2|c_k|^4>0$. Multiplying by the negative nonzero product
of the degree-one central values gives
$a_m=\zeta_K^{(m)}(1/2)/m!<0$.

Finally the root discriminant is
\[
 D_K^{1/24}=205^{3/4}751^{2/3}>4371.
\]
This comparison is an exact integer inequality: raising to
the twelfth power gives
\[
 205^9 751^8-4371^{12}
    =16062670137855473202693706441001475807620284>0.
\]
The conductor criterion proved and interval-certified in
NN-A018 therefore gives $C_K>0$. Applying its full second
Puiseux formula with $a_m<0$ and $m=2+4k$ yields
\[
 \#\{\Re s_j(t)<1/2\}=\frac{m+2}2=2+2k,
 \quad
 \#\{\Re s_j(t)>1/2\}=\frac{m-2}2=2k
\]
on a common sufficiently small interval $t>0$. The two
vertical coefficients are $\pm iA$, and their second normal
coefficient is $-C_KA^2/m<0$. This proves all the assertions.
\end{proof}

\paragraph{Scope and arithmetic consequence.}
Together NN-A033 and NN-A034 supply explicit number fields
with opposite central side imbalances under the same
ideal-norm shift. The latter field is fixed before the
finite certification and requires no nonabelian GRH. The
unknown $k$ affects the total number of branches but not the
excess of two on the left. No noncentral sign-density law,
Riemann-zeta simplicity theorem or heat-collision statistic
is inferred from these central examples.

\paragraph{AG and representation translation.}
Away from $5\cdot41\cdot751$, the specified compositum gives
an actual finite etale $Q_8\times C_3$-cover of
$\operatorname{Spec}\mathbb Z[1/(5\cdot41\cdot751)]$.
The external tensor-product decomposition of its regular
representation gives the displayed Artin factorization.
Complex conjugation pairs the two cubic degree-two twists,
so their contribution to the central analytic length is
$4k$, while the old symplectic factor contributes two.
After $t=u^{2+4k}$, the local analytic Cartier divisor has
$2+4k$ normalized sections, with the two vertical first jets
having strictly negative second real normal component.
These are actual representation multiplicities and local
Artin lengths. The exterior-power/representation-ring
formalism is SGA~6, Expos\'e~V, Example~2.2(b),
Lemma~2.2.1; the local finite-cover/discriminant dictionary
is SGA~1, Expos\'e~I, section~4, Proposition~10. The
specified analytic norm shift is not asserted to arise
from an arithmetic deformation of the cover.


\section{Multiple arithmetic zeros cannot remain tangent through the second Puiseux jet}
\label{sec:analytic-multiple-zero-bending}

This section treats zeros of multiplicity $m\ge2$. The simple
case has an additional deformation term, displayed at the end;
no conclusion about simple-zero transversality is inferred.
All statements at a specified critical-line zero are
unconditional. Assumption A makes them apply to every multiple
nontrivial zero of $\zeta$ or a primitive Dirichlet $L$-function.

\begin{lemma}[A positive logarithmic derivative at one horizontal unit]
\label{lem:analytic-reflected-logderivative}
Let $\Lambda$ be an entire order-one completion with zeros in
$0\le\Re s\le1$, infinitely many zeros, and reflection
\[
 \Lambda(s)=\epsilon\overline{\Lambda(1-\bar s)},
 \qquad |\epsilon|=1.
\]
Suppose $\rho=1/2+i\gamma$ is a zero of order $m\ge1$.
Then, unconditionally,
\begin{equation}\label{eq:analytic-reflected-positive}
 \Re\frac{\Lambda'}{\Lambda}(\rho+1)
   =m+\sum_{\rho'\ne\rho}
       \frac{3/2-\Re\rho'}
       {(3/2-\Re\rho')^2+(\Im\rho'-\gamma)^2}>m.
\end{equation}
Multiplicity is retained and the displayed real series is
absolutely convergent.
\end{lemma}

\begin{proof}
Put $F(z)=\Lambda(\rho+z)$ and write its first two coefficients
as $b_mz^m+b_{m+1}z^{m+1}+O(z^{m+2})$. Since
$1-\bar\rho=\rho$, reflection becomes
$F(z)=\epsilon\overline{F(-\bar z)}$. Therefore
\[
 b_m=\epsilon(-1)^m\bar b_m,
 \qquad b_{m+1}=\epsilon(-1)^{m+1}\bar b_{m+1},
 \qquad
 \Re\frac{b_{m+1}}{b_m}=0.
\]
The genus-one Hadamard product has the form
\[
 F(z)=z^m e^{az+b}\prod_{\alpha\ne0}
           E_1(z/\alpha),\qquad
 E_1(w)=(1-w)e^w,
 \quad \alpha=\rho'-\rho.
\]
All multiplicities are included and
$\sum |\alpha|^{-2}<\infty$. Since each $E_1$ has zero
linear coefficient at the origin, the logarithmic derivative
of $F(z)/z^m$ at zero is $a=b_{m+1}/b_m$. Thus
$\Re a=0$; this is the required control of the exponential
factor, not an assumption about its sign.

At $z=1$, which is not a zero, logarithmic differentiation gives
\[
 \frac{F'(1)}{F(1)}=m+a+
     \sum_{\alpha\ne0}\left(\frac1{1-\alpha}
                                      +\frac1\alpha\right).
\]
The combined complex terms converge absolutely. Moreover
$|\Re\alpha|\le1/2$, so
$\sum|\Re(1/\alpha)|<\infty$. Reflection pairs $\alpha$ with
$-\bar\alpha$, with the same multiplicity, and
\[
 \Re\left(\frac1\alpha+\frac1{-\bar\alpha}\right)=0.
\]
For a purely imaginary zero this cancellation holds termwise.
Taking real parts therefore leaves just
$m+\sum\Re(1/(1-\alpha))$. Each summand is
\[
 \frac{1-\Re\alpha}
 {(1-\Re\alpha)^2+(\Im\alpha)^2}>0,
\]
because $1-\Re\alpha\ge1/2$. Its tail is bounded by a fixed
multiple of $\sum|\alpha|^{-2}$, proving absolute convergence.
There are noncentral zeros, so the inequality is strict.
Substituting $\alpha=\rho'-\rho$ proves
\eqref{eq:analytic-reflected-positive}.
\end{proof}

\begin{theorem}[NN-A035: strict left bending at multiple Dirichlet zeros]
\label{thm:analytic-multiple-dirichlet-bending}
Let $L(s)=L(s,\chi)$ for a nonprincipal primitive Dirichlet
character, real or complex, of parity $a\in\{0,1\}$ and
conductor $q$. Suppose $\rho=1/2+i\gamma$ is a zero of
multiplicity $m\ge2$. Write
\[
 L(\rho+z)=A_mz^m+A_{m+1}z^{m+1}+O(z^{m+2}),
 \qquad
 D_\rho=\frac1\rho+\frac{L'}L(\rho+1)
                              -\frac{A_{m+1}}{A_m}.
\]
Then
\begin{equation}\label{eq:analytic-multiple-D-positive}
 \Re D_\rho>m+h_a(\gamma)>m,
\end{equation}
where the exact positive functions are
\[
 h_1(\gamma)=\frac{\pi}{2\cosh(\pi\gamma)},\qquad
 h_0(\gamma)=\frac1{2(\gamma^2+1/4)}
                      -\frac{\pi}{2\cosh(\pi\gamma)}>0.
\]
For the shifted deformation of NN-A004, the $m$ normalized
branches after $t=u^m$ satisfy
\begin{equation}\label{eq:analytic-multiple-full-jets}
 s_j(t)=\rho+c_jt^{1/m}
                  +\frac{D_\rho}{m}c_j^2t^{2/m}
                  +O(t^{3/m}),\qquad
 c_j^m=\frac{\rho L(\rho+1)}{A_m}\ne0.
\end{equation}
Every branch has a nonzero real normal coefficient by its
second Puiseux order. Specifically, if $\Re c_j=0$ and
$c_j=i\beta_j$, then
\[
 \Re s_j(t)-\frac12
       =-\frac{\beta_j^2}{m}\Re D_\rho\,t^{2/m}
                         +O(t^{3/m})<0
\]
for all sufficiently small positive $t$.
If vertical first directions occur, the exact left/right
counts are $(m/2+1,m/2-1)$ for even $m$, and
$((m+1)/2,(m-1)/2)$ for odd $m$.
\end{theorem}

\begin{proof}
Use the exact completion
\[
 Q_\chi(s)=\left(\frac q\pi\right)^{(s+a)/2}
                         \Gamma((s+a)/2),
 \qquad \Lambda(s)=Q_\chi(s)L(s,\chi).
\]
The primitive functional equation and conjugation give
$\Lambda(s)=\epsilon_\chi
\overline{\Lambda(1-\bar s)}$ even when $\chi$ is complex;
the root number has modulus one. These are DLMF~25.15.5--25.15.6.
The theta integral as in NN-A032, now with the conjugate theta
series on the reflected half, gives the entire order-one
completion. The classical critical-strip and infinite-zero
inputs are DLMF~25.15(ii); cancellation at the endpoint uses
$L(1,\chi)\ne0$, DLMF~25.15.9, just as in NN-A032.
Thus Lemma~\ref{lem:analytic-reflected-logderivative} applies.

The coefficient ratio in the completed germ is
\[
 \frac{b_{m+1}}{b_m}
          =\frac{A_{m+1}}{A_m}+\frac{Q_\chi'}{Q_\chi}(\rho).
\]
Its real part is zero by reflection. Consequently
\begin{equation}\label{eq:analytic-multiple-D-completed}
 \Re D_\rho=
 \Re\frac{\Lambda'}\Lambda(\rho+1)
 +\Re\left(\frac1\rho+
       \frac{Q_\chi'}{Q_\chi}(\rho)
       -\frac{Q_\chi'}{Q_\chi}(\rho+1)\right).
\end{equation}
The conductor term cancels exactly. It remains to evaluate
the real gamma difference, retaining its height dependence.

Put $z=(\rho+a)/2$. Since $\Re\rho=1/2$, one has
$\overline{(\rho+1)/2}=1-\rho/2$. The digamma reflection
identity $\psi(1-z)-\psi(z)=\pi\cot(\pi z)$, DLMF~5.5.4,
therefore gives
\[
 \Re\left(\frac12\psi(\rho/2)
                   -\frac12\psi((\rho+1)/2)\right)
             =-\frac{\pi}{2\cosh(\pi\gamma)}.
\]
Indeed $\Re\cot(\pi/4+i\pi\gamma/2)=1/\cosh(\pi\gamma)$,
as follows by writing sine and cosine in exponentials.
This yields $h_0$ after adding
$\Re(1/\rho)=1/[2(\gamma^2+1/4)]$.
For odd parity, the recurrence (DLMF~5.5.2)
$\psi(1+\rho/2)=\psi(\rho/2)+2/\rho$ reverses the preceding
difference after the same $1/\rho$ term is added, giving
$h_1=\pi/(2\cosh(\pi\gamma))$.
The positivity of $h_0$ is exact: from $2<\pi<4$,
\[
 \cosh(\pi\gamma)\ge1+\frac{\pi^2\gamma^2}{2}
              >\pi(\gamma^2+1/4).
\]
This proves both positive functions. Combining them with
\eqref{eq:analytic-reflected-positive} proves
\eqref{eq:analytic-multiple-D-positive}.

For completeness, the local deformation coefficient is
derived rather than imported from a simple-zero formula.
The shifted expansion gives
\[
 L_t(\rho+z)=A_mz^m+A_{m+1}z^{m+1}
       -t(B_0+B_1z)+O(z^{m+2})+O(tz^2)+O(t^2),
\]
where $B_0=\rho L(\rho+1)$ and
$B_1=L(\rho+1)+\rho L'(\rho+1)$. Since $\Re(\rho+1)=3/2$,
the Euler product makes $B_0\ne0$. Put $t=u^m$ and
$z=c_ju+d_ju^2+O(u^3)$. The coefficient of $u^m$ gives
$A_mc_j^m=B_0$, with distinct nonzero roots. The coefficient
of $u^{m+1}$ gives
\[
 mA_mc_j^{m-1}d_j+A_{m+1}c_j^{m+1}-B_1c_j=0,
\]
because $2m\ge m+2$ for $m\ge2$. Thus
$d_j=c_j^2D_\rho/m$, proving
\eqref{eq:analytic-multiple-full-jets}. The analytic implicit
function theorem after division by $u^m$ justifies these
convergent sections and their common error estimate.

For a vertical $c_j=i\beta_j$, its second real coefficient
is strictly negative by \eqref{eq:analytic-multiple-D-positive}.
Every other branch is decided by its nonzero first real
coefficient. If $m$ is even, vertical directions come as the
opposite pair; the remaining roots are antipodal and split
equally between left and right. If $m$ is odd, there is exactly
one vertical root, and the other roots pair with opposite
real parts by their angular positions $\pi/2+2\pi j/m$ (or
the corresponding lower vertical starting angle). Adding
the left-bending vertical branches gives the stated counts.
\end{proof}

\begin{corollary}[Dedekind zeros at sufficient height, and every multiple zeta zero under A]
\label{cor:analytic-multiple-dedekind-bending}
Let $K$ have signature $(r_1,r_2)$ and degree $n$. At a
critical-line zero $\rho=1/2+i\gamma$ of $\zeta_K$ with
order $m\ge2$, the same $D_\rho$ and Puiseux expansion obey
\begin{equation}\label{eq:analytic-dedekind-D-lower}
 \Re D_\rho>m-
 \left(\frac{3}{2(\gamma^2+9/4)}
       +\frac{r_1\pi}{2\cosh(\pi\gamma)}
       +\frac{r_2}{2(\gamma^2+1/4)}\right).
\end{equation}
In particular, the sufficient condition
\[
 \gamma^2\ge\frac1m
       \left(\frac{r_2+3}{2}+\frac{r_1}{\pi}\right)
\]
already implies $\Re D_\rho>0$ and strict left bending of
every vertical first direction. If
\[
 \gamma^2\ge\frac{r_2+3}{2}+\frac{r_1}{\pi},
\]
then the stronger bound $\Re D_\rho>m-1>0$ holds. The simpler sufficient bound
$\gamma^2\ge n+2$ also works.

For $K=\mathbb Q$, the existing NN-A014 interval certificate
implies that every critical-line zero has $\gamma^2>8$.
Hence every multiple critical-line zeta zero satisfies the
stronger bound
\[
                    \Re D_\rho>m-\frac{11}{48}>0.
\]
Under A this treats all multiple nontrivial zeta zeros, without
assuming that any such zero exists.
\end{corollary}

\begin{proof}
Apply Lemma~\ref{lem:analytic-reflected-logderivative} to the
entire pole-removed Dedekind completion of NN-A033. The
local coefficient relation and
\eqref{eq:analytic-multiple-D-completed} still hold. The
pole-removal factors contribute
\[
 \Re\left(\frac1\rho+\frac1{\rho-1}
                                  -\frac1{\rho+1}\right)
                  =-\frac{3}{2(\gamma^2+9/4)}.
\]
Each real-place gamma factor contributes
$-\pi/(2\cosh(\pi\gamma))$, as just calculated. Each
complex-place factor contributes
$\Re(\psi(\rho)-\psi(\rho+1))=-\Re(1/\rho)$.
This proves \eqref{eq:analytic-dedekind-D-lower} with the
exact stated completion and no location hypothesis beyond
the specified zero being on the line.

For $\gamma\ne0$, the bracket is strictly less than
\[
 \frac1{\gamma^2}
        \left(\frac{r_2+3}{2}+\frac{r_1}{\pi}\right),
\]
using $\cosh(\pi\gamma)>\pi^2\gamma^2/2$. Comparing this
upper bound first with $m$ and then with $1$ proves both
thresholds. Since $\pi>3$ and $n=r_1+2r_2$,
the numerator is less than $n/3+3/2<n+2$, giving the
simpler sufficient threshold.

For the zeta assertion, reuse the positive theta integral
and exact normalization $H_0(x)=\xi(1/2+ix/2)/8$ of
NN-A006. The full saved Arb ball of NN-A014 is
\[
 A_0=\frac{\xi''(0)}{4\xi(1/2)}
       \in (0.0117395,0.0117396)\subset(0,1/16).
\]
No new computation is required for this wider rational
enclosure. If $c_1$ is the normalized second theta moment,
the positive coefficient sum in NN-A014 gives
$2c_1\le A_0<1/16$. The elementary inequality
$\cos v\ge1-v^2/2$ therefore gives, unconditionally,
\[
 H_0(x)\ge H_0(0)(1-c_1x^2)>0\qquad(x^2\le32).
\]
This is a statement about the integral and does not use a
real-zero product or RH. Thus a critical-line zero has
$x=2\gamma$ with $\gamma^2>8$. In the case $r_1=1,r_2=0$,
the bracket in \eqref{eq:analytic-dedekind-D-lower} is
strictly less than
$((3/2)+(1/3))/8=11/48$. The claimed lower bound and all
second-jet conclusions follow.
\end{proof}

\paragraph{Why the simple-zero gate remains open.}
For $m=1$, write $v=\rho L(\rho+1)/L'(\rho)$. Direct
coefficient matching instead gives
\begin{equation}\label{eq:analytic-simple-extra-term}
 s(t)=\rho+vt+
 \left(v^2D_\rho-
       \frac{v(\rho+1)L(\rho+2)}{2L(\rho+1)}\right)t^2
       +O(t^3).
\end{equation}
The extra term is the contribution of
$t^2\rho(\rho+1)L(\rho+2)/2$ in the shifted series.
If $v$ is purely imaginary, positivity of $\Re D_\rho$
controls only the first term in the parentheses. The real
part of the second term has not been controlled here.
Thus this theorem resolves the second normal jet at multiple
zeros, and neither proves $\Re v\ne0$ at simple zeros nor
assigns a universal bending direction to a simple tangent.

\paragraph{AG translation.}
At a length-$m$ critical-line analytic zero, completed
reflection puts the first reduced logarithmic coefficient
in $i\mathbb R$. The genus-one exponential and reflected
zero pairs make this local Real constraint compatible with
the strictly positive real logarithmic derivative one unit
to the right. After $t=u^m$, the finite analytic Cartier
divisor normalizes into $m$ sections. For $m\ge2$ the
specified shift has no $t^2$ contribution at the second
Puiseux order, so every vertical tangent acquires the
proved nonzero left normal coefficient. The local algebra
is $\mathbb C[z]/(z^m)$; its normalization and ramified
cover have the standard differentials/discriminant
interpretation of SGA~1, Expos\'e~I, section~3,
Proposition~1 and section~4, Proposition~10. This is a
second-jet theorem for the given analytic deformation,
not a purity statement or an arithmetic realization of
the shifted denominators.


\section{A positive proportion of actual zeta zeros have a definite normal velocity}
\label{sec:analytic-arithmetic-transverse-proportion}

All zero counts in this section use positive ordinates and
count multiplicity unless a simple-zero condition is written
explicitly. Put
\[
 N(T)=\sum_{0<\gamma\le T}1,\qquad
 N_s(T)=\#\{\rho=1/2+i\gamma:0<\gamma\le T,
                                      \ \zeta'(\rho)\ne0\}.
\]
Thus a zero of multiplicity $m$ contributes $m$ to $N(T)$ and
nothing to $N_s(T)$ when $m>1$.

\begin{theorem}[NN-A036: a quantitative arithmetic transversality proportion]
\label{thm:analytic-arithmetic-transverse-proportion}
Assume A. At each simple zero, let
\[
 v_\rho=\frac{\rho\zeta(\rho+1)}{\zeta'(\rho)},\qquad
 r_\rho=\frac{|\Re v_\rho|}{|v_\rho|}\in[0,1],\qquad
 c=\frac1{2\zeta(3)}-\frac13>\frac1{15}.
\]
Then
\begin{equation}\label{eq:analytic-normal-mean}
 \liminf_{T\to\infty}\frac1{N(T)}
                 \sum_{\substack{0<\gamma\le T\\\rho\ {
m simple}}}
                           r_\rho\ \ge c.
\end{equation}
In particular, at least a lower proportion $c$ of all zero
occurrences are simple with $\Re v_\rho\ne0$.

For $0<\eta<3c/2$, let
\[
 G_\eta(T)=\#\{\rho\ {
m simple}:0<\gamma\le T,
                              \ |\Re v_\rho|\ge\eta|v_\rho|\}.
\]
The bounds with their two distinct denominators are
\begin{align}
 \liminf_{T\to\infty}\frac{G_\eta(T)}{N(T)}
       &\ge\frac{c-(2/3)\eta}{1-\eta},
         \label{eq:analytic-total-transverse-density}\\
 \liminf_{T\to\infty}\frac{G_\eta(T)}{N_s(T)}
       &\ge\frac{(3/2)c-\eta}{1-\eta}.
         \label{eq:analytic-simple-transverse-density}
\end{align}
For example, taking $\eta=1/30$ gives the convenient rational
consequences
\[
 \liminf\frac{G_{1/30}(T)}{N(T)}\ge\frac4{87},\qquad
 \liminf\frac{G_{1/30}(T)}{N_s(T)}\ge\frac2{29}.
\]
These are lower bounds for the magnitude of the relative
normal component. They assign neither its sign nor a half
proportion to left and right labels, and do not establish
transversality at every simple zero.
\end{theorem}

\begin{lemma}[The exterior Euler phase averaged at the zeros]
\label{lem:analytic-exterior-phase-average}
Under RH, for each fixed $\sigma>1$,
\begin{equation}\label{eq:analytic-exterior-ratio-mean}
 \lim_{T\to\infty}\frac1{N(T)}
       \sum_{0<\gamma\le T}
       \frac{\zeta(\sigma+i\gamma)}
            {\zeta(\sigma-i\gamma)}
                     =\frac1{\zeta(2\sigma)}.
\end{equation}
This is an unweighted zero average with multiplicity, not a
joint independence assertion involving the Hardy phase.
\end{lemma}

\begin{proof}
For every fixed $x>1$, Landau's formula is
\[
 \sum_{0<\gamma\le T}x^\rho
      =-\frac{T}{2\pi}\Lambda(x)+O_x(\log T),
\]
where $\Lambda(x)=\log p$ when $x=p^k$ with integer $k\ge1$,
and is zero for other real $x>1$. A primary locator is
S.~M.~Gonek, \emph{A Formula of Landau and Mean Values of
$\zeta(s)$}, author-hosted PDF, p.~1, equation~(1);
its uniform extension is Theorem~1, equation~(3).
The fixed-$x$ formula is all that is used here.
Under RH, $x^\rho=\sqrt x\,e^{i\gamma\log x}$ and
$N(T)\sim T\log T/(2\pi)$. Thus
\begin{equation}\label{eq:analytic-fixed-frequency-zero-average}
 \frac1{N(T)}\sum_{0<\gamma\le T}e^{i\gamma\log x}\longrightarrow0
                    \qquad(x>0,\ x\ne1).
\end{equation}
For $0<x<1$ use complex conjugation and $1/x>1$.

For $\sigma>1$, the following double series is absolutely and
uniformly convergent in real $t$:
\[
 \frac{\zeta(\sigma+it)}{\zeta(\sigma-it)}
 =\sum_{n,m\ge1}\frac{\mu(m)}{(nm)^\sigma}
                                  e^{it\log(m/n)}.
\]
Indeed the absolute coefficient sum is at most
$\zeta(\sigma)^2<\infty$. Each normalized zero average of a
term has modulus at most one, so dominated convergence for
this summable family applies to
\eqref{eq:analytic-fixed-frequency-zero-average}.
Only $m=n$ survives, and its sum is
$\sum_{n\ge1}\mu(n)n^{-2\sigma}=1/\zeta(2\sigma)$.
This proves \eqref{eq:analytic-exterior-ratio-mean} without
requiring uniformity of Landau's error in $x$.
\end{proof}

\begin{proof}[Proof of Theorem~\ref{thm:analytic-arithmetic-transverse-proportion}]
Let $\theta(t)$ be the usual Hardy--Riemann--Siegel phase and
use the Euler-product branch
\[
 \phi(t)=\arg\zeta(3/2+it),\qquad
 \delta(t)=\arctan\frac1{2t},\qquad
 \psi(t)=\phi(t)-\delta(t).
\]
The first necessary phase average is already a proved
classical statement, not full Gram equidistribution:
\begin{equation}\label{eq:analytic-theta-first-mode}
 \frac1{N(T)}\sum_{0<\gamma\le T}e^{2i\theta(\gamma)}
                                      \longrightarrow0.
\end{equation}
The exact locator is J.~Arias de Reyna,
\emph{On the distribution (mod 1) of the normalized zeros of
the Riemann Zeta-function}, arXiv:1403.7126v1,
Corollary~3.2, equation~(13), with $\kappa=1$; its proof gives
$O(1/\log T)$ in this case. That result sums
$\rho=1/2+i\alpha$ without assuming RH. Here A makes
$\alpha=\gamma$ real, with precisely the multiplicity
convention used in $N(T)$.

Since $e^{2i\phi(t)}=\zeta(3/2+it)/\zeta(3/2-it)$,
Lemma~\ref{lem:analytic-exterior-phase-average} gives
\[
 \frac1{N(T)}\sum_{0<\gamma\le T}e^{2i\phi(\gamma)}
                                  \longrightarrow\frac1{\zeta(3)}.
\]
The replacement of $\phi$ by $\psi$ has vanishing average
error. To see this with no interchange of divergent sums,
fix $Y>0$: the finitely many zeros below $Y$ contribute
$o_T(1)$, while for $\gamma>Y$ the difference of the two
unit phases is at most $2\delta(Y)$. Let $T\to\infty$ and
then $Y\to\infty$. Taking real parts, the two limits are
\begin{equation}\label{eq:analytic-two-separate-real-means}
 \langle\cos2\theta\rangle_T\longrightarrow0,
 \qquad
 \langle\cos2\psi\rangle_T\longrightarrow\frac1{\zeta(3)}.
\end{equation}
No joint distribution or decorrelation has been used.

Define at every zero occurrence, including multiple zeros,
the bounded phase observable
$f(\gamma)=|\cos(\theta(\gamma)+\psi(\gamma))|$.
The pointwise trigonometric identity gives
\[
 \cos2\theta+\cos2\psi
   =2\cos(\theta+\psi)\cos(\theta-\psi)
   \le2|\cos(\theta+\psi)|=2f.
\]
Consequently
\begin{equation}\label{eq:analytic-all-phase-lower}
 \liminf_{T\to\infty}\langle f\rangle_T
                              \ge\frac1{2\zeta(3)}.
\end{equation}
This deterministic inequality is why the separate first
moments suffice for the present lower bound, although they
do not prove the parity-weighted half-label law.

At a simple zero, $Z(t)=e^{i\theta(t)}\zeta(1/2+it)$ is
real, $Z'(\gamma)\ne0$, and
\[
 \zeta'(\rho)=-i e^{-i\theta(\gamma)}Z'(\gamma),\qquad
 v_\rho=\frac{i\rho e^{i\theta(\gamma)}
                   \zeta(3/2+i\gamma)}{Z'(\gamma)}.
\]
For $\gamma>0$, $\arg(i\rho)=\pi-\delta(\gamma)$.
The real sign of $Z'(\gamma)$ only adds a multiple of $\pi$
to this argument. Therefore the corrected exact velocity
formula of NN-A001 gives
\begin{equation}\label{eq:analytic-phase-is-relative-normal}
                   r_\rho=f(\gamma)
                       \qquad(\rho\text{ simple}).
\end{equation}
The Euler product makes the numerator of $v_\rho$ nonzero,
so the displayed normalization is well defined.

The remaining imported input is precisely Montgomery's
proved simple-proportion bound under RH:
\begin{equation}\label{eq:analytic-montgomery-simple-proportion}
                   N_s(T)\ge(2/3+o(1))N(T).
\end{equation}
The source is \emph{The Pair Correlation of Zeros of the
Zeta Function}, Corollary~2, printed p.~182, equation~(7).
Its proof in section~4, pp.~191--192, explicitly states
that the zero sums count multiplicity; a multiplicity-$m$
zero is entered $m$ times and contributes $m^2$ to the
equal-ordinate pair sum. Thus
\eqref{eq:analytic-montgomery-simple-proportion} has exactly
the denominator $N(T)$ used in both phase averages. This is
the already catalogued classical input NT-Z02, not a new
simplicity claim or a full pair-correlation conjecture.

Because $0\le f\le1$, removing the nonsimple occurrences
from \eqref{eq:analytic-all-phase-lower} loses at most
$(N(T)-N_s(T))/N(T)\le1/3+o(1)$. Together with
\eqref{eq:analytic-phase-is-relative-normal}, this proves
\eqref{eq:analytic-normal-mean}. Since $r_\rho\le1$ and
vanishes at a tangent simple zero, it also proves the lower
proportion $c$ of simple transverse zeros.

For the sharper threshold bounds, put
$p_T=N_s(T)/N(T)$ and $g_T=G_\eta(T)/N(T)$. The pointwise
upper bounds on the three subsets (nonsimple, simple below
threshold, and simple above threshold) give
\[
 \langle f\rangle_T
       \le (1-p_T)+\eta(p_T-g_T)+g_T
       =1-(1-\eta)p_T+(1-\eta)g_T.
\]
Hence
\[
 g_T\ge p_T-\frac{1-\langle f\rangle_T}{1-\eta}.
\]
Taking lower limits using $p_T\ge2/3+o(1)$ proves
\eqref{eq:analytic-total-transverse-density}. Dividing the
same inequality by $p_T$, which is bounded away from zero,
and using
\[
 \frac{1-\langle f\rangle_T}{p_T}
       \le\frac{1-1/(2\zeta(3))+o(1)}{2/3+o(1)}
\]
proves \eqref{eq:analytic-simple-transverse-density}.
There is no hidden conversion from a distinct-zero count.

Finally,
\[
 \zeta(3)=1+\frac18+\sum_{n\ge3}\frac1{n^3}
       <1+\frac18+\int_2^\infty x^{-3}\,dx=\frac54,
\]
so $c>2/5-1/3=1/15$. Inserting $\eta=1/30$ in the two
exact bounds gives respectively a number greater than
$4/87$ and a number greater than $2/29$, proving the stated
rational consequences without a numerical constant
calculation.
\end{proof}

\paragraph{Scope.}
The theorem supplies a quantitative lower density of actual
arithmetic simple zeros at which the normal velocity is a
fixed fraction of the full velocity. It uses A, the single
proved Hardy-phase mode, the fixed-frequency Landau formula,
and the classical $2/3$ simplicity proportion. It does not
use a normalized CLT, a joint phase law, full GUE, or an
assumption that all zeros are simple. The positive and
negative signs may remain correlated with index parity;
the half-label law and the exclusion of every simple tangent
remain open here.

\paragraph{AG and harmonic-analysis translation.}
The zero divisor is counted with its actual local lengths.
At a reduced point, the equation $\zeta_t(s)=0$ gives
$d\zeta_t=\zeta'(\rho)\,ds-\rho\zeta(\rho+1)\,dt$;
its section has tangent $\partial_t+v_\rho\partial_s$.
The observable $r_\rho$ is the absolute normalized real
normal component of this tangent relative to the critical
line. The theorem gives a positive lower counting density
on that reduced locus, retaining all multiplicities in the
ambient denominator. The K\"ahler-differential local model
is the one of SGA~1, Expos\'e~I, section~3, Proposition~1.
For the exterior phase, the characters $x\mapsto
e^{i\gamma\log x}$ are characters of the multiplicative
group $\mathbb Q_{>0}$, whose compact dual is the product
of prime circles. Landau's formula identifies the relevant
Haar means, and absolute Dirichlet convergence makes the
ratio's constant term $\sum\mu(n)n^{-3}=1/\zeta(3)$.
This is an actual harmonic-analysis computation, not an
asserted Frobenius equidistribution theorem for the zeros
or an arithmetic family realizing the denominator shift.


\section{Joint phase cancellation and primitive Dirichlet transversality}

The two marginal phase means used in NN-A036 can be replaced by a
proved mixed mean. This also answers the fixed primitive Dirichlet
analogue. Throughout this section a character and its conductor are
fixed as the height tends to infinity. No uniformity over conductors,
family average, full GUE conjecture, or simplicity assumption is used.
The character of conductor one denotes the Riemann zeta function.

\begin{theorem}[NN-A037: a joint arithmetic mode and definite normal motion]
Assume A. Let $L(s,\chi)$ be any fixed primitive Dirichlet $L$-function,
including $\zeta$. Write $a\in\{0,1\}$ for its parity and put
\[
 Q_\chi(s)=(q/\pi)^{(s+a)/2}\Gamma((s+a)/2),\qquad
 \Lambda_\chi(s)=Q_\chi(s)L(s,\chi),\qquad
 \Lambda_\chi(s)=\epsilon_\chi\overline{\Lambda_\chi(1-\bar s)}.
\]
Choose a square root of $\epsilon_\chi$ and the corresponding continuous
real phase $\theta_\chi$ for which
$Z_\chi(t)=e^{i\theta_\chi(t)}L(1/2+it,\chi)$ is real. Let
$N_\chi(T)$ count positive-ordinate nontrivial zeros with multiplicity,
and $N_{\chi,s}(T)$ count simple zeros. At a simple zero $\rho=1/2+i\gamma$
the specified norm shift has velocity
\[
 v_{\rho,\chi}=
 \frac{\rho L(\rho+1,\chi)}{L'(\rho,\chi)},\qquad
 r_{\rho,\chi}=\frac{|\Re v_{\rho,\chi}|}{|v_{\rho,\chi}|}.
\]
For every $\sigma>1$, let
$\phi_{\chi,\sigma}(t)=\arg L(\sigma+it,\chi)$ be the Euler-product
argument. Then
\begin{equation}\label{eq:A037-joint}
 \frac1{N_\chi(T)}\sum_{0<\gamma\le T}
 e^{2i\theta_\chi(\gamma)}
 \frac{L(\sigma+i\gamma,\chi)}
      {\overline{L(\sigma+i\gamma,\chi)}}\longrightarrow0.
\end{equation}
Every sum over zeros in this theorem retains multiplicity. In particular,
with $\delta_\gamma=\arctan(1/(2\gamma))$,
\begin{equation}\label{eq:A037-second}
 \frac1{N_\chi(T)}\sum_{0<\gamma\le T}
 \cos^2\bigl(\theta_\chi(\gamma)+
              \phi_{\chi,3/2}(\gamma)-\delta_\gamma\bigr)
 \longrightarrow\frac12.
\end{equation}
Consequently,
\[
 \liminf_{T\to\infty}\frac1{N_\chi(T)}
 \sum_{\substack{0<\gamma\le T\\\rho\ {
m simple}}}
 r_{\rho,\chi}^{\,2}\ge\frac16.
\]
At least a lower proportion $1/6$ of all zero occurrences are simple
and transverse; relative to the simple zeros the lower proportion is
$1/4$. More quantitatively, if $0<\eta<1/2$ and $G_{\chi,\eta}(T)$
counts simple zeros with $r_{\rho,\chi}\ge\eta$, then
\begin{equation}\label{eq:A037-counts}
 \liminf\frac{G_{\chi,\eta}(T)}{N_\chi(T)}
 \ge\frac{1-4\eta^2}{6(1-\eta^2)},\qquad
 \liminf\frac{G_{\chi,\eta}(T)}{N_{\chi,s}(T)}
 \ge\frac{1-4\eta^2}{4(1-\eta^2)}.
\end{equation}
For example, at $\eta=1/4$ these bounds are $2/15$ and $1/5$.
For each fixed $\varepsilon>0$, every zero counted at a fixed positive
threshold also satisfies
$|\Re v_{\rho,\chi}|\gg_{\chi,\eta,\varepsilon}
\gamma^{1-\varepsilon}$ as $\gamma\to\infty$.
No positive/negative sign proportion is asserted.
\end{theorem}

\begin{lemma}[Twisted first Hardy mode]\label{lem:A037-twist}
For every fixed $x>0$, under GRH for the specified primitive character,
\begin{align}\label{eq:A037-twist}
 \sum_{0<\gamma\le T}e^{2i\theta_\chi(\gamma)}x^{i\gamma}
 &=-\frac{i^a\epsilon_\chi^{-1}}{\sqrt{qx}}
    \sum_{n\le qxT/(2\pi)}\chi(n)\Lambda(n)e^{-2\pi i n/(qx)}
    +O_{\chi,x}(\sqrt T\log^2 T)\\
 &=O_{\chi,x}(T).\nonumber
\end{align}
Changing the convention at the upper endpoint changes the displayed
expression by an admissible error. For $q=1$, $x=1$, this is the
$\kappa=1$ formula of Arias de Reyna, Corollary~3.2, equation~(13).
The proof below gives the required twist and character extension,
including their constants; they are not inferred from that untwisted
statement.
\end{lemma}

\begin{proof}
Put
\[
 \Delta_\chi(s)=\epsilon_\chi Q_\chi(1-s)/Q_\chi(s).
\]
At a critical-line point,
$\Delta_\chi(1/2+it)^{-1}=e^{2i\theta_\chi(t)}$.
The gamma quotient is meromorphic with all of its zeros and poles real;
it is holomorphic and nonzero in the open upper half-plane. In every
fixed vertical strip, Stirling's formula gives, uniformly for $t\ge2\pi$,
\begin{equation}\label{eq:A037-stirling}
 \Delta_\chi(\sigma+it)^{-1}
 =\epsilon_\chi^{-1}(qt/(2\pi))^{\sigma-1/2}
 e^{i\{t\log(qt/(2\pi))-t+\pi(2a-1)/4\}}
 (1+O_\chi(t^{-1})).
\end{equation}
This follows directly by applying the usual Stirling expansion to
$\Gamma((s+a)/2)$ and $\Gamma((1-s+a)/2)$. The phase constant is
recorded because it supplies $i^a$ in \eqref{eq:A037-twist}.

First choose $T'$ in $(T,T+1)$ at distance $\gg_\chi1/\log T$ from
all zero ordinates. Such a choice is possible because the number of
zeros in this interval, or in a fixed enlargement of it, is
$O_\chi(\log T)$: see Iwaniec--Kowalski, Proposition~5.7(1),
equation~(5.27), pp.102--103. Their Theorem~5.6,
equations~(5.23)--(5.24), p.102, gives the Hadamard product and
partial-fraction formula. That formula gives
\[
 \frac{L'}L(\sigma+iT',\chi)=O_\chi(\log^2 T)
 \quad(-3/2\le\sigma\le 1+1/\log T).
\]
For clarity, this standard bound does not require simple zeros:
subtract the partial-fraction expansion at $2+iT'$; the poles at
distance at most one contribute at most $O(\log T)$ terms of size
$O(\log T)$, counting multiplicity. The remaining terms are bounded
by a constant times $\sum_{k\ge1}\log(T+k+2)/k^2=O(\log T)$.
The gamma and pole factors contribute $O(\log T)$. This proves the
claimed bound and also explains the admissible contour choice.

Integrate
\[
 \frac1{2\pi i}\frac{L'(s,\chi)}{L(s,\chi)}
          \Delta_\chi(s)^{-1}x^{s-1/2}
\]
around the rectangle with left side $\Re s=-3/2$, right side
$c=1+1/\log T'$, bottom at a fixed
$t_0\in(2\pi,2\pi+1)$ avoiding zero ordinates, and top at $T'$.
Any finite set of omitted lower zeros contributes $O_{\chi,x}(1)$.
There are no gamma poles or trivial zeros in this rectangle. The zeta
pole when $q=1$ is below it. The residue at a zero is its multiplicity
times $\Delta_\chi(\rho)^{-1}x^{\rho-1/2}$; GRH identifies this with
the summand in \eqref{eq:A037-twist}.

The bottom edge is $O_{\chi,x}(1)$. On the top edge the preceding
logarithmic-derivative bound and \eqref{eq:A037-stirling} give
$O_{\chi,x}(\sqrt T\log^2T)$ (in fact one logarithm can be saved).
On the left edge the functional equation and absolute convergence at
$5/2-it$ give $L'/L=O_\chi(\log(t+2))$, while
$|\Delta_\chi|^{-1}=O_\chi(t^{-2})$. Its integral is bounded.

Extending the right-edge integral down from $t_0$ to $2\pi$ costs
$O_{\chi,x}(1)$: on this fixed compact segment the integrand is
uniformly bounded as $c\downarrow1$. On the right edge use the absolutely convergent series
$L'/L=-\sum\chi(n)\Lambda(n)n^{-s}$ and Stirling. Its error is
\[
 O_{\chi,x}\left(\sum_{n\ge2}\frac{\Lambda(n)}{n^c}
                 \int_{2\pi}^{T'}t^{c-3/2}\,dt\right)
 =O_{\chi,x}(\sqrt T\log T).
\]
Here $T^{c-1}=O(1)$ and $\sum\Lambda(n)n^{-c}=O(\log T)$.
For the remaining integral set $u=t/(2\pi)$, $X=T'/(2\pi)$,
and $b=qx$. Its value is
\begin{equation}\label{eq:A037-stationary-integral}
 -\epsilon_\chi^{-1}e^{i\pi(2a-1)/4}
 b^{c-1/2}\sum_{n\ge2}\frac{\chi(n)\Lambda(n)}{n^c}
 \int_1^X u^{c-1/2}
 e^{2\pi i u\log(bu/(en))}\,du.
\end{equation}

Here are the full summation bounds for stationary phase. Set $r=n/b$.
Split $[1,X]$ into its intersection with $[r/2,2r]$ and the two
remaining intervals. On either remaining interval the derivative of
$u\log(u/(er))$ has absolute value at least $\log2$. The first
derivative test bounds each such integral by $O(X^{c-1/2})$.
Summing gives $O_{\chi,x}(\sqrt T\log T)$. The finitely many $r<1$
can be handled by the same test away from a bounded interval and add
at most this error.

On $[A,B]=[1,X]\cap[r/2,2r]$, the stationary-phase lemma of Arias
de Reyna, Lemma~2.5 (the underlying primary lemmas are Levinson and
Gonek), states that the integral is
\[
 {\bf1}_{A\le r\le B}\,r^c e^{i\pi/4}e^{-2\pi i r}
 +O\left(r^{c-1/2}
 +\frac{r^{c+1/2}}{|A-r|+\sqrt r}
 +\frac{r^{c+1/2}}{|B-r|+\sqrt r}\right).
\]
Its constants are uniform for $c-1/2$ in the compact interval in use.
Multiplication by the coefficient in
\eqref{eq:A037-stationary-integral} turns the main term into
\[
 -i^a\epsilon_\chi^{-1}b^{-1/2}
       \chi(n)\Lambda(n)e^{-2\pi i n/b}.
\]
The first remainder sums to $O_{\chi,x}(\sqrt T)$, since
$\sum_{n\le C T}\Lambda(n)n^{-1/2}\ll\sqrt T$.
Except for finitely many $n$, $A=r/2$, so the lower-endpoint
remainder has the same bound. If $r\le X/2$, then $B=2r$ and the
upper-endpoint remainder again has that bound. For $X/2<r\le2X$
it is at most
\[
 C_{\chi,x}\sqrt T\log(T+2)
 \sum_{bX/2<n\le2bX}
       \frac1{|n-bX|+\sqrt T}
 \ll_{\chi,x}\sqrt T\log^2T.
\]
The last bound follows by splitting the distances into unit intervals,
or comparing with $\int_0^{C T}(u+\sqrt T)^{-1}\,du$; it is valid
even when $bX$ is an integer. The constant number of small $n$ has
already been absorbed. This proves the displayed asymptotic at $T'$.
Chebyshev's bound $\sum_{n\le y}\Lambda(n)\ll y$
(Iwaniec--Kowalski, Section~2.2, equation~(2.12), p.33) gives its
$O_{\chi,x}(T)$ consequence without requiring cancellation in the
additive prime sum.

Finally, under GRH each summand on the left has modulus one, and
there are $O_\chi(\log T)$ zeros between $T$ and $T'$. The change
in the prime sum has $O_{q,x}(1)$ terms, each $O(\log T)$.
Thus the formula holds at every real cutoff $T$, with the stated
endpoint convention and error.
\end{proof}

\begin{lemma}[The precise fixed-character simplicity input]
Under GRH for a fixed primitive Dirichlet $L$-function,
\[
 \liminf_{T\to\infty}\frac{N_{\chi,s}(T)}{N_\chi(T)}\ge\frac23.
\]
This is a classical restricted-pair-correlation consequence, used here
as background rather than counted as a new simplicity theorem.
\end{lemma}
\begin{proof}
We specify the imported theorem and the multiplicity calculation to
avoid confusing a family average with a result for a fixed character.
Rudnick--Sarnak, \emph{Zeros of principal $L$-functions and random
matrix theory}, Theorem~1.2, specializes at degree $m=1$, correlation
order two, to every fixed primitive Dirichlet character. Its hypothesis
(1.7) is automatic here (and the theorem explicitly permits $m\le3$).
Their zeros are indexed with multiplicity; equation~(1.4) sums over
distinct indices, not necessarily distinct ordinates. The test functions
$f(x_1,x_2)=g(x_1-x_2)$ have Fourier support in
$(u,-u)$ with $|u|<1$, exactly the permitted
$|\xi_1|+|\xi_2|<2$. With $K(u)=\sin(\pi u)/(\pi u)$, this gives
\[
 \frac1N\sum_{j\ne k}g(\widetilde\gamma_j-
                         \widetilde\gamma_k)
 \longrightarrow\int_{\mathbb R}g(u)(1-K(u)^2)\,du.
\]
Here the indices run through the first $N$ positive zeros, and their
normalization is that of the cited theorem. Equal ordinates have
exactly equal normalized ordinates, which is all that is needed below.

Take nonnegative Schwartz tests $g$ with $g(0)=1$ and Fourier support
strictly inside $(-1,1)$. Every zero of multiplicity $m\ge2$ contributes
$m(m-1)\ge m$ to the sum on the left. Hence the fraction of occurrences
belonging to nonsimple zeros is bounded above by the displayed integral
in the limit. The same statement holds at cutoffs $T$ that include each
whole multiplicity block.

For completeness the endpoint optimization uses only admitted
Schwartz tests. For $0<\beta<1$, approximate
$\beta^{-1}{\bf1}_{[-\beta/2,\beta/2]}$ in $L^2$ and $L^1$ by
smooth real even functions $b_j$ supported inside that interval with
$\int b_j=1$. Then $g_j=|\widehat b_j|^2$ are nonnegative Schwartz,
$g_j(0)=1$, have Fourier support in $[-\beta,\beta]$, and converge in
$L^1$ to $g_\beta(u)=(\sin(\pi\beta u)/(\pi\beta u))^2$.
Since $0\le1-K^2\le1$, their integrals converge to
\[
 \int g_\beta(1-K^2)
 =\frac1\beta-1+\frac\beta3.
\]
Indeed $\widehat g_\beta(u)=\beta^{-1}(1-|u|/\beta)_+$ and
$\widehat {K^2}(u)=(1-|u|)_+$, so their product integral is
$1-\beta/3$. First take $N\to\infty$, then the smooth approximation,
then $\beta\uparrow1$. The nonsimple fraction is at most $1/3$.
Central zeros and finitely many endpoint choices do not affect this
positive-height limit.
\end{proof}

\begin{proof}[Proof of NN-A037]
For $\sigma>1$ the absolutely convergent expansion is
\[
 \frac{L(\sigma+it,\chi)}{\overline{L(\sigma+it,\chi)}}
 =\sum_{n,m\ge1}
 \frac{\chi(n)\mu(m)\overline{\chi(m)}}{(nm)^\sigma}
       e^{it\log(m/n)}.
\]
Its absolute coefficient sum is at most $\zeta(\sigma)^2$.
For every fixed pair $m,n$, Lemma~\ref{lem:A037-twist}, with
$x=m/n$, says that the normalized mixed sum tends to zero because
$N_\chi(T)\sim T\log T/(2\pi)$. For the positive-height counting
convention this is the calculation in the proof of
Iwaniec--Kowalski, Theorem~5.8, equations~(5.34)--(5.36), pp.104--105;
Rudnick--Sarnak equation~(2.11) records the corresponding two-sided
count. All such normalized sums have
absolute value at most one under GRH. Dominated convergence therefore
proves \eqref{eq:A037-joint}, without any uniform estimate as $m/n$
varies. This proves one mixed Fourier mode; it does not assert
independence or uniform distribution of the full joint phase.

Multiplication by $e^{-2i\delta_\gamma}$ changes the mean by $o(1)$.
To see this without a height-uniform error claim, fix $Y$, discard the
finitely many zeros below $Y$, and bound the remaining error by
$\sup_{\gamma>Y}2|\delta_\gamma|$; then take $T\to\infty$ and
$Y\to\infty$. Taking real parts and using
$\cos^2 u=(1+\cos2u)/2$ proves \eqref{eq:A037-second}.

At a simple zero,
$L'(\rho,\chi)=-i e^{-i\theta_\chi(\gamma)}Z_\chi'(\gamma)$,
where $Z_\chi'(\gamma)$ is real and nonzero. Since
$\arg(i\rho)=\pi-\delta_\gamma$, the exact velocity formula gives
\[
 r_{\rho,\chi}
 =\big|\cos(\theta_\chi(\gamma)+
              \phi_{\chi,3/2}(\gamma)-\delta_\gamma)\big|.
\]
The choice of square root of the root number changes the phase by
$\pi$ and therefore changes neither this expression nor the theorem.

Put $p=N_{\chi,s}(T)/N_\chi(T)$ and
$g=G_{\chi,\eta}(T)/N_\chi(T)$. The preceding simplicity lemma
gives $p\ge2/3-o(1)$. The contribution of all nonsimple occurrences
to \eqref{eq:A037-second} is at most $1-p$, so the simple contribution
has lower limit at least $1/6$. Also
\[
 \frac12+o(1)
 \le(1-p)+\eta^2(p-g)+g
 =1-(1-\eta^2)p+(1-\eta^2)g.
\]
It follows that
$g\ge p-[2(1-\eta^2)]^{-1}+o(1)$, which gives the first inequality
in \eqref{eq:A037-counts}. Dividing by $p$, bounded below by
$2/3-o(1)$, gives the second. Letting $\eta\downarrow0$ proves the
two asserted nonzero-velocity proportions. The numerical fractions at
$\eta=1/4$ are exact rational substitutions.

Finally, the Euler product gives
$|L(3/2+i\gamma,\chi)|\ge1/\zeta(3/2)$. The exact classical input
for the fixed-character Lindelof bound is Iwaniec--Kowalski,
\emph{Analytic Number Theory}, Section~5.7, Corollary~5.20,
equation~(5.60), printed p.116, following Theorem~5.19 on that page.
The hypotheses are GRH and the Ramanujan--Petersson coefficient
bound; the latter holds trivially for a Dirichlet character, since
$|\chi(n)|\le1$. The analytic conductor for fixed $\chi$ is
$\asymp_\chi 1+|t|$. Cauchy's derivative formula on a small
fixed circle, using the functional equation on its left half, gives
$|L'(1/2+i\gamma,\chi)|\ll_{\chi,\varepsilon}\gamma^\varepsilon$.
Explicitly, choose the circle radius smaller than $\varepsilon/4$;
Lindelof on the right and the gamma ratio bound on the left make the
supremum $O(\gamma^\varepsilon)$ after reducing the auxiliary exponent.
Thus $|v_{\rho,\chi}|\gg\gamma^{1-\varepsilon}$, and multiplication
by the fixed margin $\eta$ proves the final assertion.
\end{proof}

\paragraph{Exact scope and the Dedekind obstruction.}
This theorem is about every fixed primitive degree-one function. It
strengthens NN-A036 for $\zeta$ and proves its requested Dirichlet
analogue. A general Dedekind zeta function is not primitive. Repeated
Artin factors can produce nonsimple zeros systematically, and the
degree-one first-mode calculation cannot be imported for its product
without a new estimate. Nor do \eqref{eq:A037-joint} and
\eqref{eq:A037-second} distinguish the two signs of $\Re v$: the
sign of the real derivative $Z_\chi'(\gamma)$ is lost upon squaring.
The index-parity correlation and the magnitude-tail obstruction
recorded in Q-A036 therefore remain genuine sign-density gates.

\paragraph{Algebraic-geometric and representation translation.}
A primitive character gives its actual rank-one Artin representation
on the finite etale cyclotomic cover of
$\operatorname{Spec}\mathbb Z[1/q]$. Its root number supplies the
half-phase that makes the completed critical-line germ Real, including
for complex characters. At a reduced analytic zero the specified
deformation has differential
\[
 dL_t=L'(\rho,\chi)\,ds-\rho L(\rho+1,\chi)\,dt,
\]
so its tangent section has the stated normalized real normal
component; the ambient divisor is counted with its actual local
lengths. SGA~1, Expose~I, section~3, Proposition~1 is the local
differentials dictionary. On the compact dual of $\mathbb Q_{>0}$,
the mixed-mode theorem annihilates the first Hardy character times
each fixed prime-circle character, and absolute Euler convergence
extends that annihilation to the displayed exterior ratio. This is
a precise Fourier statement, not a full product-Haar or Frobenius
equidistribution theorem for analytic zeros. Rudnick--Sarnak's
restricted correlation theorem controls the nonreduced occurrence
mass sufficiently to obtain the explicit reduced-locus density.
No shifted arithmetic family, pure motive, or complexity separation
is supplied by this dictionary.

\paragraph{Primary locators used.}
Arias de Reyna, \emph{On the distribution (mod~1) of the normalized
zeros of the Riemann zeta-function}, arXiv:1403.7126v1,
Lemma~2.5, Corollary~3.2 equation~(13), and Section~4.4 give the
stationary-phase input and the untwisted comparison.
Rudnick--Sarnak, \emph{Zeros of principal $L$-functions and random
matrix theory}, author manuscript PDF pp.1--5, especially the
multiplicity convention on p.1, definitions (1.2)--(1.4) on p.3,
and Theorems~1.1--1.2 on pp.4--5, give precisely the fixed-character
restricted correlation input. Their distinction between primitive
and product functions is on p.6. The manuscript is available at
\url{https://empslocal.ex.ac.uk/people/staff/mrwatkin/zeta/rudnick-sarnak.pdf}.
The completion and root number use DLMF~25.15.5--25.15.6; the
Stirling expansion is DLMF~5.11.3 with its remainder formulation
in Section~5.11(ii). The absolute-speed clause uses
Iwaniec--Kowalski, \emph{Analytic Number Theory} (AMS Colloquium
Publications~53, 2004), Section~5.7, Corollary~5.20, equation~(5.60),
p.116, with Theorem~5.19's GRH and Ramanujan hypotheses checked
as above. The same book supplies the auxiliary Chebyshev bound at
Section~2.2, equation~(2.12), p.33; Hadamard and the logarithmic
derivative at Theorem~5.6, equations~(5.23)--(5.24), p.102; the
local zero count at Proposition~5.7(1), equation~(5.27), pp.102--103;
and the positive-height count in Theorem~5.8's proof,
equations~(5.34)--(5.36), pp.104--105.
The printed text was checked in an online OCR transcription
of the primary book, rather than a page image; the publisher record
is \url{https://doi.org/10.1090/coll/053}.
The new mixed-mode and transversality deductions
are proved in full above; no historical novelty claim is made.


\section{The Newman boundary and exact positive heat matrices}

Use the normalization of the supplied workbench22,
\[
 H_t(z)=\int_0^\infty e^{tu^2}\Phi(u)\cos(zu)\,du,
 \qquad H_0(z)=\tfrac18\xi(1/2+iz/2),
 \qquad \partial_tH_t=-\partial_z^2H_t.
\]
The kernel $\Phi$ is positive and decays doubly exponentially. The
de Bruijn--Newman threshold property is imported as a known theorem:
all zeros of $H_t$ are real exactly when $t\ge\Lambda$. Rodgers--Tao,
arXiv:1801.05914v5, Theorem 1.1, proves $\Lambda\ge0$; A gives
$\Lambda\le0$. We use their intersection $\Lambda=0$, without claiming
that either theorem supplies simplicity at the endpoint.

\begin{theorem}[NN-A006: endpoint collision or escape to infinity]
Under A, every zero of $H_t$ is real and simple for $t>0$.
The zeros of $H_0$ are real, with their possible multiplicities retained.
If $t_j<0$, $t_j\uparrow0$, and $z_j$ is a nonreal zero of $H_{t_j}$,
then every finite accumulation point of $z_j$ is a multiple real zero
of $H_0$. In particular, if all zeta zeros are simple, such nonreal zeros
escape every compact subset of $\mathbb C$.

At a real multiplicity-$m$ zero $(t_0,z_0)$ the local branches are
\[
 z_j(t)=z_0+\sqrt{2(t-t_0)}\,h_{m,j}+O(|t-t_0|),
\]
where $h_{m,j}$ are the distinct real roots of the probabilists' Hermite
polynomial $\operatorname{He}_m$. For $m\ge2$, a real collision can occur only
at $t_0\le0$ under A. Replacing this by $t_0<0$ requires simplicity at
$t=0$, which A does not state.
\end{theorem}
\begin{proof}
The integral is jointly entire in $(t,z)$ locally, by the kernel decay,
so all compact-set convergences and derivatives used here are uniform.
Write $H_{t_0}(z)=a_m(z-z_0)^m+O((z-z_0)^{m+1})$.
The PDE gives
$\partial_t^\ell H=(-1)^\ell\partial_z^{2\ell}H$.
Set $t-t_0=u^2/2$ and $z-z_0=uy$; divide by $a_mu^m$.
The resulting analytic function equals at $u=0$
\[
 \sum_{\ell=0}^{\lfloor m/2\rfloor}
 \frac{m!(-1)^\ell y^{m-2\ell}}
 {2^\ell\ell!(m-2\ell)!}=\operatorname{He}_m(y).
\]
Gaussian differentiation and Rolle's theorem give $m$ real roots;
the Hermite differential equation excludes repeated roots. The implicit
function theorem at each root, with Rouche's theorem for completeness,
gives the stated expansion. For negative real $t-t_0$, at least two
of the nonzero Hermite roots give nonreal branches. If $t_0>0$, nearby
negative increments still have $t>0$, contradicting $\Lambda=0$.
This proves simplicity at all positive times.

If $z_j$ has a finite accumulation point $z_*$, uniform convergence
implies $H_0(z_*)=0$, and A makes $z_*$ real. If that zero were simple,
the implicit function theorem would give a unique local zero for all
small real $t$. Reality of the coefficients means its conjugate is
another local zero, so uniqueness forces the branch to be real. This
contradicts the nonreal $z_j$. Compactness gives the final escape claim.
\end{proof}

The theorem gives an exact dichotomy and no height law. In particular,
$\Lambda=0$ alone does not prove an exponential collision front or an
independent-pair collision approximation. In the completed arithmetic
coordinate $s=1/2+iz/2$, the double-zero law becomes
$s-s_0=\pm i\sqrt{(t-t_0)/2}+O(|t-t_0|)$; the normalization cannot
be omitted when comparing gap-based heuristics. The geometric local
model is a ramified double cover of the parameter disc (higher
Hermite branching for $m>2$), not a family with unramified individual
zero sections at the collision.

\begin{theorem}[NN-A009: strict heat positivity and determined moments]
Let the complete arithmetic expression $Z(a)$ be the one in
workbench22, equation \path{eq:rh-connes-arithmetic-heat}, including
both pole terms, the archimedean integral, and the prime-power sum.
Under A it has the exact representation
\[
 Z(a)=\sum_\rho e^{-a\gamma_\rho^2},\qquad a>0,
\]
with nontrivial zeros counted with multiplicity. Consequently:
\begin{enumerate}
\item $(-1)^kZ^{(k)}(a)>0$ for every $k\ge0$, and
$(\log Z)''(a)>0$.
\item If $a_1,\ldots,a_d>0$ are distinct, the matrix
$[Z(a_i+a_j)]_{i,j=1}^d$ is positive definite.
\item For every $d\ge1$, the matrix
$[(-1)^{i+j}Z^{(i+j)}(a)]_{i,j=0}^{d-1}$ is positive definite.
\item The positive finite measure
$\mu_a=\sum_\rho e^{-a\gamma_\rho^2}\delta_{\gamma_\rho^2}$
is uniquely determined among positive measures on $[0,\infty)$ by its
moments $(-1)^kZ^{(k)}(a)$.
\end{enumerate}
No simplicity assumption is needed.
\end{theorem}
\begin{proof}
The supplied unrestricted explicit formula is
$Z(a)=\sum_\rho e^{a(\rho-1/2)^2}$, absolutely and locally uniformly
convergent with all derivatives. A replaces $\rho-1/2$ by $i\gamma$.
There are infinitely many distinct positive values $\gamma^2$, since
an entire function has finitely many zeros in each compact set and
$\xi$ has infinitely many zeros. None is zero: $\zeta(1/2)\ne0$,
for example from the alternating eta series.
Termwise differentiation proves strict complete monotonicity.
After dividing the weights by $Z(a)$,
$Z''/Z-(Z'/Z)^2$ is the variance of $\gamma^2$, and is strictly
positive because its support has more than one point.

For the first matrix the quadratic form is
$\sum_\rho|\sum_i c_i e^{-a_i\gamma^2}|^2$.
A nonzero real linear combination of $d$ exponentials with distinct
real exponents has at most $d-1$ distinct real zeros: divide by one
exponential, differentiate, and use Rolle's theorem inductively on
any chosen finite interval containing the zeros. If a complex linear
combination vanished at infinitely many points, its real and imaginary
parts would each vanish identically by this fact; independence of the
exponentials then makes every coefficient zero. Thus the quadratic form
vanishes only at $c=0$. The second matrix has quadratic form
$\sum_\rho e^{-a\gamma^2}|\sum_{i=0}^{d-1}c_i\gamma^{2i}|^2$.
A nonzero polynomial has finitely many roots, proving strictness.

For determinacy, $\int e^{(a/2)x}\,d\mu_a=Z(a/2)<\infty$.
Any positive measure $\nu$ with the same moments has the same finite
exponential moment by monotone convergence of its nonnegative Taylor
series. Its moment generating function therefore agrees with that of
$\mu_a$ in $|z|<a/2$, by dominated termwise integration. Both Laplace
transforms are holomorphic on $\Re z<a/2$; the identity theorem gives
equality on that half-plane, hence equality of characteristic functions
on the imaginary axis. Fourier uniqueness for finite measures yields
$\nu=\mu_a$.
\end{proof}

The finite matrices are genuine Gram matrices of the displayed functions
in $L^2(\mu)$ or $L^2(\mu_a)$. An approximate matrix
$\widetilde G$ with a certified operator-norm error $\epsilon$ obeys
$\lambda_{\min}(\widetilde G)\ge\lambda_{\min}(G)-\epsilon$.
Under A any rigorously certified negative exact value is impossible;
an observed negative value must therefore fail either its error
enclosure, its admissibility, or the retained complete formula.
Strictness gives no uniform positive lower bound as the matrix size,
parameters, or conditioning vary. This is an analytic Hilbert-space
Gram realization. It is not a finite-dimensional pure Weil object,
and the tautological diagonal operator with spectrum $\gamma^2$
does not construct a new arithmetic Hilbert--Polya operator.


\section{A total-positivity boundary and two actual character constructions}

\begin{theorem}[NN-A007: the Newman threshold is a PF boundary]
Define the entire function
\[
 F_t(w)=\frac{H_t(i\sqrt w)}{H_t(0)}=\sum_{k\ge0}c_k(t)w^k,
 \qquad
 c_k(t)=\frac{\int_0^\infty e^{tu^2}\Phi(u)u^{2k}\,du}
 {(2k)!H_t(0)}.
\]
The square-root notation denotes the entire even-power series, so it
introduces no branch choice. Every $c_k(t)$ is strictly positive for
real $t$. Under A, the infinite Toeplitz matrix
\[
 T_t=(c_{j-i}(t))_{i,j\ge0},\qquad c_k=0\quad(k<0),
\]
is totally nonnegative if and only if $t\ge0$. Thus for each $t<0$
some finite minor is negative, although every entry on or above its
diagonal is positive. At each $t\ge0$,
\[
 F_t(w)=\prod_j(1+w/x_j(t)^2),
\]
where $x_j(t)>0$ lists the positive zeros of $H_t$, with multiplicities.
\end{theorem}
\begin{proof}
Positivity and convergence follow directly from the theta integral.
The order estimate $\log\max_{|z|\le R}|H_t(z)|=O_t(R\log R)$
in workbench22, \texttt{sec:rh-newman-collision}, implies that $F_t$
has order at most $1/2$. When all zeros of $H_t$ are real, pairing
$x_j$ with $-x_j$ gives the displayed genus-zero product; the possible
zero-free exponential factor is constant because the order of $F_t$
is less than one, and $F_t(0)=1$ fixes it. Moreover
$\sum_jx_j^{-2}<\infty$.

Each factor $1+bw$, $b\ge0$, has upper-bidiagonal Toeplitz matrix
$I+bS$, which is totally nonnegative: every nonzero minor has its
order-preserving matching and nonnegative product. Multiplication of
generating series is multiplication of their Toeplitz matrices.
Cauchy--Binet therefore gives nonnegative minors for every finite
product. Local uniform convergence gives coefficient convergence and
hence convergence of each finite determinant, proving the forward
direction without invoking a classification theorem.

For the converse use the known Edrei theorem, in the precise original
form of Proposition A on pp.86--87 and its proof in Section 3, p.92,
of Edrei, \emph{Proof of a conjecture of Schoenberg}, Canad. J. Math.
5 (1953), 86--94, DOI 10.4153/CJM-1953-010-3: an entire generating
function of a totally nonnegative sequence has all its zeros on the
negative real axis. A zero $z$ of $H_t$ gives a zero $w=-z^2$ of
$F_t$. If all such $w$ are negative real, every $z$ is real. Thus PF
is equivalent to all-real zeros. The threshold property and A together
with Rodgers--Tao give exactly $t\ge0$. Failure of total nonnegativity
means, by its definition, failure at a finite minor; coefficients are
real, so that minor is strictly negative.
\end{proof}

The distinction between positive coefficients and all minors is material:
entrywise positivity persists throughout the nonreal-zero phase. The
negative certificate has finite size for each fixed negative time, but
this argument supplies neither its size nor an algorithmic rate as
$t\uparrow0$. It therefore gives no complexity upper bound.

\begin{theorem}[NN-A008: zeta parameters in symmetric, unitary, and flag geometry]
Under A, let $\gamma_j>0$ list the positive zeta ordinates with
multiplicities and put
\[
 C=\sum_j\gamma_j^{-2},\qquad
 \beta_j=\gamma_j^{-2}/C,\qquad
 E(w)=\prod_j(1+\beta_jw)=\sum_{k\ge0}e_k(\beta)w^k.
\]
Then $C$ is finite and positive, every $0<\beta_j<1$, and
$\sum\beta_j=1$. The following constructions are well defined.
\begin{enumerate}
\item A normalized central positive-definite character of $S_\infty$ is
\[
 \chi_S(\sigma)=\prod_{k\ge2}
 \left((-1)^{k-1}\sum_j\beta_j^k\right)^{m_k(\sigma)}.
\]
It has Thoma parameters $(\alpha,\beta,\gamma)=(0,\beta,0)$.
\item A normalized continuous central positive-definite character on
$U(\infty)=\bigcup_nU(n)$ is
\[
 \chi_U(U)=\prod_j\det((1-\beta_j)I_n+\beta_jU),\qquad U\in U(n).
\]
These are two different characters, not restriction of a single formula.
On a permutation cycle of length $k$, the latter has factor
\[
 \prod_j\{(1-\beta_j)^k-(-\beta_j)^k\},
\]
whereas the former has factor $(-1)^{k-1}\sum_j\beta_j^k$.
\item Every finite upper-triangular Toeplitz truncation
$T^{(n+1)}=(e_{j-i}(\beta))_{0\le i,j\le n}$ has strictly positive
minors whenever upper triangularity does not force them to vanish.
For $1\le i\le n$, put
\[
 D_i=\det T^{(n+1)}_{\{0,\ldots,i-1\},\{n+1-i,\ldots,n\}},
 \quad D_0=D_{n+1}=1,\quad
 q_i=\frac{D_{i-1}D_{i+1}}{D_i^2}>0.
\]
These are literal positive coordinates in the finite Toeplitz
parametrization associated with quantum cohomology of the complete flag
variety $\mathrm{Fl}_{n+1}$, by the specified Rietsch theorem below.
\end{enumerate}
\end{theorem}
\begin{proof}
The paired Hadamard product for $\Xi(z)=\xi(1/2+iz)$ and the zero-count
bound give $\sum\gamma_j^{-2}<\infty$; its infinitely many positive
terms give $0<\beta_j<1$.
For (1), let $D$ be the positive trace-class diagonal operator with
eigenvalues $\beta_j$ on $\ell^2$, and let $U_\sigma$ permute the
factors of $(\ell^2)^{\otimes n}$. Direct summation in the tensor basis
gives
\[
 \operatorname{Tr}(D^{\otimes n}U_\sigma)
 =\prod_{k\ge1}(\operatorname{Tr}D^k)^{m_k(\sigma)}.
\]
Multiplication by $\operatorname{sgn}\sigma$ gives exactly the stated
formula, since $\operatorname{Tr}D=1$. It is compatible with adjoining
fixed points. For any finite complex coefficients $a_i$ and permutations
$\sigma_i$, its positive-definiteness is
\[
 \sum_{i,j}\overline{a_i}a_j\chi_S(\sigma_i^{-1}\sigma_j)
 =\operatorname{Tr}(D^{\otimes n}A^*A)\ge0,
 \quad A=\sum_j a_j\operatorname{sgn}(\sigma_j)U_{\sigma_j}.
\]
The tensor density commutes with all permutations, so the character is
central. The displayed parameters coincide with the standard Thoma
formula; its extremality is the known classification theorem, with exact
modern primary locator Okounkov, arXiv:math/9803037, Section 2.3,
Theorem (Thoma), printed p.23. No classification is needed for the
explicit existence proof just given.

For (2), at a fixed $0\le b\le1$ the determinant expands as
\[
 \det((1-b)I_n+bU)
 =\sum_{k=0}^n b^k(1-b)^{n-k}\operatorname{Tr}(\bigwedge^kU).
\]
It is a nonnegative combination of representation characters, normalized
at the identity because the dimensions sum to $(b+1-b)^n=1$.
Products remain positive definite. For fixed $n$, uniform convergence
follows from $|\det(I_n+\beta_j(U-I_n))-1|\le C_n\beta_j$ and
$\sum\beta_j<\infty$; the finitely many early factors cause no issue,
including any zeros. Hence the limit is continuous and positive
definite. The embedding $U\mapsto U\oplus1$ leaves every factor
unchanged. For a permutation $k$-cycle the eigenvalues are the $k$th
roots of unity, whose product gives $(1-b)^k-(-b)^k$.

For (3), use the finite products of the bidiagonal factors in the
preceding proof. A minor with increasing row indices $r_i$ and column
indices $c_i$ is forced zero unless $r_i\le c_i$ for all $i$.
When these inequalities hold, realize a positive term in the iterated
Cauchy--Binet formula by moving the rightmost selected index from $r_d$
to $c_d$, one unit at a time, then the next rightmost to its target,
and so on. At every step the selected indices are distinct and the
corresponding bidiagonal minor is either one or a positive $\beta_j$.
After finitely many steps all targets are reached. Further factors
have a diagonal term equal to one, so every longer product retains a
fixed strictly positive lower bound for this minor. Passing to the
limit proves strictness. The $D_i$ therefore are positive.
Rietsch, \emph{Totally Positive Toeplitz Matrices: Classical and Modern},
ICM2026, DOI 10.1137/25M1806727, Theorem 2.1 and equations (2.1)--(2.3),
identifies the positive Toeplitz space with the positive diagonal torus
using $d_i=D_i/D_{i-1}$ and $d_{n+1}=1/D_n$. Its successive-ratio
coordinates are exactly the displayed $q_i$. Section 2.1 identifies
these with the quantum parameters for $\mathrm{Fl}_{n+1}$.
\end{proof}

For a fully explicit finite representation measure, let
$\rho_n(\sigma)=\operatorname{sgn}(\sigma)U_\sigma$ and let
$P_\lambda$ be its central orthogonal projection onto the $S_n$
isotypic component of type $\lambda$. Then
\[
 p_n(\lambda)=\operatorname{Tr}(D^{\otimes n}P_\lambda)\ge0,
 \qquad\sum_{\lambda\vdash n}p_n(\lambda)=1,
\]
and $\chi_S|_{S_n}$ is the mixture of normalized irreducible
characters with these weights. Orthogonal-projection positivity and
$\sum P_\lambda=I$ prove both assertions. The standard Frobenius formula
identifies these weights as $f^\lambda s_{\lambda'}(\beta)$; the
factor $f^\lambda$ is essential when using normalized irreducible
characters.

This gives actual algebraic-group flag varieties and representation
objects with explicit real parameters. The parameters are infinite
analytic zero data; they are not nonnegative integer structure constants,
not an asserted canonical basis for GCT, and not Frobenius eigenvalues
of a finite Weil object. The finite quantum-cohomology dictionary is a
specified imported parametrization; it does not create an arithmetic
model of $\operatorname{Spec}\mathbb Z$ or prove a GCT positivity
hypothesis. In particular it is compatible with the supplied GCT
counterexamples.


\section{An explicit asymptotic for the actual symmetric velocity trace}

\begin{theorem}[NN-A013: every fixed polynomial residue moment]
Under A there is a sequence $T_n\in(n,n+1]$ avoiding all zero ordinates
such that, for every fixed integer $r\ge0$ and every $\epsilon>0$,
\[
 \mathcal V_r(T_n):=
 \sum_{0<\Im\rho<T_n}
 \operatorname{Res}_{s=\rho}
 \frac{s^r\zeta(s+1)}{\zeta(s)}
 =\frac{i^rT_n^{r+1}}{2\pi(r+1)}
 +O_{r,\epsilon}(T_n^{r+\epsilon}).
\]
The sum is over distinct nontrivial zeros, with the full residue at a
multiple zero. No simplicity assumption is needed. If all these zeros
are simple, then
\[
 \mathcal V_r(T_n)=\sum_{0<\gamma<T_n}\rho^{r-1}v_\rho,
 \qquad v_\rho=\frac{\rho\zeta(\rho+1)}{\zeta'(\rho)}.
\]
In particular,
\[
 \mathcal V_0(T_n)=\frac{T_n}{2\pi}+O_\epsilon(T_n^\epsilon),
 \qquad
 \mathcal V_1(T_n)=\frac{iT_n^2}{4\pi}
 +O_\epsilon(T_n^{1+\epsilon}).
\]
The latter is the derivative of the sum of zeros in the fixed contour
under Hurwitz deformation, including ramified clusters as symmetric
traces. It is not an asymptotic for individual velocities at multiple
zeros.
\end{theorem}
\begin{proof}
We explicitly import the following standard consequence of RH: there
exist $T_n\in(n,n+1]$ and an absolute $A_0$ such that
\[
 \sup_{1/2\le\sigma\le2}|\zeta(\sigma+iT_n)|^{-1}
 \ll\exp(A_0\log T_n/\log\log T_n).
\]
A primary paper recording the precise statement is Milinovich--Ng,
\emph{A note on a conjecture of Gonek}, Section 3, equation (3.2),
printed p.5; its original locator is Titchmarsh, Theorem 14.16.
This imported bound assumes RH, not simplicity. The functional equation
and Stirling's formula extend it to $-1/2\le\sigma\le2$ as
\[
 |\zeta(\sigma+iT_n)|^{-1}
 \ll T_n^{\min(\sigma-1/2,0)+\epsilon}
\]
for every fixed $\epsilon>0$. A also gives the usual Lindelof bound
$\zeta(\sigma+it)\ll_\epsilon t^\epsilon$ uniformly on fixed compact
real intervals with $\sigma\ge1/2$, away from bounded $t$.

Integrate $f_r(s)=s^r\zeta(s+1)/\zeta(s)$ counterclockwise around
the rectangle with vertical sides $\Re s=2$ and $\Re s=-1/2$,
bottom at a fixed positive nonzero-ordinate height $T_0$, and top at
$T_n$. Its only poles are the nontrivial zeta zeros inside; the pole
of $\zeta(s+1)$ is at $s=0$, outside the rectangle, and $1/\zeta(s)$
has a zero rather than a pole at $s=1$. Trivial zeros are outside.
Adding the fixed finitely many omitted lower residues contributes $O(1)$.

On the right edge,
\[
 \frac{\zeta(s+1)}{\zeta(s)}=
 \sum_{m\ge1}b(m)m^{-s},\qquad
 b(m)=m^{-1}\prod_{p\mid m}(1-p),\quad b(1)=1.
\]
The bound $|b(m)|\le1$ gives absolute convergence. The $m=1$ term
contributes
\[
 \frac1{2\pi}\int_{T_0}^{T_n}(2+it)^r\,dt
 =\frac{i^rT_n^{r+1}}{2\pi(r+1)}+O_r(1+T_n^r).
\]
For $m\ge2$, integration by parts yields
$|\int_{T_0}^{T_n}(2+it)^re^{-it\log m}dt|
\le C_r(1+T_n)^r/\log m$.
Summing with $|b(m)|m^{-2}$ gives an $O_r(1+T_n^r)$ error.

On the left edge, the functional equation and absolute convergence at
$3/2-it$ give $|1/\zeta(-1/2+it)|\ll t^{-1}$.
The numerator is $\zeta(1/2+it)\ll t^\epsilon$, so
$|f_r(-1/2+it)|\ll t^{r-1+\epsilon}$.
Its integral is $O_{r,\epsilon}(T_n^{r+\epsilon})$, including $r=0$.
On the top edge the numerator lies in $\Re(s+1)\ge1/2$, hence is
$O(T_n^\epsilon)$, and the imported good-height reciprocal estimate
gives $f_r=O(T_n^{r+2\epsilon})$ uniformly. The horizontal length
is fixed. The bottom edge is a constant contribution. Replace
$\epsilon$ by $\epsilon/2$ at the start and apply the residue theorem.
At a simple zero the residue equals
$\rho^r\zeta(\rho+1)/\zeta'(\rho)=\rho^{r-1}v_\rho$.
At a multiple zero the exact interpretation is NN-A005.
\end{proof}

These formulas strengthen the elementary contour identity by supplying a
verified asymptotic at explicit good-height sequences, and repair the
missing simplicity qualification in a velocity interpretation. They do
not extend automatically to every real cutoff $T$: an individual residue
near a cutoff need not be bounded by the displayed total error. Nor does
the small real part of $\mathcal V_1$ imply equal proportions of positive
and negative real velocities. A sum can cancel with arbitrarily unequal
sign counts; $N-1$ entries equal to $1$ and one equal to $-(N-1)$ already
show this. The odd phase-correlation gate of NN-A002 remains a separate
arithmetic problem.


\section{Finite detection must escape to larger Toeplitz windows}

\begin{theorem}[NN-A014: an explicit finite-window obstruction near the Newman boundary]
Use the normalized coefficients $c_k(t)$ of NN-A007. Let
$T_t^{[R]}=(c_{j-i}(t))_{0\le i,j\le R}$, and for $t<0$ define
\[
 R_*(t)=\min\{R\ge1:T_t^{[R]}\text{ has a negative minor}\}.
\]
Under A this minimum exists for every $t<0$, and
\[
 \lim_{t\uparrow0}R_*(t)=\infty.
\]
No hypothesis on simplicity of zeta zeros is used. The following
explicit quantitative margin is available for every fixed $R$.
Let $\delta_R>0$ be the minimum at $t=0$ of all minors of
$T_0^{[R]}$ not forced to vanish by upper triangularity. Put
\[
 C_R=\max(1,c_1(0),\ldots,c_R(0)),\qquad
 A_R=\max_{0\le k\le R}(2k+2)(2k+1)c_{k+1}(0),
\]
\[
 B_R=(R+1)!(R+1)A_RC_R^R,\qquad
 \epsilon_R=\frac{\delta_R}{2B_R}.
\]
Then all non-forced minors of $T_t^{[R]}$ are at least
$\delta_R/2$ for $-\epsilon_R\le t\le0$, so
$R_*(t)>R$ on $-\epsilon_R\le t<0$.
\end{theorem}
\begin{proof}
NN-A007 gives a negative finite minor at every $t<0$, which is
contained in some initial window. At $t=0$ the positive-zero product
has infinitely many positive parameters. The nonintersecting-index
construction in NN-A008 applies without the normalization $\sum\beta=1$
and proves positivity of every non-forced minor. There are finitely
many minors in a fixed window, so $\delta_R>0$.

The theta integral gives the exact coefficient equation
\[
 c_k'(t)=(2k+2)(2k+1)c_{k+1}(t)-2c_1(t)c_k(t).
\]
Indeed numerator differentiation inserts $u^2$, and denominator
differentiation inserts its normalized mean $2c_1(t)$.
Under the probability measure proportional to
$e^{tu^2}\Phi(u)du$, this derivative is
$\operatorname{Cov}(u^{2k},u^2)/(2k)!$.
For increasing real functions $f,g$, covariance is nonnegative:
with independent $X,Y$ it equals
$\frac12\mathbb E[(f(X)-f(Y))(g(X)-g(Y))]\ge0$.
Thus $0<c_k(t)\le c_k(0)$ for $t\le0$, and
$0\le c_k'(t)\le(2k+2)(2k+1)c_{k+1}(0)$.

An $r$ by $r$ determinant is a sum of $r!$ signed products; its
derivative is a sum of at most $r!r$ products with one differentiated
entry. For $r\le R+1$, each product has absolute value at most
$A_RC_R^{r-1}$, which is bounded by the displayed $B_R$ after
multiplication by $r!r$. Entries below the diagonal and their
derivatives are zero. The fundamental theorem of calculus gives
$|D(t)-D(0)|\le B_R|t|$ for every such minor and $t\le0$.
This proves the stated margin. Since $R$ was arbitrary, it proves
divergence of the minimal detecting window.
\end{proof}

\begin{corollary}[An effective lower bound for the detecting window]
Under A, for every integer $R\ge1$ put
\[
 D_R=\left\lfloor\frac{(R+1)^2}{4}\right\rfloor,
 \qquad
 \underline\epsilon_R=
 \frac{40000^{-D_R}}
 {2^{R+1}(D_R!)^2(R+1)!(R+1)}.
\]
Every minor of $T_t^{[R]}$ not forced to vanish by upper triangularity
is strictly positive for $-\underline\epsilon_R\le t\le0$.
In particular, for all sufficiently small negative $t$,
\[
 R_*(t)\gg
 \sqrt{\frac{\log(1/|t|)}{\log\log(1/|t|)}}.
\]
The implicit constant and the meaning of ``sufficiently small'' are
effective; the displayed factorial bound is an explicit version for
every finite $R$. No simplicity assumption is required.
The later NN-A022 improves the asymptotic lower bound to
$R_*(-a)\gg a^{-1/2}$, with an explicit constant, by controlling
all minors of each bounded order at arbitrary indices.
\end{corollary}
\begin{proof}
First we establish a deliberately generous effective bound for zero
ordinates. Hasanalizade--Shen--Wong, \emph{Counting zeros of the Riemann
zeta function}, arXiv:2107.06506, Corollary~1.2, equation~(1.5), gives
for $T\ge e$
\[
 \left|N(T)-\frac{T}{2\pi}\log\frac{T}{2\pi e}\right|
 \le .1038\log T+.2573\log\log T+9.3675.
\]
The lower bound on the right, after subtracting $T/100$, is positive
at $100$ and strictly increasing for $T\ge100$. For example,
$2\pi<7$, $e<3$, $\log100<5$, and $\log\log100<2$ give its value
at $100$ greater than
$100/7-.104\cdot5-.258\cdot2-9.368-1>0$; its derivative is greater
than $1/7-.104/100-.258/100-.01>0$.
Consequently $N(100j)>j$ and $\gamma_j\le100j$, with ordinates
indexed with multiplicity. Since $x_j=2\gamma_j$ in our heat
normalization, the product parameters satisfy
\[
 b_j=x_j^{-2}\ge\frac{1}{40000j^2}=\frac{\kappa}{j^2}.
\]

Consider an allowed $r$ by $r$ minor with increasing row and column
indices $a_i\le d_i$ in $\{0,\ldots,R\}$. The path construction in
NN-A008 moves the rightmost index first, one unit at a time, and
requires exactly $D=\sum_i(d_i-a_i)$ moves. Here
\[
 D\le r(R+1-r)\le D_R.
\]
Use one move in each of the first $D$ bidiagonal factors
$I+b_jS$, and use diagonal steps in every other factor. Its positive
Cauchy--Binet summand is $\prod_{j=1}^D b_j$; all other summands are
nonnegative. Passing from finite products to the coefficientwise
limit therefore gives
\[
 \delta_R\ge\frac{\kappa^{D_R}}{(D_R!)^2}.
\]
This includes $D=0$, whose chosen summand is $1$.

Uniform coefficient bounds remove the remaining dependence on
unknown finite minima. Define
\[
 C=\sum_{k\ge0}c_k(0)
   =\frac{H_0(i)}{H_0(0)}=\frac{1}{2\xi(1/2)},
 \qquad
 A_0=\sum_{k\ge0}(2k+2)(2k+1)c_{k+1}(0)
   =\frac{\xi''(0)}{4\xi(1/2)}.
\]
The positive theta integral proves convergence of both sums and
$C>1$, $A_0>0$. The accompanying Arb certificate, using
$\xi(s)=(s-1)\pi^{-s/2}\Gamma(1+s/2)\zeta(s)$, rigorously gives
\[
 1<C<2,\qquad 0<A_0<1.
\]
One may retain the exact constants instead of this certificate:
$C=1.00579\ldots$ and $A_0=0.0117395\ldots$.
Since $C_R\le C$ and $A_R\le A_0$, the preceding theorem applies
with derivative bound
$(R+1)!(R+1)2^R$. Substituting the lower bound for $\delta_R$
gives exactly $\underline\epsilon_R$.

For completeness, elementary factorial bounds give, for $R\ge1$,
\[
 \log(1/\underline\epsilon_R)
 \le 10(R+1)^2\log(R+2).
\]
Indeed use $D_R\le(R+1)^2/4$, $\log(n!)\le n\log n$,
$\log40000<11$, and $\log(R+2)>1$ in the explicit expression;
the sum of the resulting coefficients is less than $10$.
Writing $L=\log(1/|t|)$, choose $R+1$ of order
$\sqrt{L/(20\log L)}$. For sufficiently large $L$ the last display
is at most $L/2$, so $|t|<\underline\epsilon_R$ and $R_*(t)>R$.
This proves the asserted effective asymptotic lower bound.
\end{proof}

This is a precise sense in which backward heat leaves the positive
infinite object while every fixed finite truncation initially remains
positive. It does not assert that the minimum \emph{order} of a negative
minor diverges: a bounded-order minor could move to unbounded row or
column indices. The theorem controls the full window parameter and
supplies computable analytic margins, not a polynomial complexity
classification. In representation language, every fixed finite family
of these Toeplitz/Schur tests stays positive for a negative-time interval,
although the full specialization has already left the PF cone. In flag
geometry, every fixed-rank point remains in its positive Toeplitz chart
on that interval; the obstruction to a compatible positive object
occurs at ranks tending to infinity.

\begin{proposition}[A004 refinement: transversality is additional data]
The conditions of a real entire function with all zeros on
$\Re s=1/2$ and functional equation $L(1-s)=L(s)$ do not imply that
the formal shifted-series deformation has a nonzero first normal jet
at every simple zero. An exact example is
\[
 L(s)=((s-1/2)^2+1)(1+\tfrac23(s-1/2)^2),\qquad \rho=1/2+i.
\]
All four zeros are simple and on the line, but
\[
 \rho'(0)=\frac{25i}{4},\qquad
 \rho''(0)=-208+\frac{4979i}{16}.
\]
It therefore moves to the left at second order for small positive
time, despite a purely tangential nonzero first jet.
\end{proposition}
\begin{proof}
In $x=s-1/2$, the zeros are $\pm i,\pm i\sqrt{3/2}$, distinct,
and the polynomial is even with real coefficients.
The series defining $L_t$ converges locally for $|t|<1$, since
$L(s+k)=O(k^4)$ on compact $s$ sets. Substitution into NN-A004 gives
the displayed rational complex numbers; the included symbolic check
uses exact rational arithmetic. Since the first real coefficient is
$\Re(\rho''(0)/2)=-104$, the sign claim follows.
This model has no asserted Euler product. It does not refute
transversality for the actual zeta function; it isolates the extra
phase-nonvanishing input needed for that assertion.
\end{proof}


\section{Forward heat, character dynamics, and a rational flag limit}

\begin{theorem}[NN-A015: a universal normalized character limit]
Retain the exact heat normalization of NN-A007 and set, for $t\ge0$,
\[
 p_m(t)=\sum_{j\ge1}x_j(t)^{-2m},\qquad
 \beta_j(t)=\frac{x_j(t)^{-2}}{p_1(t)},\qquad
 E_t(w)=F_t(w/p_1(t)).
\]
Under A these are positive zero weights, indexed in decreasing order
with multiplicity, and $\sum_j\beta_j(t)=1$. They obey the exact
moment equations
\[
 p_m'=4m\sum_{j=1}^{m}p_jp_{m+1-j}-m(4m+2)p_{m+1}
 \qquad(m\ge1).                                      \tag{A015.1}
\]
The moments are analytic even at a possible multiple zero at $t=0$;
no separate differentiable labeling of colliding roots is needed.

As $t\to+\infty$,
\[
 p_1(t)\sim\frac{(\log t)^2}{32},\qquad
 E_t(w)\longrightarrow E_\infty(w)=\cosh\sqrt{2w}
                                                        \tag{A015.2}
\]
locally uniformly on $\mathbb C$. The square-root expression is its
entire even-power series. In addition,
\[
 \beta_j(t)\longrightarrow
 \beta_j^*=\frac{8}{\pi^2(2j-1)^2},\qquad
 \sum_j|\beta_j(t)-\beta_j^*|\longrightarrow0.           \tag{A015.3}
\]
Consequently the characters constructed in NN-A008 have the limits
\[
 \chi_{S,t}(\sigma)\longrightarrow
 \prod_{k\ge2}\left[
 (-1)^{k-1}\frac{2^k(4^k-1)\zeta(2k)}{\pi^{2k}}
 \right]^{m_k(\sigma)},                                \tag{A015.4}
\]
and, for $U\in U(d)$ with eigenvalues $\lambda_1,\ldots,\lambda_d$,
\[
 \chi_{U,t}(U)\longrightarrow
 \prod_{a=1}^d\cosh\sqrt{2(\lambda_a-1)}.                \tag{A015.5}
\]
The limit in (A015.2) and the asymptotic for $p_1=c_1$ follow from
the kernel without A. A identifies the whole path $t\ge0$ with
positive zero weights and the indicated characters.
\end{theorem}

\begin{proof}
For every fixed real $t$, Fourier decay gives $H_t(x)\to0$ as
$x\to\pm\infty$ on the real line. If its zero product were finite,
the order-less-than-one product for $F_t$ would make $H_t$ a nonzero
polynomial, contradicting this decay. Under A its real zero product
therefore has infinitely many positive parameters at each $t\ge0$.
The heat PDE, in the coordinate $w$, and normalization at $w=0$ give
\[
 \partial_tF_t=4wF_t''+2F_t'-2p_1F_t.
\]
Near $w=0$, put $q=F_t'/F_t$ and use the convergent product to write
\[
 \log F_t(w)=\sum_{m\ge1}\frac{(-1)^{m-1}}m p_m(t)w^m,
 \qquad q(w)=\sum_{m\ge0}(-1)^m p_{m+1}(t)w^m.
\]
Then $\partial_t\log F_t=4w(q'+q^2)+2q-2p_1$.
Comparison of each coefficient proves (A015.1). Every $p_m$ is a
fixed polynomial in $c_1,\ldots,c_m$ by this logarithmic expansion,
so it is analytic in $t$ even if zero branches meet. In particular,
\[
 p_1'=4p_1^2-6p_2>0.                                  \tag{A015.6}
\]
Strictness also follows directly from the nonzero variance below.

We give a direct concentration proof of (A015.2), including the
tail control needed for local uniform convergence. The kernel is
\[
 \Phi(u)=\sum_{n\ge1}
 (2\pi^2n^4e^{9u}-3\pi n^2e^{5u})e^{-\pi n^2e^{4u}}
 \qquad(u\ge0).
\]
There are constants $0<b<B<\infty$ such that, on this whole half-line,
\[
 b e^{9u-\pi e^{4u}}\le\Phi(u)
       \le B e^{9u-\pi e^{4u}}.                        \tag{A015.7}
\]
The first summand proves the lower bound, since
$2\pi^2-3\pi e^{-4u}\ge2\pi^2-3\pi>0$.
For the upper bound discard the negative term and bound the remaining
sum by $2\pi^2\sum_{n\ge1}n^4e^{-\pi(n^2-1)}$ after factoring
out $e^{9u-\pi e^{4u}}$.

Let $U_t\ge0$ have density $H_t(0)^{-1}e^{tu^2}\Phi(u)$, and put
$a_t=(\log t)/4$ for large $t$. Integrating on
$[a_t,a_t+1/t]$ in (A015.7) gives
\[
 H_t(0)\ge\exp(t a_t^2-Ct)                            \tag{A015.8}
\]
for an absolute $C$ and all sufficiently large $t$.
Fix $0<\epsilon<1$, an integer $p\ge0$, and $A_0\ge0$.
On $0\le u\le(1-\epsilon)a_t$, the extra weight
$(u/a_t)^p e^{A_0u/a_t}$ is bounded. The integral over this interval,
divided by (A015.8), is at most a constant times
\[
 a_t\exp\bigl(-[2\epsilon-\epsilon^2]t a_t^2
                         +Ct+9a_t+A_0\bigr)\longrightarrow0.
\]
On $u\ge b_t=(1+\epsilon)a_t$, for all sufficiently large $t$,
\[
 tu^2+9u+A_0u/a_t+p\log(u/a_t)\le\tfrac\pi2e^{4u}.
\]
Indeed $u^2e^{-4u}$ and $ue^{-4u}$ decrease for large $u$, and
$t b_t^2/e^{4b_t}=O((\log t)^2t^{-\epsilon})\to0$;
the remaining linear terms are smaller. The weighted tail integral
is therefore at most
\[
 B\int_{b_t}^\infty e^{-(\pi/2)e^{4u}}du
 \le \frac{B}{2\pi e^{4b_t}}
                 e^{-(\pi/2)e^{4b_t}},
\]
which still tends to zero after division by (A015.8).
These estimates prove $U_t/a_t\to1$ in every fixed $L^p$, and
uniform integrability of all fixed exponential weights in $U_t/a_t$.
Since $p_1=c_1=\mathbb E U_t^2/2$, the claimed asymptotic for $p_1$
follows. The variance identity
$p_1'=\operatorname{Var}(U_t^2)/2>0$ justifies (A015.6).

Set $V_t=U_t/\sqrt{p_1(t)}$. We have $V_t\to\sqrt2$ in every
fixed $L^p$ and the same exponential integrability. Moreover
\[
 E_t(w)=\mathbb E\cosh(\sqrt w V_t).
\]
For $|w|\le R$ the difference from $\cosh\sqrt{2w}$ is bounded by
\[
 \sqrt R\,\mathbb E\left[
 |V_t-\sqrt2|e^{\sqrt R(V_t+\sqrt2)}\right],
\]
which tends to zero by Cauchy--Schwarz, the $L^2$ convergence and the
exponential-moment bound. This proves local uniform convergence.

The positive zeros of
$H_t(y/\sqrt{p_1(t)})/H_t(0)=E_t(-y^2)$ converge, in their order,
to the simple positive zeros of $\cos(\sqrt2y)$, by Rouche's theorem
and uniform nonvanishing on the intervening compact real intervals.
Under A all zeros being counted are real. Thus
\[
 \sqrt{p_1(t)}x_j(t)\to\frac{\pi(2j-1)}{2\sqrt2},
\]
which gives pointwise convergence in (A015.3). Both the finite-time
and limiting sequences sum to one; the latter uses
$\sum_{j\ge1}(2j-1)^{-2}=\pi^2/8$. Pointwise convergence of
nonnegative unit-mass sequences then gives $\ell^1$ convergence:
first truncate at $J$, pass to the limit in that finite set, and
bound the two omitted tails by their remaining masses.

For every fixed $k$, $|a^k-b^k|\le k|a-b|$ on $[0,1]$, hence
the power sums converge. Summing the odd reciprocal powers gives
\[
 \sum_j(\beta_j^*)^k
 =\frac{8^k}{\pi^{2k}}(1-2^{-2k})\zeta(2k)
 =\frac{2^k(4^k-1)\zeta(2k)}{\pi^{2k}}.
\]
The finite-cycle product gives (A015.4). Finally
$\chi_{U,t}(U)=\prod_{a=1}^d E_t(\lambda_a-1)$, so (A015.5)
is the already proved local uniform limit. These are actual
positive-definite characters, either by the constructions in NN-A008
with $\beta^*$ or by the pointwise limit of their finite positive
Gram matrices.
\end{proof}

\begin{corollary}[Rational normalized flag coordinates and the classical limit]
Fix $n\ge1$ and form the $(n+1)$ by $(n+1)$ upper Toeplitz matrix
from the coefficients of $E_t(w)$. Let $\widetilde\Delta_i(t)$
be its top-right $i$ by $i$ minor, with boundary values
$\widetilde\Delta_0=\widetilde\Delta_{n+1}=1$, and let
\[
 \widetilde q_i(t)=
 \frac{\widetilde\Delta_{i-1}(t)\widetilde\Delta_{i+1}(t)}
 {\widetilde\Delta_i(t)^2}.
\]
Every limit $\widetilde q_i^*$ is a strictly positive rational number,
computed from the Toeplitz coefficients $2^k/(2k)!$. In particular,
\[
 \widetilde q_1^*=n(2n-1)(4n-3).
\]
For the unnormalized coefficients of $F_t$, the corresponding
Rietsch quantum flag coordinates satisfy
\[
 q_i(t)=p_1(t)^{-2}\widetilde q_i(t)
 \sim\frac{1024\widetilde q_i^*}{(\log t)^4}.
\]
Thus this specific positive path approaches the locus $q_1=\cdots=q_n=0$
of the quantum parameters, whose specialization is the classical
cohomology ring in the usual quantum-cohomology presentation.
\end{corollary}
\begin{proof}
Coefficient convergence gives minor convergence. The limiting product
has infinitely many strictly positive parameters $\beta_j^*$, so
NN-A008 makes every allowed limiting minor strictly positive. Its
coefficients are rational, proving the first assertion. For $n\ge2$,
put $a_k=2^k/(2k)!$. Then
\[
 \widetilde\Delta_1^*=a_n,\qquad
 \widetilde\Delta_2^*=a_{n-1}^2-a_{n-2}a_n,
\]
and substitution gives
\[
 \widetilde q_1^*=
 n^2(2n-1)^2-n(n-1)(2n-1)(2n-3)
 =n(2n-1)(4n-3).
\]
For $n=1$, both boundary minors are one and $a_1=1$, giving the same
formula. Each term of $\Delta_i$ has coefficient degree
$i(n+1-i)$ under the scaling $c_k\mapsto p_1^{-k}c_k$.
Taking the second difference of this degree gives exactly
$\widetilde q_i=p_1^2q_i$. The asymptotic follows from (A015.2).
The coordinate identification is the imported Rietsch Theorem~2.1,
equations~(2.1)--(2.3), and Section~2.1, already checked in NN-A008;
the $q=0$ specialization is that theorem's standard quantum flag
presentation, not a new claim of an arithmetic motive.
\end{proof}

The limiting symmetric-group cycle values are rational because
$\zeta(2k)/\pi^{2k}$ is rational. This supplies actual rational finite
flag points, while the limiting Thoma parameters themselves contain
$\pi^{-2}$. Neither fact gives an integral positive basis for a GCT
multiplicity problem. The moment evolution and the universal limit
belong to this specified analytic heat family; they are not a theorem
about arbitrary $L$-function deformations.


\section{The first heat correction and mass at high zero indices}
\label{sec:analytic-heat-correction}

\begin{theorem}[NN-A019: Gaussian saddle and first character correction]
Retain the notation of NN-A015. For sufficiently large $t$, let $b_t>1$
be the unique large solution of
\[
 4\pi e^{4b_t}=2tb_t+9,
 \qquad \lambda_t=8tb_t-2t+36.
\]
Let $U_t$ have density $e^{tu^2}\Phi(u)/H_t(0)$ on $[0,\infty)$,
and put
\[
 \mu_t=\mathbb E U_t^2=2p_1(t),\qquad
 v_t=\frac{\operatorname{Var}(U_t^2)}{\mu_t^2}.
\]
The following assertions about the kernel hold without A or M:
\begin{align}
 b_t&=\tfrac14\bigl(\log t+\log\log t-\log(8\pi)+o(1)\bigr),
 &p_1(t)&=\tfrac12b_t^2+o(1),\\
 \sqrt{\lambda_t}(U_t-b_t)&\Longrightarrow\mathcal N(0,1),
 &v_t&\sim\frac{1}{2tb_t^3},\\
 p_1'(t)&\sim\frac{b_t}{4t}\sim\frac{\log t}{16t}.
 \label{eq:heat-first-correction-scale}
\end{align}
The normal limit holds with every absolute moment and every fixed
exponential weight. Locally uniformly for $w\in\mathbb C$,
\begin{equation}\label{eq:heat-first-correction-entire}
 E_t(w)=E_*(w)+\frac{v_t}{4}w\bigl(E_*(w)-E_*'(w)\bigr)+o(v_t),
 \qquad E_*(w)=\cosh\sqrt{2w}.
\end{equation}
Equivalently, with $\epsilon_t=v_t/4$,
\[
 E_t(w)=e^{\epsilon_tw}E_*((1-\epsilon_t)w)+o(\epsilon_t).
\]

Under A, the positive normalized zero weights satisfy, for every
fixed $j$,
\begin{equation}\label{eq:heat-first-correction-weights}
 \beta_j(t)=\beta_j^*(1-v_t/4)+o(v_t),\qquad
 \beta_j^*=\frac{8}{\pi^2(2j-1)^2}.
\end{equation}
For every fixed $k\ge2$, their power sums obey
\[
 \sum_j\beta_j(t)^k
 =\left(1-\frac{k v_t}{4}\right)\sum_j(\beta_j^*)^k+o(v_t).
\]
If $\sigma$ is a fixed finite permutation moving $d(\sigma)$ points,
then its Thoma character has the sharp relative expansion
\begin{equation}\label{eq:heat-first-correction-character}
 \chi_{S,t}(\sigma)=\chi_{S,*}(\sigma)
           \left(1-\frac{d(\sigma)v_t}{4}+o(v_t)\right).
\end{equation}
Nevertheless, the vectors $(\beta(t)-\beta^*)/v_t$ do not converge
in $\ell^1$, and
\[
 \liminf_{t\to\infty}
       \frac{\|\beta(t)-\beta^*\|_1}{v_t}\ge\frac12.
\]
Thus the first-order finite-character limit cannot be obtained by
differentiating the full weight sequence in $\ell^1$.
\end{theorem}

\begin{proof}
Write
\[
 f_t(u)=tu^2+9u-\pi e^{4u},\qquad
 h(u)=\Phi(u)e^{-9u+\pi e^{4u}}.
\]
The exact kernel and its bounds in NN-A015 give constants
$0<c<C<\infty$ with $c\le h(u)\le C$ for all $u\ge0$ and
$h(u)\to2\pi^2$ as $u\to\infty$. The latter follows by factoring
the $n=1$ summand; its negative correction is
$-3\pi e^{-4u}$ and the other summands tend to zero uniformly by
the same summable upper bound used in (A015.7).

For large $t$, $f_t'(1)>0$ whereas $f_t'(u)\to-\infty$.
Since $(f_t')''(u)=-64\pi e^{4u}<0$, there is exactly one root
$b_t>1$, the local maximum. Evaluating $4\pi e^{4b}-2tb-9$ at
$b=(\log t)/4$ and $b=(\log t+2\log\log t)/4$ shows
$b_t\sim(\log t)/4$. Taking logarithms of its defining equation
then gives the more precise displayed expansion. Also
$f_t''(b_t)=-\lambda_t$ and $\lambda_t\sim8tb_t$.

Here is a full tail justification for the Gaussian saddle method.
Fix a small constant $\delta>0$. For all sufficiently large $t$,
\[
 f_t''(u)\le-c_0\lambda_t\qquad (u\ge b_t-\delta)
\]
with $c_0>0$ independent of $t$: insert
$16\pi e^{4b_t}=8tb_t+36$ and use $b_t\to\infty$.
Since $f_t'(b_t)=0$, this implies
\[
 f_t(u)-f_t(b_t)\le-\tfrac12c_0\lambda_t(u-b_t)^2
               \qquad(u\ge b_t-\delta).
\]
On $[0,b_t-\delta]$, the maximum of $f_t$ is at an endpoint.
Indeed $f_t'$ is strictly concave, negative at $0$, positive at $1$,
and zero at $b_t$; its earlier zero is a local minimum.
Moreover $f_t(b_t)\sim tb_t^2$ whereas $f_t(0)=-\pi$.
Consequently
\[
 \sup_{0\le u\le b_t-\delta}f_t(u)
                       \le f_t(b_t)-c_1\lambda_t
\]
for some $c_1>0$. The integral on this interval remains negligible
even after multiplication by any fixed power of
$|\sqrt{\lambda_t}(u-b_t)|$ and by
$\exp(A|\sqrt{\lambda_t}(u-b_t)|)$, for fixed $A$:
its logarithmic extra cost is at most
$A b_t\sqrt{\lambda_t}+O(\log(b_t\lambda_t))=o(\lambda_t)$.

For bounded $y$, Taylor's formula gives
\[
 f_t(b_t+y/\sqrt{\lambda_t})-f_t(b_t)
                  =-y^2/2+O(|y|^3/\sqrt{\lambda_t}),
\]
because $f_t'''=O(\lambda_t)$ in a fixed neighbourhood of $b_t$.
The bounded positive factor $h$ tends there to $2\pi^2$.
The preceding Gaussian majorant and left-tail estimate justify
dominated convergence after this change of variables, including
every fixed polynomial and exponential weight. Hence, writing
$Y_t=\sqrt{\lambda_t}(U_t-b_t)$, one has
$\mathbb E Y_t\to0$, $\mathbb E Y_t^2\to1$, and uniformly bounded
absolute moments of every fixed order with exponential weights.

It follows that
\[
 \mu_t=b_t^2+o(b_t/\sqrt{\lambda_t})+O(1/\lambda_t)
          =b_t^2+o(1),
 \qquad
 \operatorname{Var}(U_t^2)\sim\frac{4b_t^2}{\lambda_t}.
\]
For the variance, expand
$U_t^2=b_t^2+2b_tY_t/\sqrt{\lambda_t}+Y_t^2/\lambda_t$;
the cross term is $O(b_t/\lambda_t^{3/2})$ and the last variance
is $O(1/\lambda_t^2)$, both smaller than the leading term.
The claims about $v_t$ and $p_1$ follow. Differentiating the
positive exponential tilt gives
$p_1'(t)=\operatorname{Var}(U_t^2)/2$, proving its asymptotic.

Set $V_t=U_t^2/\mu_t$, so that $\mathbb E V_t=1$ and
$\mathbb E(V_t-1)^2=v_t$. For each compact set $|w|\le R$,
the entire function of $v$
$g_w(v)=\cosh\sqrt{2wv}$ has, on $v\ge0$, third derivative bounded
by $C_R\exp(C_R\sqrt v)$. This follows either from its power
series on $0\le v\le1$ or from its explicit derivatives on
$v\ge1$. Taylor expansion at $v=1$ therefore has an expected
remainder bounded by
\[
 C_R\mathbb E\bigl[|V_t-1|^3
                         e^{C_R\sqrt{\max(V_t,1)}}\bigr].
\]
To bound it, put $\eta_t=1/(b_t\sqrt{\lambda_t})$. The exact
formula for $U_t^2$ and $\mu_t\sim b_t^2$ show
\[
 |V_t-1|\le C\eta_t(|Y_t|+\mathbb E|Y_t|)
       +C\eta_t^2(Y_t^2+\mathbb E Y_t^2),
 \qquad
 \sqrt{V_t}\le C+C\eta_t|Y_t|.
\]
The weighted moment bounds imply that the remainder is
$O_R(\eta_t^3)=o(v_t)$, since $v_t\sim4\eta_t^2$.
The linear Taylor term has zero mean. With $a=\sqrt{2w}$,
\[
 g_w''(1)=\frac{a^2\cosh a-a\sinh a}{4}
          =\frac w2\bigl(E_*(w)-E_*'(w)\bigr).
\]
Using $E_t(w)=\mathbb E g_w(V_t)$ proves
\eqref{eq:heat-first-correction-entire}, uniformly on the compact set.

Under A, the simple limiting roots
$w_j^*=-1/\beta_j^*$ correspond to the ordered finite-time roots
$w_j(t)=-1/\beta_j(t)$, as proved in NN-A015. The local root
expansion from \eqref{eq:heat-first-correction-entire} is
\[
 w_j(t)-w_j^*
 =-\frac{v_t}{4}\,
       \frac{w_j^*(E_*(w_j^*)-E_*'(w_j^*))}{E_*'(w_j^*)}
       +o(v_t)
 =\frac{v_t}{4}w_j^*+o(v_t).
\]
This proves \eqref{eq:heat-first-correction-weights}. On a fixed
small disc avoiding the zeros of $E_*$,
\[
 \log E_t(w)=\log E_*(w)+\frac{v_t}{4}
                w\left(1-\frac{E_*'(w)}{E_*(w)}\right)+o(v_t).
\]
Cauchy's formula permits coefficient comparison. The coefficient
of $w$ remains one; for $k\ge2$, comparison with
$\log E_t=\sum_{k\ge1}(-1)^{k-1}(\sum_j\beta_j(t)^k)w^k/k$
gives the power-sum formula. Multiplying the finitely many nonzero
cycle values gives \eqref{eq:heat-first-correction-character}.

Finally, if $(\beta(t)-\beta^*)/v_t$ converged in $\ell^1$, every
coordinate of its limit would be $-\beta_j^*/4$ by the proved
fixed-coordinate formula, so its coordinate sum would be $-1/4$.
Continuity of summation on $\ell^1$ instead makes that sum zero,
because each finite-time vector has total mass one. This is a
contradiction. For any fixed $J$, the zero total mass of
$\beta(t)-\beta^*$ also gives
\[
 \|\beta(t)-\beta^*\|_1
 \ge2\left|\sum_{j\le J}(\beta_j(t)-\beta_j^*)\right|.
\]
Divide by $v_t$, take a lower limit, and then let $J\to\infty$ to
obtain the lower bound $1/2$.
\end{proof}

\paragraph{Exact representation-theoretic translation.}
For every sufficiently small $\epsilon>0$,
\[
 e^{\epsilon w}E_*((1-\epsilon)w)
 =e^{\epsilon w}\prod_j(1+(1-\epsilon)\beta_j^*w)
\]
is the generating function of the genuine Thoma parameters
$\alpha=0$, $\beta=(1-\epsilon)\beta^*$ and $\gamma=\epsilon$.
Their total mass is one. Formula
\eqref{eq:heat-first-correction-entire} proves that this path gives
the first-order values of every fixed finite character observable
of the actual heat family, although that family has $\gamma=0$
at every finite time. The missing first-order mass at each fixed
zero index is compensated at indices tending to infinity.
The proved failure of an $\ell^1$ derivative makes this distinction
precise. The corresponding genuine $U(\infty)$ character is
\[
 e^{\epsilon\operatorname{Tr}(U-I)}
 \prod_j\det\bigl((1-(1-\epsilon)\beta_j^*)I
                         +(1-\epsilon)\beta_j^*U\bigr).
\]
Its exponential factor is positive definite on $U(d)$ by the
Poisson expansion in normalized tensor powers of the defining
character, and is compatible under $U\mapsto U\oplus1$.
Consequently the entire-function expansion gives, by finite
products, the first-order unitary character values as well.
These are exact infinite-group and analytic statements, not an
integral positive basis or an arithmetic motive. The only imported
classification is the same primary Thoma theorem used in NN-A008;
the saddle and correction proofs above are complete.
For the exterior-power dictionary, the concrete Grothendieck
locator is SGA~6, Expos\'e~V, Example~2.2(b) and Lemma~2.2.1.
The latter gives the exterior-power filtration for a locally
split exact sequence; the representation-ring case is explicitly
included after the lemma. It does not turn these analytic weights
into integral vector-bundle multiplicities.


\section{Negative heat requires minors of increasing order}
\label{sec:analytic-pf-order-escape}

\begin{lemma}[Uniform angular confinement for bounded negative heat]
\label{lem:analytic-heat-angular}
Fix $M>0$ and $0<\varepsilon<1$. Uniformly for $0\le a\le M$ and
$z$ with $-\operatorname{Im}z\ge\varepsilon|z|$, as $|z|\to\infty$,
\begin{equation}\label{eq:analytic-angular-asymptotic}
 H_{-a}(z)=H_0(z)
   \exp\left(a\left(\frac{H_0'}{H_0}(z)\right)^2\right)
       \left(1+O_{M,\varepsilon}
              \left(\frac{\log^2|z|}{|z|}\right)\right).
\end{equation}
In particular, for sufficiently large $|z|$ this sector contains
no zero of any $H_{-a}$ with $a\in[0,M]$. By conjugation the same
holds in the reflected upper sector. This lemma is unconditional.
\end{lemma}

\begin{proof}
Set $r=|z|$. In the half-plane $\operatorname{Im}w<-1$,
$s=1/2+iw/2$ satisfies $\operatorname{Re}s>1$ and
\[
 H_0(w)=\frac18\xi(s)
       =\frac1{16}s(s-1)\pi^{-s/2}\Gamma(s/2)\zeta(s)
\]
is nonzero. Every factor is nonzero there; in particular there are
no gamma poles and no zeta zeros in this half-plane. It is simply
connected, so choose a holomorphic logarithm $G(w)=\log H_0(w)$.
Write $L=G'(z)$. Stirling's expansion and the absolutely convergent
Euler series in $\operatorname{Re}s>1$ give
\begin{equation}\label{eq:analytic-angular-log-bounds}
 |L|=O_\varepsilon(\log r),\qquad
 G''(w)=O_\varepsilon(r^{-1})\quad
              (|w-z|\le 2r^{1/3}).
\end{equation}
For clarity, the gamma contribution to $G'$ is
$(i/4)\psi(s/2)$ and that to $G''$ is
$-(1/16)\psi'(s/2)$; the polynomial terms are $O(1/r)$ and
$O(1/r^2)$ respectively, and the logarithmic zeta derivatives
are $O(2^{-c_\varepsilon r})$. These estimates are uniform on
the indicated disc for large $r$.

We also need a global bound along the contours used below.
For $p=2aL$ and every real $u$, all points on the straight segment
from $z$ to $z+p+u$ have imaginary part at most
$-\varepsilon r+O_M(\log r)<-\varepsilon r/2$.
Their associated $s$ therefore remain in
$\operatorname{Re}s\ge c_\varepsilon r$. Uniform Stirling estimates
in the right half-plane and the Euler series give
\[
 |G'(w)|\le C_{M,\varepsilon}\log(r+|u|)
\]
along that segment, after enlarging the constant and taking $r$
large. Consequently the holomorphic-logarithm remainder
\[
 R(u)=G(z+p+u)-G(z)-L(p+u)
\]
satisfies
\begin{align}
 |R(u)|&\le
 \frac{C_{M,\varepsilon}}r(|u|+\log r)^2,
                  &&|u|\le r^{1/3},\label{eq:analytic-local-R}\\
 |R(u)|&\le C_{M,\varepsilon}
       (|u|+\log r)\log(r+|u|),
                  &&u\in\mathbb R.\label{eq:analytic-global-R}
\end{align}
The first bound is Taylor's formula with
\eqref{eq:analytic-angular-log-bounds}; the second is integration
of the global derivative bound and $|L|=O(\log r)$.

For $a>0$, the defining Fourier integral gives the exact Gaussian
convolution
\begin{equation}\label{eq:analytic-negative-heat-convolution}
 H_{-a}(z)=\frac1{\sqrt{4\pi a}}
       \int_{\mathbb R}e^{-u^2/(4a)}H_0(z+u)\,du.
\end{equation}
Indeed, integrating $\cos((z+u)v)$ against the Gaussian inserts
$e^{-av^2}$; Fubini is justified by
$|\cos((z+u)v)|\le e^{|\operatorname{Im}z|v}$ for real $u$
and the rapid decay of the theta kernel.

Shift the contour in \eqref{eq:analytic-negative-heat-convolution}
from $\mathbb R$ to $\mathbb R+p$. The integrand is entire in
$u$. Throughout the finite vertical displacement the point $z+u$
still lies in the zero-free half-plane just specified, so the same
branch of $G$ is valid; no gamma pole or zero is crossed.
The Gaussian dominates the growth of $H_0(z+u)$ on the vertical
ends of a large rectangle. More explicitly, for fixed $z,p,a$,
Stirling's bound is at most
$\exp(C_{z,p}|u|\log(2+|u|))$, whereas the Gaussian contributes
$\exp(-(\operatorname{Re}u)^2/(4a)+O_{p,a}(1))$.
Thus those ends tend to zero. With the contour parametrized by
$u+p$, completing the square now gives the exact identity
\[
 \frac{H_{-a}(z)}{H_0(z)e^{aL^2}}
 =\frac1{\sqrt{4\pi a}}
     \int_{\mathbb R}e^{-u^2/(4a)}e^{R(u)}\,du.
\]

The local error is uniform down to $a=0$. On $|u|\le r^{1/3}$,
\eqref{eq:analytic-local-R} tends uniformly to zero and
$|e^{R(u)}-1|\le2|R(u)|$ for large $r$. Its Gaussian mean is
$O_{M,\varepsilon}((a+\log^2r)/r)$.
On $|u|>r^{1/3}$, the right side of
\eqref{eq:analytic-global-R} is at most $u^2/(8M)$ for sufficiently
large $r$, uniformly throughout this tail. Since $a\le M$,
\[
 e^{-u^2/(4a)}|e^{R(u)}|
       \le e^{-u^2/(8a)}.
\]
After normalization, the tail integral is
$O_M(e^{-c_Mr^{2/3}})$, uniformly for $0<a\le M$; the same
bound holds for the omitted Gaussian mass. This proves
\eqref{eq:analytic-angular-asymptotic}. At $a=0$ the assertion is
the identity $H_0=H_0$. Finally, both displayed leading factors
are nonzero, so the relative error proves zero-freeness for
sufficiently large $r$.
\end{proof}

\begin{lemma}[A quadratic factor and the minor-width dictionary]
\label{lem:analytic-quadratic-pfr}
Let $b>0$, $0\le\theta<\pi/(r+1)$, and
\[
 P(w)=(1+be^{i\theta}w)(1+be^{-i\theta}w).
\]
Its coefficient Toeplitz matrix has all minors of order at most
$r$ nonnegative. Products, and locally uniform limits normalized
to have constant coefficient one, retain this property.
\end{lemma}

\begin{proof}
Write $a=be^{i\theta}$, $\bar a=be^{-i\theta}$, so the coefficients
are the elementary symmetric functions in these two variables.
For a $k$ by $k$ minor with increasing row indices $i_1,\ldots,i_k$
and column indices $j_1,\ldots,j_k$, the minor is zero if some
$i_\ell>j_\ell$, by triangularity. Otherwise set
\[
 \lambda'_p=j_{k+1-p}+p-i_1-k,\qquad
 \mu'_p=i_{k+1-p}+p-i_1-k,\qquad 1\le p\le k.
\]
These are partitions with $\mu'\subseteq\lambda'$. Transposing
and reversing both index lists identifies the determinant with
the dual Jacobi--Trudi determinant $s_{\lambda/\mu}(a,\bar a)$.
In particular $\lambda_1\le k$. The bound concerns the width of
the skew diagram, not its possibly arbitrarily large degree.
The dual Jacobi--Trudi/tableau identity follows by the usual
sign-reversing cancellation of intersecting paths in the
determinant expansion; the remaining paths are exactly
semistandard tableaux, with their variable weights.

Here the tableau formula can be evaluated directly, without any
Littlewood--Richardson positivity assumption. If a column of the
skew diagram has more than two boxes, the value is zero. Otherwise
every two-box column must contain $1$ above $2$, contributing
$a\bar a=b^2$. In each row the boxes paired with a box below form
an initial segment, and those paired with a box above form a final
segment, by the partition inequalities. The unpaired boxes form
one intervening interval of length $d\le\lambda_1\le k$. Their
weakly increasing entries can be chosen independently, with
generating function
\[
 h_d(a,\bar a)
   =b^d\frac{\sin((d+1)\theta)}{\sin\theta}>0.
\]
At $\theta=0$ this is $(d+1)b^d$. The full skew Schur value is
therefore a power of $b^2$ times a product of these nonnegative
quantities. This proves the assertion when $k\le r$.

Coefficient convolution is multiplication of the upper Toeplitz
matrices. For each fixed minor the Cauchy--Binet sum is finite,
since all intermediate indices lie between its smallest row and
largest column index. Hence products preserve the inequalities.
Local uniform convergence implies coefficient convergence and
therefore convergence of each fixed finite determinant.
\end{proof}

\begin{theorem}[NN-A020: the detecting minor order escapes]
Under A, for every integer $r\ge1$ there is $\delta_r>0$ such that
the coefficients of
\[
 F_t(w)=H_t(i\sqrt w)/H_t(0)
\]
belong to $PF_r$ for $-\delta_r\le t\le0$. In fact every allowed
minor of order at most $r$, at arbitrary indices, is strictly
positive on this interval. Here \emph{allowed} means
$i_\ell\le j_\ell$ for its increasing index lists.

Let $m_*(t)$ be the minimum order of a negative finite Toeplitz
minor. Under A and Rodgers--Tao, $m_*(t)$ is finite for each $t<0$,
but
\[
                    m_*(t)\longrightarrow\infty
                         \qquad(t\uparrow0).
\]
No simplicity hypothesis on $H_0$ is needed. This concerns minor
order uniformly over all indices; it strengthens the qualitative
finite-window result of NN-A014. The later NN-A022 supplies the
rate $m_*(-a)\gg a^{-1/2}$ by a different uniform product estimate.
\end{theorem}

\begin{proof}
For a zero $z$ in the right half-plane write its angle as
$|\arg z|\in[0,\pi/2)$. There are no zeros on the imaginary axis,
because $H_t(iy)=\int_0^\infty e^{tu^2}\Phi(u)\cosh(yu)\,du>0$.
We first show
\begin{equation}\label{eq:analytic-zero-angles-shrink}
 \sup_{\substack{H_t(z)=0\\\operatorname{Re}z>0}}|\arg z|
                         \longrightarrow0\qquad(t\uparrow0).
\end{equation}
Fix an angle $\eta>0$. Lemma~\ref{lem:analytic-heat-angular}, with
$M=1$ and $\varepsilon=\sin\eta$, excludes all zeros with
$|z|>R_\eta$ and distance in angle at least $\eta$ from the real
axis, uniformly for $-1\le t\le0$. On the remaining compact disc,
$H_t\to H_0$ uniformly. Under A, its finitely many zeros are real
and nonzero. Choose small neighbourhoods of them lying in the
$\eta$ sectors about the positive or negative real axes, and a
zero-free compact complement. Uniform convergence excludes zeros
on that complement for sufficiently small $|t|$. This argument
allows arbitrary multiplicity in the neighbourhoods: it requires
their locations to be real, not their splitting to remain real.
It also excludes a neighbourhood of zero, since $H_0(0)>0$.
This proves \eqref{eq:analytic-zero-angles-shrink}.

Choose $|t|$ small enough that every right-half-plane zero has
$|\arg z|<\pi/(2(r+1))$. The genus-zero product for $F_t$, whose
order is less than one, is absolutely locally convergent and has
constant coefficient one. Its zeros are $w=-z^2$. Real zeros of
$H_t$ give factors $1+bw$ with $b>0$; each nonreal quadruple gives
a quadratic factor as in Lemma~\ref{lem:analytic-quadratic-pfr},
with parameter angle $|\theta|=2|\arg z|<\pi/(r+1)$.
The real linear factors are totally nonnegative. The lemma and
product convergence prove $PF_r$ at all indices.

There are infinitely many factors. Otherwise $F_t$, and hence
$H_t$, would be a polynomial by its order-less-than-one product,
contrary to Fourier decay $H_t(x)\to0$ along the real axis and
$H_t(0)>0$. Every factor has constant term one and strictly
positive linear coefficient. For any allowed minor of order
$k\le r$, move its row configuration to its column configuration
one unit at a time, rightmost particle first, using a different
factor for each move. Each selected single-factor minor is its
positive linear coefficient, while a stationary configuration
has minor one. The particles never collide. This is a positive
term in the finite Cauchy--Binet expansion once sufficiently many
factors are used. Additional factors preserve this lower bound
by their stationary term, and the limiting determinant is
strictly positive.

At $t=0$ the same assertion is NN-A008. Shrinking the chosen
neighbourhood if necessary gives a closed interval
$[-\delta_r,0]$. For each $t<0$, NN-A007, using A and
Rodgers--Tao, says that $F_t$ fails $PF_\infty$, so some finite
minor is negative. But for any fixed $r$ none of orders at most
$r$ is negative near zero. This is precisely the claimed
divergence of $m_*(t)$.
\end{proof}

\begin{corollary}[Effective finite certificates for each fixed order]
Assume A. There is a terminating algorithm which, given $r\ge1$,
returns a positive rational $\delta_r$ and finite analytic bounds
certifying the conclusion of NN-A020. The algorithm does not
need to decide simplicity or the multiplicities of any zero.
No polynomial-time bound is asserted.
\end{corollary}

\begin{proof}
Choose a positive rational $q$, for example $q=1/(4(r+1))$, with
$\arctan q<\pi/(2(r+1))$, and set
$\varepsilon=q/\sqrt{1+q^2}$. The proof of
Lemma~\ref{lem:analytic-heat-angular} is effective: the gamma
estimates have computable remainder bounds and the zeta bounds
come from an absolutely convergent Dirichlet series. Thus its
constants and a rational threshold $R$ for zero-freeness in
$|\operatorname{Im}z|\ge\varepsilon|z|$ can be computed for
$0\le a\le1$. There is no compactness-only constant in that lemma.
For example, the only estimate uniform over an unbounded real
tail is \eqref{eq:analytic-global-R}; after division by $u^2$ its
majorant is
\[
 C\frac{\log(R+|u|)}{|u|}
          \left(1+\frac{\log R}{|u|}\right),
\]
decreasing for $|u|\ge R^{1/3}$ when $R>e$.
The resulting finite elementary inequalities, and the local
$C\log^2R/R$ error, have computable thresholds. Increasing $R$
to a rational value preserves a threshold established by these
eventually decreasing majorants. Explicit digamma and trigamma
bounds follow either from the remainder form of Stirling or
from Euler summation of their absolutely convergent series.

Consider the compact semialgebraic set
\[
 K_{R,q}=\{z:|z|\le R,\ |\operatorname{Im}z|
                                \ge q|\operatorname{Re}z|\}.
\]
Under A, $H_0$ has no zero on this set; the only real point of the
set is zero, where $H_0(0)>0$. Its positive minimum can be
certified without listing its zeros. Cover the set by rational
grid boxes, evaluate $H_0$ at their centers by complex-ball
arithmetic, and subtract a box-radius times a certified derivative
bound on $|z|\le R+1$. Refine the grid and the evaluations until
all these lower bounds are positive. Continuity and the positive
compact minimum guarantee termination. Intersection of a rational
box with $K_{R,q}$ is a real-algebraic decision, or one may use
an effective finite overcover converging to this compact set.
The Fourier integral provides a computable derivative bound.
Denote the resulting rational lower bound by $m_R>0$.

Compute a rational upper bound $B_R$ for
\[
 \int_0^\infty u^2e^{Ru}\Phi(u)\,du.
\]
Such a bound is effective from the explicit theta sum and its
superexponential tail. Since $0\le1-e^{-au^2}\le au^2$,
\[
 \sup_{|z|\le R}|H_{-a}(z)-H_0(z)|\le aB_R.
\]
It follows that
$\delta_r=\min(1,m_R/(2B_R))$ has the required property:
there are no zeros on $K_{R,q}$ for $0\le a\le\delta_r$, and
the angular lemma excludes them on the unbounded part of the
same cones. Every zero therefore lies in
$|\operatorname{Im}z|<q|\operatorname{Re}z|$.
The proved factor criterion and positive paths certify all
orders at most $r$, at arbitrary indices, simultaneously.
\end{proof}

\paragraph{Quantitative interface.}
If
$\alpha(t)=\sup_{H_t(z)=0,\,\operatorname{Re}z>0}|\arg z|$
is bounded explicitly, the same proof certifies $PF_r$ whenever
$2(r+1)\alpha(t)<\pi$. Thus a quantitative angular-confinement
bound would give a quantitative minimum-order lower bound.
The fixed-sector estimate proved here does not control the rate
at which the finitely many bounded zeros approach the real line,
nor its interaction with $R_\eta$ as $\eta\downarrow0$.
It gives neither an exponential collision-height law nor a
lower bound for all algorithms.
The corollary gives an effective fixed-order certificate. The later
NN-A022 strengthens it to the explicit interval
$-2^{-36}/(r+1)^2\le t\le0$ and the corresponding
inverse-square-root order lower bound.

\paragraph{Classical locators and exact geometric scope.}
The sector implication is classical Schoenberg theory. A precise
published statement is Adm--Garloff--Tyaglov, \emph{Total
nonnegativity of finite Hurwitz matrices and root location of
polynomials}, JMAA~467 (2018), 148--170, Theorem~5.8,
equation~(5.10), printed p.~161, and the quadratic case
(5.11)--(5.12), p.~162; DOI~10.1016/j.jmaa.2018.06.065.
The original reference given there is Schoenberg,
\emph{A note on multiply positive sequences and the Descartes
rule of signs}, Rend.\ Circ.\ Mat.\ Palermo~4 (1955), 123--131.
The proof above supplies the needed two-variable assertion
directly; no unverified converse for products is used.
Stirling estimates may be read in DLMF~5.11.1--5.11.2 and
their differentiated form; the actual Gaussian-contour
comparison needed here is proved above.

The Toeplitz minors are the skew Schur/exterior-power tests with
diagram width bounded by their order. The theorem proves that
every fixed such width, even with unbounded polynomial degree or
matrix index, retains positivity for a negative-time interval.
The full infinite Schur cone has nevertheless been left at every
negative time. This is a concrete failure of bounded-order
positivity tests to detect the boundary; no finite-dimensional
pure motive, integral GCT basis, or complexity-class separation
is asserted.
The concrete exterior-power dictionary is SGA~6, Expos\'e~V,
Example~2.2(b) and Lemma~2.2.1, with its representation-ring
example. Those operations supply the symmetric-function
formalism; the actual analytic positivity and width bound
are proved here.


\section{Negative heat remains positive on every fixed unitary group}
\label{sec:analytic-finite-rank-characters}

Keep the normalization
\[
 p_1(t)=F_t'(0)>0,\qquad
 E_t(w)=F_t(w/p_1(t)),\qquad
 \Psi_t(z)=E_t(z-1)=\sum_{k\ge0}a_k(t)z^k.
\]
Thus $E_t(0)=E_t'(0)=1$ and $\Psi_t(1)=1$ for every real $t$.
For $U\in U(d)$ with eigenvalues $z_1,\ldots,z_d$, define
\[
             \chi_{t,d}(U)=\prod_{i=1}^d\Psi_t(z_i).
\]
Here a normalized character means a continuous central function of
positive type with value one at the identity. It need not be the
character of a finite-dimensional representation.

\begin{theorem}[NN-A021: positivity at fixed unitary rank]
Assume A. For every fixed $d\ge1$ there is $\eta_d>0$ such that
$\chi_{t,d}$ is a normalized character of $U(d)$ for
$-\eta_d\le t\le0$. Its expansion is
\begin{equation}\label{eq:analytic-fixed-unitary-expansion}
 \chi_{t,d}(U)=
 \sum_{\ell(\lambda)\le d}c_{\lambda,d}(t)s_\lambda(U),
 \qquad
 c_{\lambda,d}(t)=
       \det\bigl[a_{\lambda_i-i+j}(t)\bigr]_{i,j=1}^d>0,
\end{equation}
where $a_k=0$ for $k<0$ and partitions are padded with zero parts.
The series converges absolutely and uniformly on $U(d)$, and
\[
       \sum_{\ell(\lambda)\le d}
          c_{\lambda,d}(t)s_\lambda(1^d)=1.
\]
In particular every polynomial irreducible representation has
strictly positive probability weight, with no bound on its degree.

Under A and Rodgers--Tao, for every $t<0$ the compatible functions
$\chi_{t,d}$ fail positive type at some finite rank. If
\[
 d_*(t)=\min\{d\ge1:\chi_{t,d}\text{ is not of positive type}\},
\]
then $d_*(t)<\infty$ for $t<0$ and
\[
                         d_*(t)\longrightarrow\infty
                                  \qquad(t\uparrow0).
\]
For $t\ge0$, all ranks have positive type by NN-A008. No
simplicity assumption is needed at $t=0$.
\end{theorem}

\begin{proof}
We first establish positivity of the Taylor coefficients in the
required Toeplitz sense. The order of $E_t$ is less than one, and its
genus-zero product is
\[
 E_t(w)=\prod_j(1+\beta_j(t)w),\qquad
 \sum_j|\beta_j(t)|<\infty,\qquad
 \sum_j\beta_j(t)=1.
\]
The factors are real positive or occur in conjugate pairs when
$t$ is sufficiently close to zero, with angles tending uniformly
to zero by NN-A020. More explicitly, a zero $x$ of $H_t$ in the
right half-plane gives $\beta=1/(p_1(t)x^2)$, so that
$|\arg\beta|=2|\arg x|$ near zero.

Under A, all $\beta_j(0)$ are positive. There are infinitely many,
and their sum is one, hence $\beta_{\max}(0)<1$. Choose
$b$ with $\beta_{\max}(0)<b<1$. The entire function $E_0$ is
zero-free on $|w|\le1/b$. Local uniform convergence
$E_t\to E_0$ makes this disc zero-free for all sufficiently small
$|t|$. Consequently
\begin{equation}\label{eq:analytic-unitary-beta-bound}
                |\beta_j(t)|<b\qquad\text{for every }j.
\end{equation}
This argument does not require individually labelling any multiple
zero. Define
\[
       \widetilde\beta_j(t)=\frac{\beta_j(t)}{1-\beta_j(t)}.
\]
For $\beta=\rho e^{i\theta}$, $\rho\le b<1$, with $|\theta|<\pi/2$,
\[
 |\arg\widetilde\beta|
 \le |\theta|+
       \arctan\frac{\rho|\sin\theta|}{1-\rho\cos\theta}
 \le \frac{|\theta|}{1-b}.
\]
Thus the angles of all $\widetilde\beta_j$ tend uniformly to zero,
and their absolute values are summable.

The real factors $1-\beta_j$ are positive; each conjugate pair
has product $|1-\beta_j|^2>0$. Absolute convergence therefore gives
\begin{equation}\label{eq:analytic-shifted-PF-product}
 \Psi_t(z)=E_t(-1)\prod_j(1+\widetilde\beta_j(t)z),
                  \qquad E_t(-1)>0.
\end{equation}
For a fixed $d$, make $|t|$ small enough that every conjugate-pair
angle is less than $\pi/(d+1)$. The quadratic factor lemma in
NN-A020 proves $PF_d$ for the coefficients of the normalized
product in \eqref{eq:analytic-shifted-PF-product}. Positive
constant scaling preserves every minor sign. Moreover every
allowed minor of order at most $d$ is strictly positive: the
infinite product has infinitely many factors, all with positive
linear coefficient, and the one-step particle paths in the
Cauchy--Binet proof of NN-A020 apply unchanged. This establishes
strict positivity at arbitrary indices, not merely in a bounded
coefficient window. Shrinking the neighbourhood gives a closed
interval $[-\eta_d,0]$.

We give the character expansion explicitly. For any entire
function $\Psi(z)=\sum_{k\ge0}a_kz^k$, multiply
$\prod_i\Psi(x_i)$ by
$\Delta(x)=\det[x_i^{d-j}]_{i,j=1}^d$ and expand the rows:
\[
 \det\bigl[x_i^{d-j}\Psi(x_i)\bigr]_{i,j=1}^d.
\]
Cauchy--Binet, with the exponents chosen in strictly decreasing
order $m_i=\lambda_i+d-i$, gives the coefficient of the alternant
$\det[x_i^{\lambda_j+d-j}]$ as
$\det[a_{\lambda_i-i+j}]$. Dividing by $\Delta$ gives
\begin{equation}\label{eq:analytic-alternant-formal}
 \prod_{i=1}^d\Psi(x_i)
     =\sum_{\ell(\lambda)\le d}
        \det[a_{\lambda_i-i+j}]s_\lambda(x_1,\ldots,x_d)
\end{equation}
as a degree-wise formal identity. Each degree involves only
finitely many terms. The Schur polynomials are the irreducible
polynomial characters by the Weyl character formula.

For $\Psi=\Psi_t$, the coefficient determinant is the Toeplitz
minor with row indices $i-\lambda_i$ and column indices $j$,
after adding the same sufficiently large integer to both lists.
The row indices increase strictly, and
$i-\lambda_i\le i$, so this minor is allowed. It is therefore
strictly positive. Set every $x_i=u$, $0\le u<1$, in
\eqref{eq:analytic-alternant-formal}. Nonnegative coefficients
and degree-wise convergence give
\[
 \Psi_t(u)^d
      =\sum_{\ell(\lambda)\le d}
        c_{\lambda,d}(t)u^{|\lambda|}s_\lambda(1^d).
\]
Monotone convergence as $u\uparrow1$ proves that the positive
weights $c_{\lambda,d}(t)s_\lambda(1^d)$ sum to one.
For $U\in U(d)$,
$|s_\lambda(U)|\le s_\lambda(1^d)$, since this is the trace of
a unitary representation. The series thus converges absolutely
and uniformly. The formal identity inside the unit polydisc,
followed by radial convergence, identifies its value with
$\prod_i\Psi_t(z_i)$. A positive mixture of normalized unitary
characters is of positive type, proving the first assertion.

For completeness we prove the required all-ranks converse,
including a possible zero constant coefficient, rather than
postulating an infinite-group classification. Suppose an entire
function $\Psi$ with real Taylor coefficients and $\Psi(1)=1$
has $\prod_i\Psi(z_i)$ of positive type on every $U(d)$.
At rank one its Fourier coefficients are nonnegative; these
are exactly its Taylor coefficients $a_k$, with the negative
coefficients zero. Let $m$ be the first index with $a_m>0$ and
write $\Psi(z)=z^m B(z)$. Then $B(0)>0$, $B(1)=1$, and
every product $\prod_i B(z_i)$ is again of positive type:
it is the original function multiplied by the one-dimensional
unitary character $\det(U)^{-m}$.

The coefficient of the irreducible polynomial character
$s_\lambda$ in this central function is
\begin{equation}\label{eq:analytic-necessary-Schur-coefficient}
                      \det[b_{\lambda_i-i+j}]_{i,j=1}^d,
                       \qquad B(z)=\sum_{k\ge0}b_kz^k.
\end{equation}
One may obtain this identity directly from the Weyl integration
formula: multiply the two alternants, expand their determinants,
and use
$\int_{S^1}z^n\,dz/(2\pi iz)=\mathbf1_{n=0}$.
Since a continuous central positive-type function on a compact
group has nonnegative irreducible-character coefficients,
\eqref{eq:analytic-necessary-Schur-coefficient} is nonnegative.
This last fact also follows directly from a positive-type
quadratic form integrated against matrix coefficients, followed
by Schur orthogonality. Thus the specialization
\[
             \rho(h_k)=b_k/b_0,\qquad h_0=1,
\]
is nonnegative on every ordinary Schur function, because every
partition occurs at some finite rank.

Here ordinary Schur positivity implies skew Schur positivity.
We recall a finite-dimensional proof of the needed implication.
In two disjoint alphabets $X,Y$, splitting a semistandard
tableau according to its alphabet gives
\[
        s_\lambda(X,Y)=
             \sum_{\mu\subseteq\lambda}s_\mu(X)s_{\lambda/\mu}(Y).
\]
On the other hand, restrict the polynomial irreducible
representation with character $s_\lambda$ of $U(p+q)$ to
$U(p)\times U(q)$, where $p,q$ exceed the diagram size.
Complete reducibility gives a nonnegative integer sum
\[
        s_\lambda(X,Y)=
             \sum_{\mu,\nu}N_{\mu,\nu}^\lambda
                           s_\mu(X)s_\nu(Y),
                         \qquad N_{\mu,\nu}^\lambda\ge0.
\]
The summands are polynomial representations, so only partitions
occur. Linear independence of the $s_\mu(X)$ identifies
$s_{\lambda/\mu}(Y)=\sum_\nu N_{\mu,\nu}^\lambda s_\nu(Y)$.
This is an identity of symmetric functions by the choice of
large alphabets. Applying $\rho$ proves nonnegativity of every
skew Schur function. The ordinary Jacobi--Trudi formula
identifies every allowed Toeplitz minor of $(b_k/b_0)$ with
one of these values; forbidden minors are zero. Thus $(b_k)$
is $PF_\infty$.

Edrei's theorem now says that the entire generating function
$B$ has only nonpositive real zeros: its possible denominator
factors would have positive real poles and cannot cancel its
nonpositive numerator zeros. The factor $z^m$ adds only the
zero at zero. We have proved that an entire $\Psi$ of the stated
all-ranks positive type cannot have a nonreal zero.

Apply this conclusion to $\Psi_t$. For every $t<0$, A and
Rodgers--Tao give a nonreal zero $x$ of $H_t$. There is no
zero on the imaginary axis since $H_t(iy)>0$. Hence
$x^2$ is nonreal, and
\[
                        z=1-p_1(t)x^2
\]
is a nonreal zero of $\Psi_t$. All-ranks positive type is
therefore impossible, so $d_*(t)$ is finite. Compatibility
\[
        \chi_{t,d+1}(U\oplus1)=\chi_{t,d}(U)
\]
shows that positivity at rank $d+1$ implies positivity at
rank $d$. The first part of the theorem then proves
$d_*(t)\to\infty$ as $t\uparrow0$.
\end{proof}

\begin{corollary}[The exact finite-rank dictionary]
Let $\Psi(z)=\sum_{k\ge0}a_kz^k$ be entire with real coefficients,
$a_0>0$ and $\Psi(1)=1$. For a fixed $d\ge1$,
\[
 \bigl[\,U\longmapsto\det\Psi(U)
                 \text{ is of positive type on }U(d)\,\bigr]
 \quad\Longleftrightarrow\quad
                  (a_k)_{k\ge0}\in PF_d.
\]
Consequently, for $t<0$ sufficiently close to zero, $d_*(t)$ is
exactly the minimum order of a negative Toeplitz minor of the
coefficients of $\Psi_t$. These are the coefficients of the
translated normalized function $\Psi_t(z)=E_t(z-1)$, not the
untranslated coefficients of $F_t$.
\end{corollary}

\begin{proof}
The forward construction from $PF_d$ to a positive normalized
Schur expansion was proved above and did not require strict
inequalities. Conversely positive type gives the nonnegative
coefficient determinants
\eqref{eq:analytic-necessary-Schur-coefficient} for every
partition of length at most $d$. Normalize $a_0$ to one in the
symmetric-function specialization. In the preceding skew Schur
expansion, the coefficient $N_{\mu,\nu}^\lambda$ vanishes unless
$\mu\subseteq\lambda$, by the two-alphabet tableau identity.
Taking equally large alphabets and interchanging their two
unitary blocks proves
$N_{\mu,\nu}^\lambda=N_{\nu,\mu}^\lambda$; hence it also
vanishes unless $\nu\subseteq\lambda$.

A Toeplitz minor of order $k\le d$, with increasing indices
$i_1,\ldots,i_k$ and $j_1,\ldots,j_k$, is zero when some
$i_\ell>j_\ell$. Otherwise transpose and reverse the index
lists, and put
\[
 \lambda_p=j_{k+1-p}+p-i_1-k,\qquad
 \mu_p=i_{k+1-p}+p-i_1-k,\qquad 1\le p\le k.
\]
These are partitions with $\mu\subseteq\lambda$ and
$\ell(\lambda)\le k$. The normalized determinant is
$\rho(s_{\lambda/\mu})$ by the ordinary Jacobi--Trudi identity.
Its expansion uses only shapes $\nu\subseteq\lambda$, so every
$\ell(\nu)\le k\le d$. All their specialized Schur values are
nonnegative by the character coefficients already obtained.
Thus every minor of order at most $d$ is nonnegative.
The minimum-order assertion follows from the equivalence and
$E_t(-1)>0$ near zero.
\end{proof}

\paragraph{Effectivity at fixed rank.}
The interval $\eta_d$ can be certified effectively under A.
First certify that $E_0$ is zero-free on a rational-radius disc
$|w|\le1/b$ for some rational $b<1$. Such a choice exists by
$\beta_{\max}(0)<1$, and rational grid evaluation with a
derivative bound certifies a positive compact minimum.
Uniform effective bounds for $\partial_t E_t$ on a slightly
larger compact set, obtained from the defining theta integrals
and the positive denominator $p_1(t)$, then give a rational
time interval on which \eqref{eq:analytic-unitary-beta-bound}
holds. In the effective corollary to NN-A020 choose its integer
parameter $r$ large enough that
\[
       \frac{2\arctan(1/[4(r+1)])}{1-b}
                                      <\frac{\pi}{d+1}.
\]
That corollary certifies exactly the angular confinement used
above. The smaller of the two time intervals is a certified
$\eta_d$. This asserts a terminating finite computation at
each fixed rank; it does not assert an asymptotic rate or
polynomial running time.

\paragraph{Primary character locators and geometric translation.}
The Weyl character formula used here is, for example, Etingof
et al., \emph{Introduction to Representation Theory},
Theorem~5.22.1, printed p.~132; the neighbouring
Theorem~5.23.2 gives complete reducibility of algebraic
representations. For the precise compact-group convention,
Faraut's \emph{Infinite Dimensional Spherical Analysis},
notes written by Sho Matsumoto, Kyushu (September 2007), section~4.1,
printed p.~56, states the positive normalized-character
expansion, and section~4.2, printed p.~57, gives the Weyl
integration and character formulas. The new conclusion is
the application to bounded negative heat time at every fixed
rank, not those classical character formulas.

Exterior powers and their additive filtration formula have
the concrete Grothendieck locator SGA~6, Expos\'e~V,
Example~2.2(b) and Lemma~2.2.1, including the representation-ring
case. The coefficients in
\eqref{eq:analytic-fixed-unitary-expansion} give positive
probability measures on actual irreducible polynomial
$U(d)$-representations. They are real weights, not asserted
integer multiplicities in an algebraic vector bundle.
The analytic heat path lies in each fixed-rank positive
character cone on some negative interval, while it lies
outside the compatible positive cone for all ranks at every
negative time. This is an unbounded-degree statement at
fixed rank, stronger than finitely many character tests;
it neither supplies an integral positive basis nor a GCT
obstruction or a lower bound against all algorithms.


\section{A square-root angular law and quantitative rank escape}
\label{sec:analytic-angular-heat-rate}

The following argument uses A to retain the full real zero product
of $H_0$, with multiplicities. It strengthens NN-A020's fixed-sector
argument quantitatively. Write $x_j>0$ for the positive zeros of
$H_0$, including multiplicity, and
\[
 H_0(z)=H_0(0)\prod_j(1-z^2/x_j^2),\qquad
 N(X)=\#\{j:x_j\le X\}.
\]
The product has no residual exponential factor: as in NN-A007,
$H_0(i\sqrt w)/H_0(0)$ has order at most $1/2<1$, so its genus-zero
Hadamard factorization has constant zero-free factor, fixed at one.
In particular $\sum_jx_j^{-2}<\infty$. The classical zero-count
bound supplies an effectively computable constant $C_0$ with
\[
                       N(X)\le C_0 X\log(X+3)\qquad(X>0).
\]
For the explicit primary bound one can use
Hasanalizade--Shen--Wong, Corollary~1.2, equation~(1.5),
printed p.~2, with $x_j=2\gamma_j$, as already checked in
NN-A014. A compact initial interval is covered by an effective
bound and a certified zero-free neighbourhood of zero.

\begin{lemma}[A global canonical-product remainder]
\label{lem:analytic-product-remainder}
For a nonreal $z$, set
\[
 L(z)=\frac{H_0'}{H_0}(z),\qquad
 S(z)=\sum_j\left(\frac1{|z-x_j|^2}+\frac1{|z+x_j|^2}\right)>0.
\]
The entire function
\[
 R_z(h)=\frac{H_0(z+h)}{H_0(z)e^{L(z)h}}
\]
satisfies, for every $h\in\mathbb C$,
\begin{align}
 |R_z(h)|&\le \exp\bigl(S(z)|h|^2/2\bigr),
                  \label{eq:analytic-R-global}\\
 |R_z(h)-1|&\le
     8S(z)|h|^2\exp\bigl(2S(z)|h|^2\bigr).
                  \label{eq:analytic-R-difference}
\end{align}
\end{lemma}

\begin{proof}
The paired product and logarithmic derivative give the locally
uniformly convergent canonical product
\[
 R_z(h)=\prod_{\xi\in\{x_j,-x_j\}}
           \left(1+\frac h{z-\xi}\right)
                       \exp\left(-\frac h{z-\xi}\right).
\]
The sum of $|z-\xi|^{-2}$ is finite; pairing makes the relation
with $L(z)$ explicit. No unrecorded exponential term remains.
For every complex $v$,
\[
 |1+v|^2\le\exp(2\operatorname{Re}v+|v|^2),
 \qquad |(1+v)e^{-v}|\le e^{|v|^2/2}.
\]
Multiply this inequality over finite products and take their
locally uniform limit to obtain \eqref{eq:analytic-R-global}.

Also $R_z(0)=1$ and $R_z'(0)=0$. Put $S=S(z)$ and
$Q(v)=R_z(v/\sqrt S)$. On $|v|=1$, $|Q(v)|\le e^{1/2}$.
Cauchy's coefficient estimate and the vanishing linear
coefficient give, for $|v|\le1/2$,
\[
 |Q(v)-1|
       \le e^{1/2}\sum_{k\ge2}|v|^k
       \le2e^{1/2}|v|^2<4|v|^2.
\]
For $|v|\ge1/2$, the first bound instead gives
$|Q(v)-1|\le2e^{|v|^2/2}\le8|v|^2e^{2|v|^2}$.
The same latter bound covers $|v|\le1/2$, proving
\eqref{eq:analytic-R-difference}.
\end{proof}

\begin{lemma}[A numerical zero-free test for the Gaussian convolution]
\label{lem:analytic-gaussian-nonzero}
For $a>0$ and nonreal $z$, put
\[
                A_z=aS(z),\qquad B_z=a^2|L(z)|^2S(z).
\]
If $A_z\le2^{-10}$ and $B_z\le2^{-10}$, then $H_{-a}(z)\ne0$.
More precisely,
\begin{equation}\label{eq:analytic-relative-Gaussian-bound}
 \left|\frac{H_{-a}(z)}{H_0(z)e^{aL(z)^2}}-1\right|
 \le16e^{16B_z}
 \left\{\frac{4B_z}{\sqrt{1-16A_z}}+
             \frac{2A_z}{(1-16A_z)^{3/2}}\right\}
 \le\frac7{16}<1.
\end{equation}
\end{lemma}

\begin{proof}
The theta integral gives the exact convolution
\[
 H_{-a}(z)=\frac1{\sqrt{4\pi a}}
             \int_{\mathbb R}e^{-u^2/(4a)}H_0(z+u)\,du.
\]
Let $p=2aL(z)$ and translate the integration line to
$\mathbb R+p$. The integrand is entire. On the vertical ends,
the Gaussian quadratic decay dominates the order-one bound
$|H_0(w)|\le\exp(C|w|\log(|w|+3))$ for fixed $z,p,a$.
Thus the end integrals tend to zero. A real part of $p$ only
reparametrizes the translated horizontal line. There is no
logarithm to continue and no obstruction from zeros crossed
in this deformation: $R_z$ itself is entire.
Completing the square gives
\[
 \frac{H_{-a}(z)}{H_0(z)e^{aL(z)^2}}
          =\mathbb E R_z(p+U),\qquad U\sim N(0,2a).
\]
With $S=S(z)$, \eqref{eq:analytic-R-difference} and
$|p+U|^2\le2|p|^2+2U^2$ show that the error is at most
\[
 16S e^{4S|p|^2}
       \mathbb E\bigl[(|p|^2+U^2)e^{4SU^2}\bigr].
\]
The two elementary Gaussian integrals are
\[
 \mathbb E e^{4SU^2}=(1-16aS)^{-1/2},\qquad
 \mathbb E U^2e^{4SU^2}
                       =2a(1-16aS)^{-3/2}.
\]
Since $S|p|^2=4B_z$, these give the first inequality in
\eqref{eq:analytic-relative-Gaussian-bound}.
For $A_z,B_z\le2^{-10}$, bound $e^{16B_z}<2$,
$(1-16A_z)^{-1/2}<2$, and
$(1-16A_z)^{-3/2}<3$. The right side is then at most
$256B_z+192A_z\le448/1024=7/16$.
The denominator is nonzero under A, and this strict relative
bound proves the assertion.
\end{proof}

\begin{lemma}[Uniform cone bounds with a vanishing radial tail]
\label{lem:analytic-cone-logderivative}
There is an effectively computable $C\ge1$ such that, whenever
$0<q\le1$, $r=|z|$ and $|\operatorname{Im}z|\ge qr$,
\begin{align}
 S(z)&\le Cq^{-2}\frac{\log(r+3)}{r+1},
                 \label{eq:analytic-S-cone}\\
 |L(z)|&\le Cq^{-1}\log(r+3).
                 \label{eq:analytic-L-cone}
\end{align}
Consequently
\[
 |L(z)|^2S(z)
       \le C^3q^{-4}\frac{\log^3(r+3)}{r+1}.
\]
\end{lemma}

\begin{proof}
By evenness and conjugation we may assume
$\operatorname{Re}z\ge0$. Choose a fixed positive $r_0$ with
$2r_0<x_1$. For $r\le r_0$, the product gives bounded $S(z)$
and $L(z)=O(r)$, which are absorbed by the claimed constants.
For $r\ge r_0$ split the positive zeros into
$x_j<r/2$, $r/2\le x_j\le2r$, and $x_j>2r$.

In the first range, $|z\pm x_j|\ge r/2$. The contribution to
$S$ is at most $8N(r/2)/r^2\ll\log(r+3)/r$.
Using the paired logarithmic derivative
\[
                  L(z)=-2z\sum_j\frac1{x_j^2-z^2},
\]
the same range contributes $O(N(r/2)/r)=O(\log(r+3))$.
In the middle range, $|z-x_j|\ge qr$ and, since
$\operatorname{Re}z\ge0$, $|z+x_j|\ge r$.
Its contributions are at most
$2q^{-2}N(2r)/r^2$ to $S$ and
$2q^{-1}N(2r)/r$ to $|L|$.
Finally partial summation of the count bound gives
\[
             \sum_{x_j>2r}x_j^{-2}
                          \ll\frac{\log(r+3)}r.
\]
This last range contributes at most eight times that sum to
$S$, and $O(r)$ times that sum to $|L|$.
Combine the bounds and compare $r$ with $r+1$ for $r\ge r_0$.

All constants can be made effective from the stated zero-count
bound, the computable sum $\sum_jx_j^{-2}=F_0'(0)$, and a
certified zero-free disc at zero. To obtain the last disc one
can bound $|H_0(z)-H_0(0)|$ by $|z|$ times the computable theta
integral bound for $H_0'$ on $|z|\le1$, and use $H_0(0)>0$.
No simplicity or root-separation estimate is used.
\end{proof}

\begin{theorem}[NN-A022: angular rate and detecting-order lower bounds]
Assume A and define
\[
 \alpha(-a)=
 \sup_{\substack{H_{-a}(z)=0\\\operatorname{Re}z>0}}|\arg z|,
                          \qquad a>0.
\]
There are effectively computable constants $a_0,c,c_U,C>0$ such
that for $0<a\le a_0$,
\begin{equation}\label{eq:analytic-root-angle-rate}
                         \alpha(-a)\le C\sqrt a.
\end{equation}
More precisely the high zeros have a uniformly vanishing
scaled contribution:
\begin{equation}\label{eq:analytic-angular-tail-small}
 \lim_{R\to\infty}
 \sup_{0<a\le1}\ 
 \sup_{\substack{H_{-a}(z)=0,\ \operatorname{Re}z>0\\|z|\ge R}}
                       \frac{|\arg z|}{\sqrt a}=0,
\end{equation}
where the supremum over an empty set is zero.

For every integer $r\ge1$, all allowed coefficient Toeplitz
minors of $F_t$ of order at most $r$, at arbitrary indices,
are strictly positive throughout
\[
                         -\frac c{(r+1)^2}\le t\le0.
\]
For every $d\ge1$, $\chi_{t,d}$ from NN-A021 is of positive
type, with every polynomial Schur weight strictly positive,
throughout
\[
                         -\frac {c_U}{(d+1)^2}\le t\le0.
\]
Thus the least negative-minor order and the least failing
unitary rank satisfy
\[
                 m_*(-a)\gg a^{-1/2},
                 \qquad d_*(-a)\gg a^{-1/2}\qquad(a\downarrow0).
\]
The least detecting index window of NN-A014 also satisfies
$R_*(-a)\ge m_*(-a)-1\gg a^{-1/2}$, since a minor in a window
of size $R+1$ has order at most $R+1$.
Rodgers--Tao is needed only to know that these detecting
quantities are finite at every negative time.
\end{theorem}

\begin{proof}
The functions
$\log(r+3)/(r+1)$ and $\log^3(r+3)/(r+1)$ are bounded on
$[0,\infty)$ and tend to zero at infinity. The last lemma
therefore gives absolute effective constants $C_1,C_2$ with
\[
                  S(z)\le C_1q^{-2},\qquad
                  |L(z)|^2S(z)\le C_2q^{-4}
\]
on the stated cones. Choose an effective $K\ge1$ such that
$C_1/K^2\le2^{-10}$ and $C_2/K^4\le2^{-10}$.
For $q=K\sqrt a\le1/2$, Lemma~\ref{lem:analytic-gaussian-nonzero}
excludes every zero with $|\operatorname{Im}z|\ge q|z|$.
There are no imaginary-axis zeros, since $H_{-a}(iy)>0$.
Hence
\[
                 \alpha(-a)\le\arcsin(K\sqrt a)\le2K\sqrt a
                    \qquad(0<a\le1/(4K^2)).
\]
This proves \eqref{eq:analytic-root-angle-rate}.

For the sharper tail statement, fix $0<\varepsilon<1/2$.
In the cone with $q=\varepsilon\sqrt a$, $0<a\le1$, the
two parameters in the Gaussian nonvanishing test obey
\[
 A_z\le\frac C{\varepsilon^2}
               \frac{\log(r+3)}{r+1},\qquad
 B_z\le\frac{C^3}{\varepsilon^4}
               \frac{\log^3(r+3)}{r+1}.
\]
Both are at most $2^{-10}$ for $r\ge R_\varepsilon$,
uniformly in $a$. A zero of that size must therefore have
$|\arg z|\le\arcsin(\varepsilon\sqrt a)
\le2\varepsilon\sqrt a$. Let $\varepsilon$ decrease to zero
to obtain \eqref{eq:analytic-angular-tail-small}.

The quadratic factor criterion in NN-A020 applies whenever
$2\alpha(t)<\pi/(r+1)$. The displayed angular bound supplies
this inequality on $-c/(r+1)^2\le t\le0$ after choosing a
fixed smaller positive effective $c$. The infinite positive
particle paths give strictness for every allowed minor.
For example, with the preceding $K$, taking
$c\le\pi^2/(64K^2)$ gives a factor-two margin in the angle
inequality and lies inside the small-time range.

For the unitary statement retain the effective constants
$b<1$ and $\eta_0>0$ in the proof of NN-A021, so
$|\beta_j(t)|<b$ when $-\eta_0\le t\le0$. Then
\[
          |\arg\widetilde\beta_j(t)|
               \le\frac{2\alpha(t)}{1-b}
               \le\frac{4K\sqrt{|t|}}{1-b}.
\]
Choose $c_U>0$ with $c_U\le\eta_0$ and
$c_U\le\pi^2(1-b)^2/(64K^2)$. The same quadratic-factor
argument and the exact finite-rank dictionary of NN-A021
give the conclusion for every $d$. Taking the largest
integers $r,d$ permitted by these intervals gives the two
lower bounds for the detecting order and rank.
\end{proof}

\begin{corollary}[One explicit numerical certificate]
The constant in the $PF_r$ interval may be chosen as
$c=2^{-36}$. Thus, under A, every allowed minor of order at
most $r$ is strictly positive for
\[
                       -\frac{2^{-36}}{(r+1)^2}\le t\le0.
\]
For $0<a\le2^{-38}$ this implies
\[
 m_*(-a)\ge
       \left\lfloor\frac{2^{-18}}{\sqrt a}\right\rfloor,
 \qquad
 R_*(-a)\ge
       \left\lfloor\frac{2^{-18}}{\sqrt a}\right\rfloor-1.
\]
The numerical input is the existing certified inequality
$A_0<1$ from NN-A014; no new floating-point experiment is used.
\end{corollary}

\begin{proof}
The coefficient sum from NN-A014 obeys
\[
 A_0=\sum_{k\ge0}(2k+2)(2k+1)c_{k+1}(0)
                              \ge2c_1(0)=2p_1(0).
\]
Its saved Arb certificate gives $0<p_1(0)<1/2$. For real
$|x|\le1$, the elementary inequality $\cos v\ge1-v^2/2$
therefore implies
\[
 H_0(x)\ge H_0(0)(1-p_1(0)x^2)>0.
\]
In particular $x_1>1$, and $\sum_jx_j^{-2}=p_1(0)<1/2$.

Write $N_\zeta(T)$ for the positive-ordinate zeta count.
The exact source bound, for $T\ge e$, is
\[
 \left|N_\zeta(T)-\frac{T}{2\pi}
                      \log\frac{T}{2\pi e}\right|
       \le .1038\log T+.2573\log\log T+9.3675.
\]
This is Hasanalizade--Shen--Wong, Corollary~1.2,
equation~(1.5), printed p.~2. At $T=e$ the main term is
negative, so $N_\zeta(e)<10$. For $T\ge e$, the main term
is at most $T\log(T+3)$, and the error is less than
\[
 .362\log T+9.368
       \le .362T+4T
       \le4.362T\log(T+3).
\]
Here $\log\log T\le\log T$, $e>5/2$, and
$\log(T+3)>1$ suffice. Hence
$N_\zeta(T)\le6T\log(T+3)$. Together with $x_1>1$ and
$N(X)=N_\zeta(X/2)$, these bounds give the deliberately loose
but uniform inequality
\begin{equation}\label{eq:analytic-explicit-count-100}
                         N(X)\le100X\log(X+3)\qquad(X>0).
\end{equation}
Indeed for $X<1$ the left side is zero, for
$1\le X<2e$ it is less than ten, and for $X\ge2e$
it is at most $3X\log(X+3)$.

We now make every constant in the cone lemma explicit.
For $r\ge1/2$, the range $x_j<r/2$ contributes at most
\[
       400\frac{\log(r+3)}r\quad\hbox{to }S,
       \qquad \frac{400}3\log(r+3)\quad\hbox{to }|L|.
\]
The middle range contributes at most
\[
       800q^{-2}\frac{\log(r+3)}r\quad\hbox{to }S,
       \qquad800q^{-1}\log(r+3)\quad\hbox{to }|L|,
\]
using $\log(2r+3)\le2\log(r+3)$.
For the last range, partial summation and
$\log(2r+3)>1$ give
\begin{align*}
 \sum_{x_j>2r}x_j^{-2}
 &\le200\int_{2r}^\infty\frac{\log(u+3)}{u^2}\,du\\
 &\le200\,\frac{\log(2r+3)+1}{2r}
 \le400\,\frac{\log(r+3)}r.
\end{align*}
It contributes at most
$3200\log(r+3)/r$ to $S$ and
$(3200/3)\log(r+3)$ to $|L|$.
Thus
\[
 S\le4400q^{-2}\frac{\log(r+3)}r
       \le13200q^{-2}\frac{\log(r+3)}{r+1},
       \qquad |L|\le2000q^{-1}\log(r+3).
\]
For $r\le1/2$, $x_1>1$ and $p_1(0)<1/2$ give
$S<4$ and $|L|\le4r/3\le2/3$, which satisfy the same
bounds after taking $C=2^{14}$ in the cone lemma.

For every $r\ge0$,
\[
 \frac{\log(r+3)}{r+1}\le2,\qquad
 \frac{\log^3(r+3)}{r+1}\le16.
\]
For the second inequality put $y=r+3\ge3$, note
$r+1=y-2\ge y/3$, and maximize $(\log y)^3/y$:
its maximum is $27/e^3$, so the ratio is at most
$81/e^3<81/8<16$. The first follows from
$\log(r+3)\le r+2$.
Consequently one can use $C_1=2^{15}$ and $C_2=2^{46}$
in the proof of the theorem. With $K=2^{16}$,
\[
             C_1/K^2=2^{-17},\qquad C_2/K^4=2^{-18},
\]
both smaller than the required $2^{-10}$. Thus
\[
 \alpha(-a)\le2^{17}\sqrt a\qquad(0<a\le2^{-34}).
\]
If $a\le2^{-36}/(r+1)^2$, then
$2\alpha(-a)\le1/(r+1)<\pi/(r+1)$.
The quadratic factor criterion proves the explicit interval.
Finally take
$r=\lfloor2^{-18}/\sqrt a\rfloor-1\ge1$.
All minors of order at most this $r$ are positive, so
$m_*(-a)\ge r+1$, and $R_*(-a)\ge m_*(-a)-1$.
\end{proof}

\begin{corollary}[The limiting angle detects endpoint multiplicities]
For each distinct positive zero $x$ of $H_0$, let $m_x$ be its
multiplicity, and let $h_m$ be the largest absolute value of a
root of the probabilists' Hermite polynomial $\operatorname{He}_m$;
in particular $h_1=0$. Under A,
\begin{equation}\label{eq:analytic-endpoint-angular-limit}
       \lim_{a\downarrow0}\frac{\alpha(-a)}{\sqrt a}
          =\kappa_0,\qquad
       \kappa_0=\sup_{H_0(x)=0,\ x>0}\frac{\sqrt2\,h_{m_x}}x.
\end{equation}
The supremum is finite and is attained if it is positive.
In particular all nontrivial zeta zeros are simple if and only
if $\alpha(-a)=o(\sqrt a)$. Under that additional simplicity
condition the two detection bounds improve to
\[
                     \sqrt a\,m_*(-a)\longrightarrow\infty,
              \qquad \sqrt a\,d_*(-a)\longrightarrow\infty.
\]
This does not determine the height of the first nonreal zero.
\end{corollary}

\begin{proof}
The Hermite recurrence gives its roots as the eigenvalues of
the real symmetric tridiagonal matrix with off-diagonal
entries $\sqrt1,\ldots,\sqrt{m-1}$. Gershgorin's bound yields
$h_m\le2\sqrt m$. Since
$m_x\le N(2x)\ll x\log(x+3)$, it follows that
$h_{m_x}/x\to0$ as $x\to\infty$. There are finitely many
distinct zeros in each compact set, so $\kappa_0$ is finite
and any positive supremum is a maximum.

At each fixed positive zero $x$, NN-A006 gives the complete
local splitting for negative time:
\[
             z_j(-a)=x+i\sqrt{2a}\,h_{m_x,j}+O(a).
\]
Consequently the largest angle of these local branches,
divided by $\sqrt a$, tends to $\sqrt2\,h_{m_x}/x$.
This also covers a simple zero, whose branch remains real
for small real time.

Choose a radius $R$ whose boundary contains no zero of $H_0$.
Uniform compact convergence, local root splitting and
Rouch\'e's theorem account for every zero in $|z|\le R$.
There are finitely many endpoint germs there, so the scaled
maximum of their angles tends to the maximum of their
Hermite constants. The remaining zeros contribute arbitrarily
little, uniformly for $0<a\le1$, by
\eqref{eq:analytic-angular-tail-small}. Combining the finite
maximum with this vanishing tail proves
\eqref{eq:analytic-endpoint-angular-limit}.
Since $h_m>0$ precisely when $m\ge2$, $\kappa_0=0$ is
equivalent to simplicity of every endpoint zero.

If $\alpha(-a)/\sqrt a\to0$, then for any prescribed $M>0$
the quadratic angle criterion applies for all orders
$r\le M/\sqrt a$ once $a$ is sufficiently small. The
M\"obius angle bound from NN-A021 gives the same assertion
at unitary ranks $d\le M/\sqrt a$. Since $M$ is arbitrary,
the two strengthened lower bounds follow.
\end{proof}

\paragraph{Exact geometric and representation meaning.}
The multiplicity $m_x$ is the length of the completed local
zero algebra. The parameter change $t=-u^2/2$ splits its
initial heat polynomial into the distinct Hermite directions;
the positive real direction $x$ converts their imaginary
normal displacement into the coefficient
$\sqrt2\,h_{m_x}/x$. The global theorem shows that these
finite local Artin lengths determine the limiting maximal
angle, because the high-radius contribution vanishes uniformly.
The lower bounds concern every skew Schur test of bounded
width and every polynomial irreducible at bounded unitary
rank, even when its degree is unrestricted. Their
exterior-power dictionary is SGA~6, Expos\'e~V,
Example~2.2(b) and Lemma~2.2.1, with the exact finite-rank
equivalence proved in NN-A021. The bounds are for these
specified certificate families; no general complexity lower
bound, pure arithmetic motive, or full collision-height law
is asserted.


\section{An isolated pair: collision time and the surrounding zeros}
\label{sec:analytic-isolated-collision}

\begin{theorem}[NN-A023: an isolated-pair collision with its first interaction correction]
Let $f$ be a real entire function of order at most one, with
only real zeros, and let $b-\delta/2,b+\delta/2$ be adjacent
simple zeros, $\delta>0$. Put
\[
 G(z)=\frac{f(z)}{(z-b)^2-\delta^2/4},\qquad
 L=\frac{G'}G(b)\in\mathbb R,\qquad
 S=\sum_{G(\xi)=0}\frac1{(b-\xi)^2},\qquad
                       \kappa=\delta^2S.
\]
Zeros in $S$ are counted with multiplicity. Assume the
order-one Hadamard product for $G$; its possible linear
exponential factor is allowed and is retained in $L$.
The case $S=0$ is included.

For all sufficiently small $\kappa$, with an absolute
threshold, the two zeros continued from this pair under
$f_{-a}=e^{a\partial_z^2}f$ undergo a unique double collision
on the local time scale $0\le a\le\delta^2/4$. Its time and
real position satisfy
\begin{align}
 a_c&=\delta^2\left(\frac18+\frac{\kappa}{32}
                                      +O(\kappa^2)\right),
                        \label{eq:analytic-isolated-time}\\
 z_c&=b-2a_cL+O(\delta\kappa^{3/2}).
                        \label{eq:analytic-isolated-center}
\end{align}
The $O$ constants are absolute. Before the collision the
tracked zeros are real, and after it they are a nonreal
conjugate pair. The local neighbourhood moves with the real
center $b-2aL$; no smallness assumption on $\delta L$ is needed.

Under A this applies to every adjacent simple pair of $H_0$.
In particular any sequence of such arithmetic pairs with
$\kappa\to0$ has actual collision times
$t_c=-\delta^2(1+\kappa/4+O(\kappa^2))/8$.
This is a conditional statement about the stated pair
parameter, not a proof of how often these pairs occur.
\end{theorem}

\begin{proof}
Hadamard factorization has no exponential term of degree two
or higher, by the order assumption. Therefore
\[
 R(h)=\frac{G(b+h)}{G(b)e^{Lh}}
       =\prod_{G(\xi)=0}
         \left(1+\frac h{b-\xi}\right)
                           \exp\left(-\frac h{b-\xi}\right)
\]
converges locally uniformly, and the proof of
Lemma~\ref{lem:analytic-product-remainder} gives
$|R(h)|\le e^{S|h|^2/2}$. Its first Taylor coefficients are
\begin{equation}\label{eq:analytic-pair-R-expansion}
 R(h)=1-\frac S2h^2+\frac{T_3}{3}h^3+\mathcal E_4(h),
          \qquad T_3=\sum_{G(\xi)=0}(b-\xi)^{-3},
          \qquad |T_3|\le S^{3/2}.
\end{equation}
There is an absolute constant $C$ such that
\begin{equation}\label{eq:analytic-pair-R-fourth}
 |\mathcal E_4(h)|\le
               CS^2|h|^4e^{2S|h|^2}\qquad(h\in\mathbb C).
\end{equation}
Indeed, after putting $v=h\sqrt S$, Cauchy's estimate on
$|v|=1$ bounds the remainder by $2e^{1/2}|v|^4$ for
$|v|\le1/2$. For $|v|\ge1/2$, use the global growth bound,
the explicit quadratic coefficient and
$|T_3|/(3S^{3/2})\le1/3$; each of the four terms is bounded
by a fixed multiple of $|v|^4e^{2|v|^2}$.
If $S=0$, the product is empty and $R=1$, so the claimed
collision and formulas are exact.

Write $a=\delta^2\tau$ and use the moving coordinate
\[
                         z=b-2aL+\delta w.
\]
Completing the Gaussian square, or conjugating the heat
operator by $e^{L(z-b)}$, gives the exact identity
\begin{equation}\label{eq:analytic-pair-scaled-family}
 \frac{f_{-\delta^2\tau}(b-2\delta^2\tau L+\delta w)}
      {\delta^2G(b)\exp(\delta wL-\delta^2\tau L^2)}
  =J(\tau,w)
  =\mathbb E\left[(Y^2-\tfrac14)R(\delta Y)\right],
                \qquad Y=w+\sqrt{2\tau}\,V,\quad V\sim N(0,1).
\end{equation}
For real $\tau\ge0$ this is an ordinary Gaussian convolution
with a real shift $2aL$. Both sides extend holomorphically
in $\tau$ to a complex neighbourhood of $[0,1/4]$:
one may choose a square-root branch locally away from zero;
the expectation is even in that square root and extends
across zero. The order-one growth and Gaussian decay
justify the continuation. The denominator is nonzero.

For bounded complex $w,\tau$, the growth bound for $R$
is dominated by $\exp(C\kappa(|V|^2+1))$.
Consequently \eqref{eq:analytic-pair-R-fourth}, integrated
against the Gaussian, gives a uniform holomorphic expansion
with an $O(\kappa^2)$ remainder:
\begin{align}
 J(\tau,w)
  &=w^2+2\tau-\tfrac14-\frac{\kappa}{2}P_4(\tau,w)
                  +\eta_3 P_5(\tau,w)+O(\kappa^2),
                    \label{eq:analytic-pair-J-expansion}\\
 \eta_3&=\delta^3T_3/3,\qquad |\eta_3|\le\kappa^{3/2}/3,\nonumber\\
 P_4(\tau,w)
  &=w^4+(12\tau-\tfrac14)w^2+12\tau^2-\tfrac12\tau,\nonumber\\
 P_5(\tau,w)
  &=w^5+(20\tau-\tfrac14)w^3
                         +(60\tau^2-\tfrac32\tau)w.\nonumber
\end{align}
All derivatives on smaller compact sets have the same
remainder order by Cauchy's formula. In particular the
quadratic first correction is even in $w$, while the
cubic correction is odd. Uniform constants depend only
on the fixed compact sets, not on $f,\delta,L,S$.

For $0\le\tau\le1/4$, the limit polynomial
$w^2+2\tau-1/4$ has exactly two roots in $|w|<2$ and
none on the boundary. Rouch\'e's theorem therefore tracks
exactly two roots of $J$ in this moving disc for small
$\kappa$, starting with the given pair at $\tau=0$.
Its derivative is $2w+O(\kappa)$ uniformly there and has
a single zero near zero. The implicit function theorem
gives this zero as a real analytic function $w_*(\tau)$
for real $\tau$, with
$w_*(\tau)=O(\kappa^{3/2})$: evaluate
$\partial_wJ(\tau,0)$ using the parity in
\eqref{eq:analytic-pair-J-expansion}, and use
$\partial_w^2J=2+O(\kappa)$.

Set $D(\tau)=J(\tau,w_*(\tau))$. The same expansion yields
\[
 D(\tau)=2\tau-\tfrac14
              -\frac{\kappa}{2}(12\tau^2-\tfrac12\tau)
              +O(\kappa^2),\qquad D'(\tau)=2+O(\kappa).
\]
It has one real zero $\tau_c$ near $1/8$. First
$\tau_c=1/8+O(\kappa)$; evaluating the correction at $1/8$,
where $12\tau^2-\tau/2=1/8$, then gives
\[
                       \tau_c=\frac18+\frac{\kappa}{32}
                                                  +O(\kappa^2).
\]
The vanishing conditions $J=J_w=0$, together with
$J_{ww}=2+O(\kappa)\ne0$, give a double zero, not a higher
collision. Since $D'$ and $J_{ww}$ are positive for real
arguments here, its two branches are real for
$\tau<\tau_c$ and nonreal conjugate for $\tau>\tau_c$.
No other local collision is possible on $[0,1/4]$:
the derivative has only the one critical point and $D$
has only the one zero. Returning to $a,z$ proves
\eqref{eq:analytic-isolated-time} and
\eqref{eq:analytic-isolated-center}. This proof also
accounts for the full pair throughout its local interval;
it does not merely locate a nonreal zero at one time.
\end{proof}

\paragraph{Relation to the classical Lehmer criterion.}
Csordas--Smith--Varga, \emph{Lehmer Pairs of Zeros, the
de Bruijn--Newman Constant $\Lambda$, and the Riemann
Hypothesis}, Constructive Approximation~10 (1994), 107--129,
Definition~1, equations~(1.10)--(1.12), printed pp.~109--110,
uses the endpoint sum
\[
 g=\sum_{\xi\ne b-\delta/2,b+\delta/2}
       \left\{(b-\delta/2-\xi)^{-2}
                       +(b+\delta/2-\xi)^{-2}\right\}.
\]
Under A this is positive. For $\kappa\to0$,
$g=2S+O(\delta^2S^2)$, by expanding the two summands
symmetrically and using $\sum(b-\xi)^{-4}\le S^2$.
Thus $\delta^2g=2\kappa+O(\kappa^2)$, and the pairs above
eventually satisfy their $\delta^2g<4/5$ criterion.
Their Theorem~1, equations~(1.13)--(1.14), printed p.~110,
states the lower bound
\[
 \frac{(1-\tfrac54\delta^2g)^{4/5}-1}{8g}\le\Lambda.
\]
It does not assert that this displayed lower bound is an
exact collision time. Its first two expansion terms agree
with \eqref{eq:analytic-isolated-time} after changing the
sign of time. The argument here identifies an actual
unique local double collision and its interaction correction.
The original printed pages were inspected directly.

\paragraph{Geometric and representation meaning.}
The chosen two-point divisor determines a finite local
rank-two zero family in the moving coordinate. Its
discriminant has a simple zero at $\tau_c$, so after a
quadratic base change the two branches are separated.
For the resulting finite locally free rank-two algebra,
the exact \'{e}tale/discriminant criterion is SGA~1,
Expos\'e~I, section~4, Proposition~10; section~3,
Proposition~1 gives the equivalent vanishing of relative
differentials. This criterion is applied to the actual
local finite algebra, not to all of the infinite divisor.
The first correction is the inverse-square trace of the
complementary zero divisor, while its logarithmic first
trace produces only the displayed real translation.
This is a concrete local discriminant computation; it
does not make the complementary analytic divisor into a
finite pure Weil object. In the Toeplitz/character
dictionary of NN-A020--A022, each such collision creates
one conjugate pair of complex product parameters.
The relevant exterior-power operations have the locator
SGA~6, Expos\'e~V, Example~2.2(b) and Lemma~2.2.1.
No integral canonical-basis statement follows from
the analytic inverse-square parameter.


\section{The first CUE heat collision is an extreme-gap event}
\label{sec:analytic-cue-first-collision}

Let $0\le\theta_1<\cdots<\theta_N<2\pi$ be the eigenangles
of Haar $U(N)$, and extend them periodically by
$\theta_{j+N}=\theta_j+2\pi$. The real characteristic
trigonometric polynomial is
\[
 Z_N(z)=\prod_{j=1}^N2\sin\frac{z-\theta_j}{2},\qquad
 Z_{N,-a}=e^{a\partial_z^2}Z_N .
\]
For odd $N$ it is antiperiodic, and its zero divisor is
still $2\pi$-periodic. Negative heat time has exactly the
same sign convention as $H_{-a}$ in the arithmetic family.
Define
\[
 \delta_N=\min_j(\theta_{j+1}-\theta_j),\qquad
 A_N=\inf\{a>0:Z_{N,-a}\text{ has a nonreal zero}\}.
\]
This section concerns this specified CUE model, with no
arithmetic transfer assumption.

\begin{theorem}[NN-A024: first CUE collision law]
For $N\ge2$, $A_N$ is positive and finite almost surely.
As $N\to\infty$,
\[
                    \frac{8A_N}{\delta_N^2}
                                  \longrightarrow1
                                  \quad\hbox{in probability}.
\]
More precisely the ratio is $1+O_{\mathbb P}(N^{-1/3})$.
With probability tending to one, the first collision is
the one continued from the initially smallest cyclic gap.
For $u\ge0$,
\begin{equation}\label{eq:analytic-cue-first-time-law}
 \lim_{N\to\infty}
 \mathbb P\left\{
       \frac{8N^{8/3}A_N}{(72\pi)^{2/3}}>u\right\}
                              =e^{-u^{3/2}}.
\end{equation}
Equivalently, after unfolding the mean spacing to one,
the first negative heat time is of order $N^{-2/3}$,
smaller than the mean-spacing-squared scale, which is one.
\end{theorem}

\begin{proof}
We first prove a deterministic global lower bound for
the first loss of reality. Write
\[
 Z_N(z)=\sum_{k=0}^Nc_ke^{i(k-N/2)z}.
\]
The coefficients $c_0,c_N$ are nonzero. Under
$e^{a\partial_z^2}$ the $k$th coefficient is multiplied
by $e^{-a(k-N/2)^2}$. After multiplication by
$e^{iNz/2}$ this is a degree-$N$ polynomial in $w=e^{iz}$.
Its leading and constant coefficients remain nonzero,
and on a compact $a$ interval the normalized polynomial
coefficients and their reciprocals at the two ends stay
bounded. The $N$ roots therefore remain in a compact
annulus in $\mathbb C^*$.

While all these roots are on the unit circle and simple,
their angles have real analytic continuations $y_j(a)$.
There is no loss of a zero at imaginary infinity at a
finite time; and a simple real zero remains real by the
implicit function theorem and conjugation. Thus a first
loss of reality must occur at a finite real collision.
Before that collision the factorization into real sine
factors gives
\begin{equation}\label{eq:analytic-circular-backward-ode}
 y_i'(a)=-\sum_{j\ne i}\cot\frac{y_i(a)-y_j(a)}2.
\end{equation}
For an empty adjacent arc of length
$g_i=y_{i+1}-y_i\in(0,2\pi)$, the two mutual terms give
$-2\cot(g_i/2)$ to $g_i'$. Every remaining cotangent
difference is negative before the minus sign: along
this empty arc the cotangent is strictly decreasing
and crosses no pole. Therefore
\[
                  g_i'(a)\ge-2\cot(g_i(a)/2),
                  \qquad (g_i^2)'(a)\ge-8.
\]
There is a slightly sharper integrated form:
\[
       \frac d{da}\cos(g_i(a)/2)
                       \le\cos(g_i(a)/2).
\]
If this gap closes at $A$, its left side integrates to
$e^{-A}\le\cos(g_i(0)/2)$. Since
$g_i(0)\ge\delta_N$ and $\delta_N<\pi$ almost surely,
\begin{equation}\label{eq:analytic-first-time-lower}
       A_N\ge-\log\cos(\delta_N/2)\ge\delta_N^2/8.
\end{equation}
This rules out an earlier collision anywhere on the
circle, not merely near the smallest pair.

For completeness, the first collision is almost surely
finite. When $N$ is even the constant Fourier coefficient
is almost surely nonzero and eventually dominates all
others, so for large $a$ there are no real zeros.
When $N$ is odd, the almost surely nonzero frequencies
$\pm1/2$ eventually dominate, leaving one simple real
zero modulo $2\pi$, fewer than $N$ for $N\ge3$.
Uniform $C^1$ convergence after rescaling gives this
last zero count. Each asserted Fourier coefficient is
a nonzero real analytic function of the eigenangles:
one can check nonvanishing when the angles are all
close to zero, where the limiting function is
$(2\sin(z/2))^N$. Its vanishing set has Lebesgue
measure zero, and the CUE law has an absolutely
continuous density. The preceding annulus argument
then supplies a finite first collision.

We next prove a uniform isolation estimate for the
smallest pair. Put $\ell_N=N^{-7/6}$. With probability
tending to one, every circular interval of length
$\ell_N$ contains at most two eigenangles.
Here is a direct determinantal proof. Let
\[
 v_N(x)=(1,e^{ix},\ldots,e^{i(N-1)x})/\sqrt{2\pi}.
\]
The three-point intensity is
$\rho_3(x_1,x_2,x_3)=
 \|v_N(x_1)\wedge v_N(x_2)\wedge v_N(x_3)\|^2$.
Successive divided differences give
\[
 \rho_3(x_1,x_2,x_3)
       \le C N^9\prod_{i<j}|x_i-x_j|^2.
\]
Indeed the norms of $v_N,v_N',v_N''$ are respectively
$O(N^{1/2}),O(N^{3/2}),O(N^{5/2})$. The determinant
of the three divided-difference columns extracts the
Vandermonde product, and their norms are bounded by
these derivative norms (with a factor $1/2$ in the
second divided difference). This applies to real
lifts of any short arc, uniformly around the circle.
The expected number of triples lying in a forward
arc of length $\ell_N$ is therefore at most
\[
 C N^9\ell_N^8=O(N^{-1/3}).
\]
If such an arc contains three points, one of its points
can be chosen as the start of a forward arc containing
the other two. Markov's inequality proves the claim.

The classical extreme-gap theorem of Ben Arous--Bourgade
gives
\begin{equation}\label{eq:analytic-CUE-min-gap}
 \frac{N^{4/3}\delta_N}{(72\pi)^{1/3}}
          \ \Longrightarrow\ Y,\qquad
                    \mathbb P\{Y>y\}=e^{-y^3}\quad(y\ge0).
\end{equation}
In particular $\delta_N=O_{\mathbb P}(N^{-4/3})$ and
$\delta_N/\ell_N\to0$ in probability.
Choose the minimum pair, using real lifts
$b-\delta_N/2,b+\delta_N/2$. On the preceding isolation
event, for $\delta_N\le\ell_N/2$ every other zero on the
periodic real line has distance at least $\ell_N/2$
from $b$. Each interval of length $\ell_N/2$ contains
at most two zeros. Summing successive distance shells
therefore gives
\[
 S_N=\sum_{\substack{Z_N(\xi)=0\\
                   \xi\ne b-\delta_N/2,b+\delta_N/2}}
                     (b-\xi)^{-2}
                          \le C\ell_N^{-2}.
\]
The periodic copies are included in this sum. Thus the
pair parameter satisfies
\[
                    \kappa_N=\delta_N^2S_N
                             =O_{\mathbb P}(N^{-1/3}).
\]
This includes the exceptional probability
$O(N^{-1/3})$ of the isolation event failing:
outside that event no bound is asserted, while its
probability tends to zero, as required for the stated
bound in probability.

The real entire function $Z_N$ has order one and only
real zeros. Removing the selected pair leaves an
order-one function with the required canonical product.
NN-A023 applies, including its unrestricted real drift.
It gives an actual collision of that pair at
\[
 a_{\min}=\frac{\delta_N^2}{8}
                    \{1+\kappa_N/4+O(\kappa_N^2)\}.
\]
The first global collision can occur no later than this
one. Combine this upper bound with
\eqref{eq:analytic-first-time-lower} to obtain
$8A_N/\delta_N^2=1+O_{\mathbb P}(N^{-1/3})$.

To identify the colliding pair, let $\delta_N^{(2)}$
be the second-smallest initial gap. The same source's
joint extreme-gap theorem implies that
$\delta_N^{(2)}/\delta_N$ converges in distribution
to a variable strictly larger than one almost surely.
Since $\kappa_N\to0$ in probability, with probability
tending to one
\[
                    (\delta_N^{(2)})^2/8>a_{\min}.
\]
Every other initially labelled cyclic gap has its
first possible closing time at least its initial
squared length divided by eight, by the universal
gap inequality. Hence none can collide first before
$a_{\min}$, proving the assertion about the smallest pair.
Finally combine \eqref{eq:analytic-CUE-min-gap} with
$8A_N/\delta_N^2\to1$ and the continuous mapping theorem:
the variable in \eqref{eq:analytic-cue-first-time-law}
converges to $Y^2$, whose survival function is
$e^{-u^{3/2}}$.
\end{proof}

\paragraph{Exact two-angle normalization check.}
For just two angles $b\pm\delta/2$,
\[
 Z_2(z)=2\cos(\delta/2)-2\cos(z-b),\qquad
 Z_{2,-a}(z)=2\cos(\delta/2)-2e^{-a}\cos(z-b).
\]
Consequently
\[
              A_2=-\log\cos(\delta/2)
                  =\delta^2/8+\delta^4/192+O(\delta^6)
\]
when $0<\delta<\pi$. The complementary periodic zero
sum in NN-A023 is
\[
          S=\frac1{2\sin^2(\delta/4)}-\frac8{\delta^2}
                                  =\frac16+O(\delta^2).
\]
Its correction $\delta^2\kappa/32$ is exactly
$\delta^4/192$ to first order. This checks both the
heat-sign normalization and the interaction coefficient
by a closed-form example.

\paragraph{Primary locator and arithmetic scope.}
The only imported random-matrix asymptotic is
Ben Arous--Bourgade, \emph{Extreme gaps between eigenvalues
of random matrices}, Annals of Probability~41 (2013),
2648--2681, DOI~10.1214/11-AOP710, Theorem~1.1
printed p.~2650 and Corollary~1.2 p.~2651.
Their full-circle gap intensity is
$u^2\,du/(24\pi)$ for the scale $N^{4/3}\delta$,
giving the constant $72\pi$ in
\eqref{eq:analytic-CUE-min-gap}. The no-triples estimate,
heat dynamics, global earliest-collision bound and
passage from gaps to collision times were proved above.
The theorem uses neither A nor M.

The model also distinguishes two different heuristic
questions in the supplied consultation. Matching
$\sqrt a$ to the mean arithmetic gap gives a
\emph{typical-gap} or bulk scale
$\log X\asymp a^{-1/2}$. It is not the scale of the
first rare collision in the CUE model. If one
additionally postulates the analogous cubic small-gap
Poisson law for arithmetic zeros, together with
isolation and exhaustion of early collisions by
those pairs, the heat-coordinate density
$D(x)\sim\log x/(4\pi)$ predicts a count
\[
 \frac{\pi^2}{9}(8a)^{3/2}\int_0^X D(x)^4\,dx
           \sim\frac{\sqrt2}{144\pi^2}
                              a^{3/2}X(\log X)^4.
\]
The first-event scale in that stronger model is
$X\asymp a^{-3/2}(\log(1/a))^{-4}$, rather than an
exponential scale. This displayed arithmetic conversion
is explicitly a conditional model prediction.
Neither A, Montgomery's restricted pair-correlation
theorem, nor a fixed-scale sine-kernel limit supplies
the required extreme-gap process or the exhaustion
of all early arithmetic collisions. NN-A023 proves
the local dynamics once the isolation parameter is small.

\paragraph{An actual algebraic family and its discriminant.}
Put $q=e^{-a/4}$ and $w=e^{iz}$. The preceding Fourier
formula gives the polynomial
\[
              P_N(q,w)=
                    \sum_{k=0}^Nc_kq^{(2k-N)^2}w^k
                 \in\mathbb C[q,q^{-1},w,w^{-1}].
\]
Its leading and constant coefficients are units over
$\mathbb C[q,q^{-1}]$. Dividing by the leading
coefficient exhibits
\[
 \operatorname{Spec}
       \mathbb C[q,q^{-1},w,w^{-1}]/(P_N)
                      \longrightarrow\mathbb G_{m,q}
\]
as a finite free family of degree $N$; inversion of $w$
is automatic because its constant coefficient is a unit.
It is \'{e}tale precisely off its polynomial
discriminant, by the derivative or trace-pairing
criterion. The exact Grothendieck locator is
SGA~1, Expos\'e~I, section~4, Proposition~10
(invertibility of the discriminant for a finite
locally free cover), with Definition~1 in that
section and section~3, Proposition~1 for the
flat/unramified and differential formulations.
At the first collision identified above with
probability tending to one, the local rank-two
component has a simple discriminant zero and
ramified double-cover normalization as in NN-A023.
The circle is the fixed locus of
$w\mapsto1/\overline w$ on the real parameter ray.
This is a genuine complex algebraic zero family
for each CUE sample; it is not an arithmetic
Frobenius construction or a canonical-basis model.


\section{The complete early CUE collision process}
\label{sec:analytic-cue-collision-process}

The first-event result NN-A024 does not by itself control the
other zeros after its first pair becomes nonreal. The following
uniform local argument supplies that missing exhaustion step.
The heat normalization is unchanged:
$Z_{N,-a}=e^{a\partial_z^2}Z_N$ and
$Z_N(z)=\prod_{j=1}^N2\sin((z-\theta_j)/2)$.

\begin{lemma}[Uniform one-root and two-root tracking]
\label{lem:analytic-uniform-clusters}
Let $f$ be real entire of order at most one with all zeros real.
Remove either a simple zero $b$, or two adjacent simple zeros
$b\pm\delta/2$, and call the remaining entire factor $G$.
Put $L=G'(b)/G(b)$ and $S=\sum_{G(\xi)=0}(b-\xi)^{-2}$.
Possible linear exponential factors are included in $L$.
There are absolute positive constants $c,C$ with the following
properties for $T>0$.

If one zero was removed and $TS<c$, then for every $a\in[0,T]$
exactly one zero of $e^{a\partial_z^2}f$ lies in the moving disc
\[
                  |z-(b-2aL)|<2\sqrt T.
\]
It is real and simple throughout the interval.

If two zeros were removed, put $R^2=\delta^2+T$. If
$R^2S<c$, exactly two zeros lie in
$|z-(b-2aL)|<2R$ for every $a\in[0,T]$. They have at most
one collision there. The collision is double and crosses from
two real zeros to a nonreal conjugate pair; the pair cannot
return to the real axis during $[0,T]$. Its collision time,
when in this interval, is the time supplied by NN-A023:
\[
             a_c=\frac{\delta^2}{8}
                \bigl(1+\delta^2S/4+O((\delta^2S)^2)\bigr).
\]
The constants are independent of $L$, and of the ratio
$\delta^2/T$.
\end{lemma}

\begin{proof}
Use the canonical remainder
\[
 \mathcal R(h)=G(b+h)/(G(b)e^{Lh}),\qquad
 |\mathcal R(h)|\le e^{S|h|^2/2},\qquad
 |\mathcal R(h)-1|\le 8S|h|^2e^{2S|h|^2},
\]
proved in NN-A022. Conjugating the Gaussian heat operator by
$e^{L(z-b)}$ removes the common drift exactly. In the one-root
case, with $z=b-2aL+\sqrt T\,u$ and $\tau=a/T$, division by
$\sqrt T G(b)e^{\sqrt T uL-aL^2}$ gives
\[
 I(\tau,u)=\mathbb E\bigl[Y\mathcal R(\sqrt T Y)\bigr],
                    \qquad Y=u+\sqrt{2\tau}\,V,
                    \qquad V\sim N(0,1).
\]
On any fixed bounded complex $u,\tau$ set this is
$I=u+O(TS)$, with every derivative on smaller compact sets
bounded in the same way. To justify this including complex
$\tau$, the Gaussian expectation is even in its square-root
parameter. Its absolute integrand is dominated by
$C(1+|V|)^C e^{C TS|V|^2}$; for small $TS$ the usual Gaussian
density dominates it uniformly. The expression therefore
extends holomorphically through $\tau=0$, and the stated
bounds and their derivatives follow by Cauchy's formula.
Rouch\'e on $|u|=2$ gives exactly one zero. The derivative
$I_u=1+O(TS)$ there is nonzero at that zero. Real conjugation
and uniqueness make the zero real for real $\tau$.

For the two-root case scale by $R$ instead, put
$d=\delta/R\in(0,1]$, $\tau=a/R^2\in[0,1]$, and divide by
$R^2G(b)e^{RuL-aL^2}$. The resulting exact family is
\[
 J(\tau,u)=\mathbb E\bigl[(Y^2-d^2/4)\mathcal R(RY)\bigr]
            =u^2+2\tau-d^2/4+O(R^2S).
\]
The same Gaussian argument is uniform in $0<d\le1$ and
controls all derivatives on fixed compact sets. On $|u|=2$
for real $0\le\tau\le1$, the leading polynomial has modulus
at least $4-2-1/4=7/4$. Thus exactly two zeros remain inside.
Its derivative has exactly one zero there, near zero, and
$J_{uu}=2+O(R^2S)>0$ on the real portion of this disc.
The critical point $u_*(\tau)$ is real analytic. Define
$D(\tau)=J(\tau,u_*(\tau))$. Because $J_u=0$ at $u_*$,
\[
                         D'(\tau)=J_\tau(\tau,u_*)
                                      =2+O(R^2S)>0.
\]
At $\tau=0$ the two roots are exactly $\pm d/2$; strict
convexity gives $D(0)<0$, even when $d$ is very small.
Consequently $D$ has at most one zero in this time interval.
At that zero $J_{uu}$ and $D'$ are positive and nonzero, so
it is a simple discriminant zero with the stated real-to-complex
crossing. The two roots cannot return without another zero of
$D$. Applying NN-A023 on its own local time interval identifies
the collision time; whenever both intervals apply they track
the same two zeros by their common initial sections and local
uniqueness. No condition on $\delta L$ or $\sqrt T L$ was used.
\end{proof}

\begin{theorem}[NN-A025: Poisson process of early CUE collisions]
\label{thm:analytic-cue-collision-process}
For Haar CUE, let $(a_{N,j},b_{N,j})$ be the real double-zero
collisions of $Z_{N,-a}$, with $a_{N,j}>0$ and real position
$b_{N,j}$ taken modulo $2\pi$. On any fixed rescaled time
window, with probability tending to one, all multiple-zero
events are such simple double collisions; each creates one
nonreal conjugate pair, and these pairs remain nonreal on that
window. The marked point process
\[
 \sum_j\delta_{(v_{N,j},\,b_{N,j}\bmod2\pi)},
 \qquad v_{N,j}=\frac{8N^{8/3}a_{N,j}}{(72\pi)^{2/3}},
\]
converges on bounded time windows to a Poisson point process
with intensity
\begin{equation}\label{eq:analytic-collision-ppp-intensity}
                   \frac32\sqrt v\,dv\,\frac{db}{2\pi}.
\end{equation}
In particular the number of nonreal conjugate pairs at time
$a=(72\pi)^{2/3}v/(8N^{8/3})$ converges in distribution to
$\operatorname{Poisson}(v^{3/2})$. The pair-count processes
converge in finite-dimensional distributions to the corresponding
inhomogeneous Poisson counting process. If $V_k$ is the limiting
$k$th collision time, then
\[
 V_k^{3/2}\sim\operatorname{Gamma}(k,1),\qquad
 f_{V_k}(v)=\frac{3}{2(k-1)!}
                       v^{3k/2-1}e^{-v^{3/2}}\quad(v>0).
\]
The limiting collision positions are independent uniform circular
marks, independent of the times. This is a CUE theorem; it
asserts no transfer to arithmetic zero statistics.
\end{theorem}

\begin{proof}
Fix $A>0$ and put
\[
 s_N=N^{-4/3},\quad T_N=A s_N^2,\quad
 \ell_N=N^{-7/6},\quad B_N=N^{1/24},\quad h_N=B_Ns_N.
\]
The three-point divided-difference estimate in NN-A024 gives
\[
 \mathbb P(\text{some circular arc of length }\ell_N
                        \text{ contains three zeros})
       \le C N^9\ell_N^8=O(N^{-1/3}).
\]
Work on the complementary event. Call each adjacent gap
$\delta\le h_N$ a small-pair cluster and each remaining root
a singleton. Since $2h_N<\ell_N$ eventually, no root belongs
to two small pairs. Every other root lies at distance at least
$\ell_N-h_N\ge\ell_N/2$ from a small pair. Two singleton
roots are at least $h_N$ apart. These assertions apply also
across the circular cut, by using periodic real lifts.

We need uniform control of the other periodic zeros. The
no-triples condition gives at most two zeros in each interval
of length $\ell_N$ on the periodic real line. Shell summation
therefore gives, for a pair center with its two selected roots
removed,
\begin{equation}\label{eq:analytic-cue-uniform-pair-SL}
 S_{\rm pair}\le C\ell_N^{-2},\qquad
 |L_{\rm pair}|\le
                  C\ell_N^{-1}\log(2/\ell_N).
\end{equation}
For the inverse-square sum, bound the shells at distances
$k\ell_N/2$ by $C(k\ell_N)^{-2}$ and sum $k^{-2}$,
including all periodic copies. For the logarithmic drift,
the original trigonometric product gives exactly
\[
 L_{\rm pair}=\frac12\sum_{j\ {
m outside\ pair}}
                         \cot((b-\theta_j)/2).
\]
The selected two cotangents cancel at their midpoint, as do
the two linear-factor logarithmic derivatives. On circular
distances $d\in(0,\pi]$, $|\cot(d/2)|\le C/d$; summing the
finitely many shells gives the logarithm in
\eqref{eq:analytic-cue-uniform-pair-SL}. This avoids assigning
a nonconvergent unsigned first trace to all periodic lifts.

For a singleton at $b$, its other zeros are at least $h_N$
away. At most a fixed number can be closer than $\ell_N$;
the remaining shells give
\begin{equation}\label{eq:analytic-cue-uniform-single-SL}
 S_{\rm single}\le C(h_N^{-2}+\ell_N^{-2}),\qquad
 |L_{\rm single}|\le
               C(h_N^{-1}+\ell_N^{-1}\log(2/\ell_N)).
\end{equation}
Here $L_{\rm single}=\frac12\sum_{j\ne i}
\cot((b-\theta_j)/2)$; the regular term from its own sine
factor vanishes at $b$.

Uniformly over all pair clusters,
\[
 (\delta^2+T_N)S_{\rm pair}
       \le C(h_N^2+T_N)\ell_N^{-2}=O(N^{-1/4})=o(1).
\]
Uniformly over all singletons,
\[
 T_NS_{\rm single}
       \le C_A(B_N^{-2}+N^{-1/3})=o(1).
\]
Lemma~\ref{lem:analytic-uniform-clusters} therefore applies
simultaneously to every cluster on $[0,T_N]$. Their moving
discs remain pairwise disjoint. Indeed pair discs have radii
$2\sqrt{\delta^2+T_N}=O_A(h_N)=o(\ell_N)$ and their center
shifts are at most $C_AT_N\ell_N^{-1}\log(2/\ell_N)
=o(\ell_N)$. Singleton discs have radius $2\sqrt{T_N}=O_A(s_N)$;
their shifts, by \eqref{eq:analytic-cue-uniform-single-SL}, are
\[
 O_A(T_N/h_N+T_N\ell_N^{-1}\log(2/\ell_N))=o(s_N).
\]
They are separated from one another by at least $h_N=B_Ns_N$,
and from any pair by at least $\ell_N/2$ before these vanishing
relative displacements. Each disc has diameter tending to zero,
so disjointness also holds modulo $2\pi$.

The finite polynomial in $w=e^{iz}$ from NN-A024 has degree
$N$, with nonzero constant and leading coefficients at every
finite $a$. The disjoint discs contain a total of exactly $N$
zeros modulo $2\pi$, counted with multiplicity. They therefore
exhaust the entire zero divisor throughout $[0,T_N]$. This is
the step that remains valid after some pairs become nonreal;
no sign inequality for the later global root flow is required.
The singleton roots stay real and simple. Each pair has exactly
its unique collision, if that time is in the window, and cannot
return during the window. In particular there are no other
early collisions or undetected nonreal roots.

For every small pair the uniform collision expansion gives
\[
 a_c=\frac{\delta^2}{8}(1+O(N^{-1/4})).
\]
Only gaps with $\delta=O_A(s_N)$ can thus collide before $T_N$.
For those gaps the sharper bound $\delta^2S=O_A(N^{-1/3})$
is available. The collision mark differs from the left endpoint
of its initial gap by at most
\[
 \delta/2+2T_N|L_{\rm pair}|+O(\delta(\delta^2S)^{3/2})=o(1).
\]
Consequently on each bounded rescaled time window the collision
process is, up to a uniformly vanishing displacement, the image
of the initial extreme-gap process under
\[
                  (y,b)\longmapsto(y^2,b),\qquad
                  y=N^{4/3}\delta/(72\pi)^{1/3}.
\]
Ben Arous--Bourgade, Theorem~1.1, printed p.~2650, states
that the process in the unnormalized coordinate
$u=N^{4/3}\delta$ converges to the Poisson process of intensity
$u^2du/(24\pi)\,db/(2\pi)$. Thus the $y$ coordinate has
intensity $3y^2dy\,db/(2\pi)$, and squaring gives
\eqref{eq:analytic-collision-ppp-intensity}.
Its mass on $[0,v]\times\mathbb R/(2\pi\mathbb Z)$ is
$v^{3/2}$. The limit has no atoms at deterministic time or
spatial boundaries. Tightness of the number of extreme gaps
in bounded windows, supplied by that same point-process
convergence, justifies the vanishing-displacement comparison
and all asserted count limits. The formula for $V_k$ follows
by the deterministic time change $v\mapsto v^{3/2}$ to a
rate-one Poisson process. The product intensity yields the
independent uniform marks.
\end{proof}

\paragraph{Primary source and scope.}
The imported random-matrix result is exactly Ben Arous--Bourgade,
\emph{Extreme gaps between eigenvalues of random matrices},
Annals of Probability~41 (2013), 2648--2681,
Theorem~1.1 p.~2650 and Corollary~1.2 p.~2651,
\url{https://math.nyu.edu/~bourgade/papers/ExtremeGaps.pdf},
DOI~10.1214/11-AOP710. The source's full marked intensity,
including $db/(2\pi)$, was checked. The new argument supplies
the heat evolution, exhaustion after the first collision,
and absence of recollisions on bounded extreme-time windows.
A, M and ordinary fixed-scale pair correlation do not supply
an arithmetic version of the imported extreme-gap process.
There is no expectation convergence or error rate for the
limiting Poisson law claimed here.

\paragraph{Geometric and representation meaning.}
The actual finite free family
\[
 \operatorname{Spec}\mathbb C[q,q^{-1},w,w^{-1}]/
       \left(\sum_{k=0}^Nc_kq^{(2k-N)^2}w^k\right)
                           \longrightarrow\mathbb G_{m,q},
                  \qquad q=e^{-a/4},
\]
is the family already identified in NN-A024. On each bounded
extreme-time interval, with probability tending to one, its
nearby analytic zero algebra splits into the tracked singleton
and rank-two cluster factors. Each rank-two factor has at
most one simple trace-discriminant zero there, with a quadratic
local normalization; each singleton remains etale. SGA~1,
Expos\'e~I, section~4, Proposition~10 is the exact finite-locally-free
trace-discriminant criterion, and section~3, Proposition~1 gives
the vanishing-differentials criterion. The resulting simple
branch points obey the proved marked Poisson law. Analytic
monodromy can be taken on a common small tubular parameter
neighborhood of the real interval, preserving the disjoint
cluster factors. Around the separate early branch points it
exchanges disjoint pairs of sheets, so these local transpositions
commute.
This is a concrete random complex algebraic discriminant family,
not an arithmetic Frobenius equidistribution theorem, a pure
motive, or a positive integral basis.


\section{Why an angular bound alone has no converse in finite order}
\label{sec:analytic-angle-only-obstruction}

\begin{theorem}[NN-A026: fixed nonreal angles with arbitrarily high finite positivity]
For every integer $n\ge1$ put
\[
 \vartheta_n=\frac{\pi}{3n},\qquad
 g_n(z)=\prod_{j=0}^{n}
                (1+e^{i\vartheta_n(n-2j)}z),\qquad
 Q(z)=\cosh\sqrt z,\qquad F_n(z)=g_n(z)Q(z).
\]
Then $F_n$ is real entire of order $1/2$, $F_n(0)=1$, and all
its Taylor coefficients are strictly positive. Every allowed
Toeplitz minor of order at most $2n$ is strictly positive.
Nevertheless $F_n$ has the two fixed nonreal zeros
$-e^{i\pi/3}$ and $-e^{-i\pi/3}$, and its maximal zero-angle
deviation from the negative real ray is exactly $\pi/3$.
Its least negative-minor order is finite and exceeds $2n$.

Rescaling to $\widehat F_n(z)=F_n(z/F_n'(0))$ also imposes
$\widehat F_n'(0)=1$ and preserves the same angular and finite
positivity conclusions. Separately, put
$\Psi_n(z)=F_n(z)/F_n(1)$. Then
$U\mapsto\det\Psi_n(U)$ is a normalized central positive-type
function on every $U(d)$ with $d\le2n$, with strictly positive
coefficients for every polynomial irreducible, yet fails to be
positive type at some larger finite rank.

Thus the maximal zero angle, the entire growth order $1/2$,
and the displayed coefficient normalizations alone give no
upper bound on necessary detecting order or unitary rank.
The positive lower bounds in NN-A022 use an angular sufficient
condition; a matching upper bound for the actual heat family
requires additional information about that family.
\end{theorem}

\begin{proof}
Let $m=n+1$, $k=2n$, and more generally consider the centered
geometric alphabet
\[
 x_j=e^{i\vartheta(m+1-2j)},\quad 1\le j\le m,
                \qquad 0<\vartheta\le\frac{\pi}{m+k-1}.
\]
For a partition $\lambda$ with at most $m$ rows, the Weyl
alternant evaluated at this geometric progression gives
\begin{equation}\label{eq:analytic-geometric-schur-sines}
 s_\lambda(x_1,\ldots,x_m)
     =\prod_{1\le i<j\le m}
        \frac{\sin(\vartheta(\lambda_i-\lambda_j+j-i))}
             {\sin(\vartheta(j-i))}.
\end{equation}
For completeness, write the alternant columns as geometric
powers with exponents $\lambda_j+m-j$. Both determinants
become Vandermonde products. Replacing
$e^{iu}-e^{-iu}$ by $2i\sin u$ gives the quotient above;
the overall phases cancel because the alphabet is centered.
The denominator is nonzero, since $0<\vartheta(j-i)<\pi$.
The formula also follows directly from the Weyl character
formula, Etingof et al., Theorem~5.22.1, printed p.~132.

If $\lambda_1\le k$, every numerator argument lies in $[0,\pi]$:
\[
 1\le\lambda_i-\lambda_j+j-i\le k+m-1.
\]
Hence all ordinary Schur values of width at most $k$ are
nonnegative. Values with more than $m$ rows are zero by the
alternating definition. Every skew Schur function of width
at most $k$ is a nonnegative integer linear combination of
ordinary Schur functions of that width: the restriction
multiplicities used and proved in NN-A021 are nonnegative,
and a nonzero coefficient entails containment of the relevant
partition in the outer partition. The exact minor/width
correspondence in NN-A020 now proves that the coefficients
of $\prod_j(1+x_jz)$ are $PF_k$.

Taking $m=n+1,k=2n$ and
$\vartheta=\pi/(m+k-1)=\pi/(3n)$ gives the stated $g_n$.
Its roots are $-x_j^{-1}$; the endpoint phases are exactly
$\pm\pi/3$, with all other phases between them. Its
coefficients are nonnegative by the already proved order-one
minors. Conjugate pairing of the alphabet makes $g_n$ real.

The entire function $Q(z)=\cosh\sqrt z$ has Taylor coefficients
$1/(2r)!$, and the genus-zero product
\[
 Q(z)=\prod_{j\ge1}\left(1+
                       \frac{4z}{\pi^2(2j-1)^2}\right).
\]
Thus $Q$ is $PF_\infty$ and has strictly positive allowed
Toeplitz minors, by the explicit positive-path proof in
NN-A008. Its coefficients are strictly positive. Multiplying
by $g_n$ preserves $PF_{2n}$ by Cauchy--Binet. It actually
preserves strict positivity of every allowed minor: in
$T_QT_{g_n}$, select the intermediate index set equal to the
column index set of the target minor. The selected
$T_{g_n}$ minor is upper triangular with diagonal one,
while the selected $T_Q$ minor is the strictly positive
allowed minor. Every other Cauchy--Binet term is nonnegative.
The sum is finite for each fixed index set, by upper
triangularity. The coefficients of the product are positive
because the constant coefficient of $g_n$ is one.

A polynomial factor changes neither the order $1/2$ of $Q$
nor its negative-real infinite zero set. The fixed nonreal
zeros of $g_n$ remain, so the maximal angle is exactly
$\pi/3$. Edrei's entire-function converse, with the exact
original locator in NN-A007, rules out $PF_\infty$;
therefore some finite minor is negative. None has order
at most $2n$. Positive scaling of the variable multiplies
any Toeplitz minor by a positive number and preserves angles,
which proves the derivative-normalized assertion.

Finally $F_n(1)>0$, so $\Psi_n(0)>0$ and $\Psi_n(1)=1$.
The exact $PF_d$ versus positive-$U(d)$ criterion in NN-A021
proves the asserted characters, including strict positivity
of all polynomial-irreducible coefficients. The same
criterion and Edrei give failure at some finite larger rank.
This character normalization and the derivative normalization
are separate operations; no preservation under a negative
translation of the variable is assumed.
\end{proof}

\paragraph{Classical locator and the new use.}
The finite sharp family is the Schoenberg example recorded in
Adm--Garloff--Tyaglov, \emph{Total nonnegativity of finite
Hurwitz matrices and root location of polynomials},
Theorem~5.7 and equation~(5.9), printed p.~161, immediately before the sector
sufficiency Theorem~5.8, equation~(5.10). The displayed Weyl
calculation supplies a complete proof of precisely the
finite-order assertion needed here. The order-$1/2$ infinite
extension, strict-minor observation, normalizations and
finite-unitary-rank interpretation are given above. This is
an obstruction to an unrestricted converse, not a counterexample
to a heat-specific upper estimate.

\paragraph{An explicit correction to the source's product inference.}
Immediately after equations~(5.11)--(5.12), printed p.~162,
the same published source asserts that a product of quadratics
whose angles are exterior to the quadratic $PF_k$ interval
cannot belong to $PF_k$. That inference is false, even when
every factor has the same outside angle. Set
\[
 q(z)=1+\frac9{10}z+z^2,\qquad
 \theta=\arccos\frac9{20}>\frac\pi3.
\]
Then $q=(1+e^{i\theta}z)(1+e^{-i\theta}z)$ is not $PF_2$,
since its central order-two minor is $81/100-1<0$.
On the other hand
\[
 q(z)^2=1+\frac95z+\frac{281}{100}z^2+\frac95z^3+z^4
\]
is $PF_2$. Its three nontrivial adjacent log-concavity
differences are $43/100$, $46561/10000$, and $43/100$,
all positive. For a positive finite consecutive coefficient
sequence, these differences make the successive ratios
nonincreasing. Therefore every order-two Toeplitz minor
\[
       a_{h+d}a_{h+e}-a_ha_{h+d+e}\quad(d,e>0)
\]
is nonnegative by multiplying the corresponding ratio
inequalities; when an index is outside the support,
upper triangularity and the consecutive support give the
same conclusion directly. This proves the required $PF_2$
assertion without a numerical test. The exact published
source is
\url{https://www-home.htwg-konstanz.de/~garloff/JMAA-3_Adm_Garloff_Tyaglov.pdf},
DOI~10.1016/j.jmaa.2018.06.065. The valid sector sufficiency
and valid single-quadratic criterion used in NN-A020 are
unaffected; they were also proved there independently.

\paragraph{Geometric and representation meaning.}
The values in \eqref{eq:analytic-geometric-schur-sines} are
literal Weyl characters evaluated at a centered diagonal
unitary matrix. In the stated rectangle of partition widths,
they are nonnegative real numbers despite its nonreal
eigenvalues. Their skew restrictions give the positive
Toeplitz tests, and the independent positive $Q$ specialization
makes all allowed tests strict. The representation-ring
exterior-power and exact-sequence formalism has the exact
locator SGA~6, Expos\'e~V, Example~2.2(b) and Lemma~2.2.1.
The functions on each $U(d)$ are actual mixtures of its
polynomial irreducible characters; they fail to extend
positively through all ranks for the fixed analytic function.
Their real coefficients are not integral representation
multiplicities. The example separates finite character-cone
positivity from a constraint depending only on the support
angle of the analytic zero divisor.


\section{A sharp detecting scale when a double endpoint controls the angle}
\label{sec:analytic-dominant-double}

The following conditional case closes a matching upper-bound gate for
the actual heat family. The extra multiplicity hypothesis is displayed;
it is not part of A. We first prove the infinite-alphabet estimate that
makes the negative minor uniform near a coalescing pair.

\begin{lemma}[A coalescing dominant pair in rectangular Schur minors]
\label{lem:analytic-rectangular-pair}
Let an $\ell^1$ complex alphabet, closed under conjugation, depend on
a positive parameter tending to zero. Suppose its entries consist of
\[
 b_1>\cdots>b_p>b,\qquad
 \beta_+=be^{i\theta},\quad\beta_-=be^{-i\theta},\qquad
                         (\eta_j)_{j\ge1},
\]
where $\theta\downarrow0$. Assume $b_i\to b_i^0$,
$b\to b^0>0$, with strictly separated limits
$b_1^0>\cdots>b_p^0>b^0$ when $p>0$ (vacuous for $p=0$), and
$\sup_j|\eta_j|\le c_0<b^0$ uniformly. Assume also
$\sum_j|\eta_j|$ is uniformly bounded, and that the tail products
\[
             C(z)=\prod_j(1-\eta_jz)^{-1}
\]
converge locally uniformly on $|z|<1/c_0$ to the corresponding product of an endpoint
alphabet of positive real entries less than $b^0$. Write $k=p+1$,
$P=\prod_{i=1}^p b_i$, and $h_r$ for the complete symmetric functions
of the full alphabet. Then there are constants $C_1>0$, $0<q<1$,
independent of $n$ and of the sufficiently small parameter, such that
\begin{equation}\label{eq:analytic-rectangular-asymptotic}
 \frac{s_{(n^k)}}{(Pb)^n}
  =\frac{|D|\sin((n+1)\theta+\arg D)}{\sin\theta}
                         +O(C_1\theta^{-k}q^n),\qquad n\ge k,
\end{equation}
where
\[
 D=\left\{\prod_{i=1}^p
              \frac{C(1/b_i)}{1-\beta_-/b_i}\right\}C(1/\beta_+)
                             \longrightarrow D_0>0.
\]
In particular if $n\theta\to u\in(\pi,2\pi)$, the displayed
rectangular Schur value is strictly negative eventually.
The same value is the elementary Toeplitz minor
\begin{equation}\label{eq:analytic-rectangular-elementary-minor}
                  s_{(n^k)}=det[e_{k-i+j}]_{i,j=1}^n.
\end{equation}
It has order $n$ and fits in the index window $0,\ldots,k+n-1$.
\end{lemma}

\begin{proof}
Choose $c$ strictly between $c_0$ and $b^0$, and smaller than
all the limiting $b_i$. For sufficiently small parameter, the generating
function
\[
 H(z)=\sum_{r\ge0}h_rz^r
       =\prod_{i=1}^p(1-b_i z)^{-1}
           (1-\beta_+z)^{-1}(1-\beta_-z)^{-1} C(z)
\]
has exactly these $p+2$ simple poles inside $|z|=1/c$.
Its modulus on that circle is uniformly bounded. Remove the simple
partial fractions there:
\[
 H(z)=\sum_{\alpha\in\mathcal B}
          \frac{A_\alpha}{1-\alpha z}+G(z),\qquad
 \mathcal B=\{b_1,\ldots,b_p,\beta_+,\beta_-\},
\]
where
\[
 A_\alpha=\prod_{\gamma\ne\alpha}(1-\gamma/\alpha)^{-1}.
\]
The product includes the whole tail. The tail product converges
absolutely and is bounded above and away from zero at the required
points, by its uniform $\ell^1$ bound and strict modulus separation.
The only vanishing denominator is
$1-\beta_-/\beta_+$ or its conjugate. Hence
$|A_\alpha|\le C\theta^{-1}$, and Cauchy's formula gives
\begin{equation}\label{eq:analytic-h-poles-remainder}
 h_r=\sum_{\alpha\in\mathcal B}A_\alpha\alpha^r+g_r,
                         \qquad |g_r|\le C\theta^{-1}c^r.
\end{equation}
The remainder $G$ is analytic through all removed poles. Its circle
bound follows from those of $H$ and the partial fractions, so no
uniform separation of the two coalescing poles was assumed.

Use the ordinary Jacobi--Trudi determinant
$s_{(n^k)}=\det[h_{n-i+j}]_{i,j=0}^{k-1}$.
The pole part of its matrix is a sum of rank-one matrices
\[
        A_\alpha\alpha^n
              (\alpha^{-i})_{i=0}^{k-1}
              (\alpha^{j})_{j=0}^{k-1}.
\]
Cauchy--Binet for that finite sum gives one term for each
$k$-element subset $J\subset\mathcal B$. The two Vandermonde
determinants cancel all internal factors in
$\prod_{\alpha\in J}A_\alpha$, leaving exactly
\begin{equation}\label{eq:analytic-rectangular-subset-coefficient}
 \left(\prod_{\alpha\in J}\alpha\right)^n
       \prod_{\substack{\alpha\in J\\\gamma\notin J}}
                                      (1-\gamma/\alpha)^{-1}.
\end{equation}
Here the complement includes the infinite tail, only through its
absolutely convergent products. To see the cancellation directly,
each pair $\alpha,\alpha'$ within $J$ contributes
$-(\alpha\alpha')/(\alpha-\alpha')^2$ from the two partial
fractions, the inverse of its contribution from the two Vandermondes.

The two subsets containing every $b_i$ and exactly one of
$\beta_\pm$ have modulus weight $(Pb)^n$. Their coefficients are
$D/(1-e^{-2i\theta})$ and its conjugate. Since
\[
             (1-e^{-2i\theta})^{-1}
                         =e^{i\theta}/(2i\sin\theta),
\]
their sum is precisely the first term of
\eqref{eq:analytic-rectangular-asymptotic}. The endpoint tail is
positive real, all $b_i^0>b^0$, and the displayed products therefore
give $D\to D_0>0$.

Every other $k$-element subset omits at least one $b_i$ and
contains both members of the pair. Its modulus weight divided by
$Pb$ is $b/b_i$, bounded by a constant less than one. Its coefficient
is uniformly bounded: the coalescing pair is internal to $J$, so its
potential singular factors already cancelled in
\eqref{eq:analytic-rectangular-subset-coefficient}.

It remains to control determinants that use remainder columns from
\eqref{eq:analytic-h-poles-remainder}. This cannot be done by a
naive bound using the largest pole to the $k$th power. Instead use
multilinearity. A term with $m\ge1$ remainder columns has $k-m$
remaining rank-one pole columns, whose pole labels must be distinct
or the determinant is zero. Its modulus is at most a fixed constant
 times
\[
 \theta^{-k}\left(c^m\prod_{\alpha\in J}|\alpha|\right)^n,
                              \qquad |J|=k-m.
\]
Fixed shifts $j-i$ affect only the constant. The largest product of
$k-m\le p$ pole moduli uses that many of the $b_i$. Because each
remaining $b_i$ is larger than $b$,
\[
          \frac{c^m\prod_{\alpha\in J}|\alpha|}{Pb}
                                      \le(c/b)^m\le q_1<1.
\]
There are finitely many such multilinear terms, their number depending
only on $k$. Combining this bound with the nonleading subsets proves
\eqref{eq:analytic-rectangular-asymptotic}, after increasing $q<1$.
If $n\theta\to u\in(\pi,2\pi)$, its first term is asymptotic to
$D_0\sin u/\theta<0$, and its exponential remainder is smaller than
any power of $\theta$. Finally dual Jacobi--Trudi transposes
$(n^k)$ to $(k^n)$ and gives
\eqref{eq:analytic-rectangular-elementary-minor}. Explicitly use
row indices $0,\ldots,n-1$ and column indices $k,\ldots,k+n-1$
in the upper-triangular Toeplitz matrix $(e_{j-i})$.
\end{proof}

\begin{theorem}[NN-A027: exact detecting constants in the dominant-double case]
\label{thm:analytic-dominant-double}
Assume A. For a positive zero $x$ of $H_0$, let $m_x$ be its
multiplicity and $h_{m_x}$ the largest absolute root of the
probabilists' Hermite polynomial $\operatorname{He}_{m_x}$,
with $h_1=0$. Make the following additional hypothesis:
there is a positive \emph{double} zero $x_0$ such that
\begin{equation}\label{eq:analytic-dominant-double-hypothesis}
                    \frac{h_{m_x}}x\le\frac1{x_0}
                    \qquad\hbox{for every positive zero }x.
\end{equation}
This is a precise conditional case; neither this multiplicity pattern
nor existence of any multiple zero is inferred from A.

Let $m_*(-a)$ be the least order of a negative Toeplitz minor of
$F_{-a}$, and $R_*(-a)$ its least detecting index window as in
NN-A014 and NN-A020. Then
\begin{equation}\label{eq:analytic-sharp-order-window}
       \lim_{a\downarrow0}\sqrt a\,m_*(-a)
        =\lim_{a\downarrow0}\sqrt a\,R_*(-a)
                       =\frac{\pi x_0}{2\sqrt2}.
\end{equation}
For the normalized unitary characters of NN-A021, put
\[
 p_0=F_0'(0)=\sum_{x>0}x^{-2},\qquad
 \beta_0=\frac1{p_0x_0^2}\in(0,1),
\]
and let $d_*(-a)$ be the least rank at which
$U\mapsto\det E_{-a}(U-I)$ fails to be positive type. Then
\begin{equation}\label{eq:analytic-sharp-unitary-rank}
       \lim_{a\downarrow0}\sqrt a\,d_*(-a)
                       =\frac{\pi x_0(1-\beta_0)}{2\sqrt2}.
\end{equation}
There are explicit negative rectangular witnesses: if $p$ is the
number of positive endpoint zeros below $x_0$ and $k=p+1$, take
$n\theta(a)\to u\in(\pi,2\pi)$, where $\theta(a)$ is the positive
argument of the corresponding complex product parameter. The minor
$\det[e_{k-i+j}]_{i,j=1}^n$ is negative. For the translated unitary
alphabet this is the negative Fourier coefficient of the actual
one-dimensional representation $\det(U)^k$ on $U(n)$.
\end{theorem}

\begin{proof}
For $m\ge2$, $h_m\ge h_2=1$, by the principal Jacobi-matrix
interlacing of the Hermite recurrence. Therefore
\eqref{eq:analytic-dominant-double-hypothesis} forces every endpoint
zero below $x_0$ to be simple. There are finitely many of them,
$x_1<\cdots<x_p<x_0$. Their continuations at time $-a$ remain
simple and real for small $a$. The double zero at $x_0$ splits by
NN-A006 into
\[
             z_\pm(a)=x_0\pm i\sqrt{2a}+O(a).
\]
The two parameters $z_\pm(a)^{-2}$ are conjugate and have arguments
$\pm\theta(a)$ with
\begin{equation}\label{eq:analytic-pair-theta-sharp}
                         \frac{\theta(a)}{\sqrt a}
                                      \longrightarrow\frac{2\sqrt2}{x_0}.
\end{equation}
The $p$ real parameters from lower zeros have strictly larger moduli.
All remaining parameters have uniformly smaller modulus than
$x_0^{-2}$: choose a circle between $x_0$ and the next larger positive
endpoint zero, and use compact zero convergence. Zeros outside a
large fixed circle contribute parameters tending uniformly to zero.

Here and below the infinite products and their uniform bounds are
legitimate without simplicity of the other zeros. For $t$ in a compact
interval the theta integral gives
$\log\max_{|z|\le R}|H_t(z)|\le C R\log(R+3)$, with $C$ uniform,
and $H_t(0)$ is bounded away from zero near $t=0$. Jensen's formula
therefore gives a uniform zero count $N_t(R)\le C R\log(R+3)$.
Partial summation bounds uniformly the tails of
$\sum|z_j(t)|^{-2}$ and makes them tend to zero as the radius tends
to infinity. The genus-zero products of $F_t$, after removing the
fixed finite cluster, consequently converge locally uniformly to the
endpoint tail product. In particular the hypotheses of
Lemma~\ref{lem:analytic-rectangular-pair} hold for the alphabet
$z_j(-a)^{-2}$, with $k=p+1$.

Take $n=\lceil(\pi+\varepsilon)/\theta(a)\rceil$ for any fixed
$0<\varepsilon<\pi$. The lemma gives a negative minor of order $n$
in window $k+n-1$ for all sufficiently small $a$. Hence
\[
 m_*(-a)\le n,\qquad R_*(-a)\le k+n-1,
\]
and \eqref{eq:analytic-pair-theta-sharp} gives the corresponding
upper limits with $\pi+\varepsilon$ in place of $\pi$.

For the lower limits, NN-A022 and the extra hypothesis give
\[
 \frac{\alpha_H(-a)}{\sqrt a}\longrightarrow\frac{\sqrt2}{x_0}.
\]
The maximal parameter angle is exactly $2\alpha_H(-a)$ when these
angles are small. The strict quadratic sector proof in NN-A020
shows that every allowed minor of order at most $r$ is positive
whenever $(r+1)2\alpha_H(-a)<\pi$. Thus for any fixed
$\varepsilon>0$ one may take
$r=\lfloor(\pi-\varepsilon)/(2\alpha_H(-a))\rfloor-1$.
It follows that
\[
 \liminf_{a\downarrow0}\sqrt a\,m_*(-a)
                           \ge\frac{\pi x_0}{2\sqrt2}.
\]
Any minor in the window $0,\ldots,R$ has order at most $R+1$,
so $R_*(-a)\ge m_*(-a)-1$. Letting $\varepsilon\downarrow0$
in the upper bounds proves \eqref{eq:analytic-sharp-order-window}.

For the unitary statement put $p_a=F_{-a}'(0)>0$. The translated
alphabet in NN-A021 is
\[
 \widetilde\beta_j(a)
   =\frac{\beta_j(a)}{1-\beta_j(a)}
   =\frac1{p_a z_j(-a)^2-1},\qquad
                    \beta_j(a)=\frac1{p_a z_j(-a)^2}.
\]
At the endpoint all its entries are positive and are strictly decreasing
as functions of the positive zero $x$. Also $\sup_j|\beta_j(a)|<b<1$
for small $a$, by the compact zero separation already proved in
NN-A021. Its $\ell^1$ tail and local convergence therefore satisfy
the same lemma. At an endpoint zero $x$ of multiplicity $m_x$,
the largest local angle coefficient becomes
\begin{equation}\label{eq:analytic-translated-angle-coefficient}
 \frac{2\sqrt2\,h_{m_x}}{x(1-1/(p_0x^2))}.
\end{equation}
This follows by substituting the local Hermite expansion into
$(p_a z^2-1)^{-1}$; $p_a=p_0+O(a)$ is real.

The maximum of these coefficients is attained at $x_0$.
Indeed the nonzero coefficients occur only at $x\ge x_0$,
$h_{m_x}/x\le1/x_0$, and
$x\mapsto(1-1/(p_0x^2))^{-1}$ decreases there. The uniform
radial-tail estimate in NN-A022 also applies after this transformation:
for high zeros $|\beta_j(a)|\le C/|z_j(-a)|^2$, while
\[
 \left|\arg\frac{\beta}{1-\beta}\right|
                         \le\frac{|\arg\beta|}{1-|\beta|}
                         \qquad(|\beta|<1).
\]
One proves this last bound by integrating
$\operatorname{Re}(1-re^{i\theta})^{-1}\le(1-r)^{-1}$
along the argument, starting at a positive real number.
Consequently the maximal translated angle, divided by $\sqrt a$,
has the exact limit
$2\sqrt2/[x_0(1-\beta_0)]$.

The sector lower bound and the same rectangular upper bound now give
\eqref{eq:analytic-sharp-unitary-rank}, using the exact $PF_d$
versus $U(d)$ equivalence of NN-A021. Multiplication by the positive
constant $\Psi_a(0)=E_{-a}(-1)$ does not affect minor signs.
For the upper witness, the coefficient formula from that equivalence
is
\[
 c_{(k^n),n}=\det[a_{k-i+j}]_{i,j=1}^n
             =\Psi_a(0)^n s_{(n^k)}(\widetilde\beta(a))<0.
\]
The highest weight $(k,\ldots,k)$ is precisely $\det(U)^k$,
so this is the claimed concrete irreducible witness.
\end{proof}

\paragraph{Geometric and representation meaning.}
The additional hypothesis states that a local Artin algebra of length
two controls the maximal scaled angular motion of the endpoint divisor.
Its quadratic normalization produces the two conjugate dominant
parameters. The rectangular Schur computation cancels the internal
Vandermonde factors and isolates their sine oscillation with a uniformly
controlled infinite complementary divisor. Dual Jacobi--Trudi turns
this into a finite minor whose order and index window have the exact
asymptotic above. After the actual unitary translation, that same
rectangle is a negative coefficient of $\det^k$ at rank $n$.
The exterior-power/representation-ring dictionary is SGA~6, Expos\'e~V,
Example~2.2(b) and Lemma~2.2.1; the rank-two discriminant criterion is
SGA~1, Expos\'e~I, section~4, Proposition~10. These are sharp bounds
for the specified analytic and representation certificate families
under the displayed multiplicity hypothesis. They are not a proof
of that hypothesis, a pure motive, or a lower bound against all
algorithms.


\section{Endpoint simplicity and the sharp order of finite detection}
\label{sec:analytic-multiplicity-detection-equivalence}

The dominant-double hypothesis in NN-A027 is needed for its displayed
\emph{exact constants}, but not for a matching order of growth when
some endpoint is multiple. We prove this using an actual compact
unitary orbit and an explicit Schur isolation argument.

\begin{lemma}[A negative Hermitian orbital Fourier value]
\label{lem:analytic-negative-orbital}
Let $D\ne0$ be a Hermitian $m$ by $m$ matrix with trace zero and
spectrum invariant under multiplication by $-1$. Then
\[
 \phi_D(X)=\int_{U(m)}e^{i\operatorname{Tr}(DUXU^*)}\,dU,
                           \qquad X=X^*,
\]
is real and takes a strictly negative value. One may choose a
negative value at a diagonal matrix with distinct positive eigenvalues.
\end{lemma}

\begin{proof}
The orbit of $D$ equals that of $-D$. Haar invariance therefore makes
$\phi_D$ real. It is continuous, conjugation-invariant, and
$\phi_D(0)=1$. Give the real vector space $\mathcal H_m$ of Hermitian
matrices the Euclidean inner product $\operatorname{Tr}(XY)$ and
Lebesgue measure for an orthonormal basis; its dimension is $d=m^2$.
Fubini and the elementary Gaussian Fourier integral give, for every
$t>0$,
\begin{equation}\label{eq:analytic-orbital-Gaussian}
 \int_{\mathcal H_m}e^{-t\|X\|^2/2}\phi_D(X)\,dX
       =\left(\frac{2\pi}{t}\right)^{d/2}
                         e^{-\|D\|^2/(2t)}.
\end{equation}
The right side tends to zero as $t\downarrow0$, because $D\ne0$.
If $\phi_D$ were nonnegative everywhere, a fixed ball about zero on
which $\phi_D\ge1/2$ would give a positive lower limit for the left
side, a contradiction. Thus some $X$ has $\phi_D(X)<0$.
Diagonalize it, perturb its eigenvalues to be distinct while preserving
strict negativity, and add a scalar matrix large enough to make every
eigenvalue positive. This last operation does not change $\phi_D$
because $\operatorname{Tr}D=0$.
\end{proof}

\paragraph{Exact compact integral locator.}
We use the classical Harish--Chandra--Itzykson--Zuber formula with Haar
probability measure. If $A=\operatorname{diag}(a_1,\ldots,a_m)$ and
$B=\operatorname{diag}(b_1,\ldots,b_m)$ have distinct eigenvalues, then
\begin{equation}\label{eq:analytic-HCIZ}
 \mathcal I(A,B):=\int_{U(m)}e^{\operatorname{Tr}(AUBU^*)}\,dU
       =\left(\prod_{j=1}^{m-1}j!\right)
                    \frac{\det[e^{a_i b_j}]}{\Delta(a)\Delta(b)}.
\end{equation}
It extends to complex diagonal $A$ and coincident eigenvalues by
holomorphic continuation; the integral is entire in the matrix entries
on compact sets. A precise primary locator is J.-B.~Zuber,
\emph{Horn's problem and Harish-Chandra's integrals. Probability
density functions}, Ann. Inst. Henri Poincar\'e Comb. Phys. Interact.~5
(2018), 309--338, orbital definition (6), p.~311, formula (11), p.~312,
and constant (12), p.~313, DOI~10.4171/AIHPD/56,
\url{https://www.numdam.org/item/10.4171/aihpd/56.pdf}.
Only this exact classical identity is imported; the sign and isolation
arguments needed below are proved here.

\begin{lemma}[A whole coalescing cluster gives a bounded scaled-width witness]
\label{lem:analytic-cluster-Schur-witness}
Let an $\ell^1$ alphabet, closed under conjugation, have the following
form as $\varepsilon\downarrow0$:
\[
 b_1(\varepsilon)>\cdots>b_p(\varepsilon)>0,\quad
 \xi_j(\varepsilon)=b(\varepsilon)
                  e^{\varepsilon a_j(\varepsilon)}\quad(1\le j\le m),
                       \quad (\eta_j(\varepsilon))_{j\ge1}.
\]
Suppose $b>0$, $b\to b^0>0$,
$b_1^0>\cdots>b_p^0>b^0$, and
\[
 a_j(\varepsilon)=i d_j+O(\varepsilon),\qquad
 \sum_j a_j(\varepsilon)=0,
\]
where the $d_j$ are distinct real numbers, not all zero, with a
sign-symmetric multiset. Suppose the tail has uniformly bounded
$\ell^1$ norm, $\sup_j|\eta_j|\le c_0<b^0$, and its product
$C_\varepsilon(z)=\prod_j(1-\eta_jz)^{-1}$ converges locally uniformly
on $|z|<1/c_0$ to a product of positive real endpoint parameters.
The assertions about $p$ are vacuous when $p=0$.

Then there is a constant $C>0$ and a family of partitions
$\lambda(\varepsilon)$ of fixed length at most $q=p+m$, with
$\lambda_1\le C/\varepsilon+1$, such that
\[
                        s_{\lambda(\varepsilon)}<0
\]
for all sufficiently small $\varepsilon$. This is a negative
Toeplitz minor of order at most $C/\varepsilon+1$ in an index
window of size at most $C/\varepsilon+q$.
\end{lemma}

\begin{proof}
Put $D=\operatorname{diag}(d_1,\ldots,d_m)$. The preceding lemma
provides distinct numbers
$L_1>\cdots>L_m>0$ with
$\phi_D(\operatorname{diag}(L))<0$. Choose $C>L_1$ and $\eta>0$
so small that $L_m>2\eta$ and $C-L_1>2\eta$. Put
\[
 \mu_i=\lfloor L_i/\varepsilon\rfloor,\qquad
 n=\lceil C/\varepsilon\rceil,\qquad
 \lambda=(\underbrace{n,\ldots,n}_{p\text{ times}},
                                  \mu_1,\ldots,\mu_m).
\]
These are partitions for small $\varepsilon$, with
$\mu_m\ge\eta/\varepsilon$ and
$n-\mu_1\ge\eta/\varepsilon$. All row lengths are $O(1/\varepsilon)$.

First consider just the cluster, divided by its positive common scale
$b$. Write $A_\varepsilon=\operatorname{diag}(a_j(\varepsilon))$,
$\rho=(m-1,m-2,\ldots,0)$,
$U_\varepsilon=\operatorname{diag}(\varepsilon(\mu+\rho))$ and
$V_\varepsilon=\operatorname{diag}(\varepsilon\rho)$.
The Weyl alternant and \eqref{eq:analytic-HCIZ} give the \emph{exact}
identity
\begin{equation}\label{eq:analytic-Schur-orbital-limit}
 \frac{s_\mu(e^{\varepsilon a_1},\ldots,e^{\varepsilon a_m})}
      {s_\mu(1,\ldots,1)}
   =\frac{\mathcal I(A_\varepsilon,U_\varepsilon)}
          {\mathcal I(A_\varepsilon,V_\varepsilon)}
      \longrightarrow\phi_D(\operatorname{diag}(L))<0.
\end{equation}
Indeed the Vandermonde quotient
$\Delta(\mu+\rho)/\Delta(\rho)$ is the positive Weyl dimension,
and the remaining determinants give the displayed ratio. Its denominator
tends to one. Thus the cluster Schur value is negative, of magnitude
at least $c\,b^{|\mu|}s_\mu(1^m)$ for a fixed $c>0$.
This conclusion is stable under the full $O(\varepsilon)$ perturbation
of $A_\varepsilon$; no exact geometric spacing of the cluster is used.

We now show that the finite positive prefix and infinite tail change
this sign only by a positive factor and an exponentially small error.
All constants below may depend on the fixed $p,m,L,C$ and the stated
spectral separations. Choose $c_1$ with $c_0<c_1<b^0$.
The complete-symmetric generating function has exactly $q=p+m$
simple top poles $\alpha\in\mathcal A=\{b_i,\xi_j\}$ inside
$|z|=1/c_1$. The cluster pole differences are bounded below by a
constant times $\varepsilon$, since the $d_j$ are distinct. The
partial-fraction argument of NN-A027 therefore gives
\begin{equation}\label{eq:analytic-full-cluster-poles}
 h_r=\sum_{\alpha\in\mathcal A}A_\alpha\alpha^r+g_r,
 \qquad |A_\alpha|\le C_2\varepsilon^{-(m-1)},\qquad
 |g_r|\le C_2\varepsilon^{-(m-1)}c_1^r.
\end{equation}
The analytic remainder bound follows on the same fixed circle after
removing all top poles. Uniform tail product bounds ensure its constant
is independent of $r$ and $\varepsilon$.

Apply \eqref{eq:analytic-full-cluster-poles} to the ordinary
Jacobi--Trudi determinant $\det[h_{\lambda_i-i+j}]_{i,j=1}^q$.
The part using all $q$ top poles is exactly
\begin{equation}\label{eq:analytic-full-cluster-leading}
                  s_\lambda(\mathcal A)
                              \prod_{\alpha\in\mathcal A}C_\varepsilon(1/\alpha).
\end{equation}
This is the same internal Vandermonde cancellation as in NN-A027,
now with a general partition instead of a rectangle. The product on
the right is real and tends to a positive number.

Here is a uniform error estimate for all terms with remainder columns.
Let $P=\prod_{i=1}^p b_i$ and $W=P^n b^{|\mu|}>0$.
A term with $r\ge1$ remainder columns can use only $q-r$ distinct
top pole labels; repeated labels give zero. After determinant expansion,
its exponential weight assigns the corresponding moduli, and $r$ copies
of $c_1$, to the nonincreasing row exponents $\lambda_i$.
The rearrangement inequality maximizes that weight by matching their
orders. Compared with the full set of top moduli, replacing any one
by $c_1$ loses a factor at most
$\widetilde q^{\mu_m}$ for a fixed $\widetilde q<1$:
all exponents are at least $\mu_m$, and every top modulus is uniformly
larger than $c_1$. More explicitly the sorted logarithms decrease
coordinatewise, their total decrease is at least
$\log(\min|\alpha|/c_1)$, and multiplication by exponents at least
$\mu_m$ gives this loss. The full top weight is bounded by a constant
times $W$, because $|\xi_j|/b=e^{O(\varepsilon^2)}$ and
$|\lambda|=O(1/\varepsilon)$. Fixed shifts $j-i$ in the matrix entries
affect only a constant. There are finitely many terms, and the residue
bounds in \eqref{eq:analytic-full-cluster-poles} introduce only a fixed
power of $\varepsilon^{-1}$. Thus the total remainder is
\begin{equation}\label{eq:analytic-tail-exponential-error}
                  O\bigl(W\varepsilon^{-K}
                              \widetilde q^{\mu_m}\bigr)
\end{equation}
for some fixed $K$. This explicitly avoids a determinant estimate that
would repeat the largest pole in all rows.

To isolate the prefix inside $s_\lambda(\mathcal A)$, expand its Weyl
numerator along the first $p$ rows. The leading choice assigns the
$p$ positive prefix columns to those rows. Dividing by the full
Vandermonde gives exactly
\begin{equation}\label{eq:analytic-prefix-factorization}
 P^n s_\mu(\xi_1,\ldots,\xi_m)
            \prod_{i=1}^p\prod_{j=1}^m(1-\xi_j/b_i)^{-1}.
\end{equation}
For example the prefix numerator block is
$P^{n+m}\Delta(b_1,\ldots,b_p)$, and the cluster block is
$s_\mu(\xi)\Delta(\xi)$; these cancel the corresponding full
Vandermonde factors. The remaining product is real and tends to a
positive number.

Every other prefix assignment puts at least one larger $b_i$ into a
cluster row, whose length is at most $\mu_1$, and a cluster modulus
into a row of length $n$. It therefore loses a factor at most
$q_2^{n-\mu_1}$, with $q_2<1$ fixed. Pair the misplaced prefix and
cluster columns; the ratio $|\xi|/b_i$ is uniformly less than one,
and apply the rearrangement inequality to any other columns.
The cluster Vandermonde in the denominator has magnitude bounded
below by a constant times $\varepsilon^{m(m-1)/2}$, while all other
Vandermonde factors stay away from zero. Consequently the total wrong
assignment contribution is
\begin{equation}\label{eq:analytic-prefix-exponential-error}
                   O\bigl(W\varepsilon^{-K'}q_2^{n-\mu_1}\bigr).
\end{equation}
For $p=0$ this error is absent. Combining
\eqref{eq:analytic-full-cluster-leading}--\eqref{eq:analytic-prefix-exponential-error},
the full Schur value is a positive factor, bounded away from zero,
times $P^n s_\mu(\xi)$, plus $W$ times an exponentially small quantity
in $1/\varepsilon$. Equation~\eqref{eq:analytic-Schur-orbital-limit}
therefore makes it negative eventually.

Finally put $r=\lambda_1\le n$; equality need not hold when $p=0$.
Dual Jacobi--Trudi gives a minor of order $r$:
\[
                     s_\lambda=\det[e_{\lambda'_i-i+j}]_{i,j=1}^r.
\]
Transpose this determinant. Use row indices $r-i$ and column indices
$\lambda'_j+r-j$, both in decreasing order, and then reverse both.
The resulting increasing index sets lie in $0,\ldots,r+q-1$,
which is contained in $0,\ldots,n+q-1$.
This proves the stated order and window bounds.
\end{proof}

\begin{theorem}[NN-A028: a finite-detection characterization of endpoint simplicity]
\label{thm:analytic-simplicity-detection-equivalence}
Assume A. If $H_0$ has any multiple zero, then there are positive
constants $c,C,a_0$ such that for $0<a<a_0$,
\begin{equation}\label{eq:analytic-any-multiple-detection-scale}
 c a^{-1/2}\le m_*(-a),\ R_*(-a),\ d_*(-a)\le C a^{-1/2}.
\end{equation}
Here the three quantities are respectively the least negative-minor
order, least detecting index window for $F_{-a}$, and least failing
rank of the actual translated unitary characters from NN-A021.
The chained notation means that \emph{each} of the three quantities
lies between the displayed bounds.

Consequently the following four statements are equivalent:
\begin{enumerate}
\item Every nontrivial zero of $\zeta$ is simple.
\item $\sqrt a\,m_*(-a)\longrightarrow+\infty$ as $a\downarrow0$.
\item $\sqrt a\,R_*(-a)\longrightarrow+\infty$ as $a\downarrow0$.
\item $\sqrt a\,d_*(-a)\longrightarrow+\infty$ as $a\downarrow0$.
\end{enumerate}
The equivalence proves none of these equivalent assertions.
When a multiple endpoint exists, the upper witnesses have partition
height bounded by a fixed number depending on its first occurrence;
only their width and corresponding unitary rank must grow.
NN-A027 supplies sharper exact constants in its dominant-double case.
\end{theorem}

\begin{proof}
Suppose a multiple positive endpoint exists and choose the smallest
one, $x_0$, of multiplicity $m\ge2$. Such a first one exists because
the endpoint divisor is discrete and bounded away from zero. All
$p$ positive endpoint zeros below it are simple, so their inverse-square
parameters are a strictly separated positive prefix for sufficiently
small negative time. The local Hermite expansion from NN-A006 gives
\[
             z_j(-a)=x_0+i\sqrt{2a}\,h_{m,j}+O(a),
                             \qquad 1\le j\le m.
\]
These local roots are distinct for $a>0$ small. Their inverse squares
form a conjugation-stable cluster. Take its positive geometric mean
$b(a)=(\prod_j z_j(-a)^{-2})^{1/m}$ and put $\varepsilon=\sqrt a$.
Branches of the logarithm near the positive endpoint give
\[
 z_j(-a)^{-2}=b(a)e^{\varepsilon a_j(\varepsilon)},\qquad
 a_j(\varepsilon)=-\frac{2i\sqrt2}{x_0}h_{m,j}+O(\varepsilon),
                         \qquad\sum_j a_j(\varepsilon)=0.
\]
The limiting Hermite diagonal is nonzero, traceless, sign-symmetric
and has distinct entries. The uniform zero-count and tail argument
in NN-A027 gives the required smaller-modulus infinite tail and its
locally uniform positive endpoint product. Therefore
Lemma~\ref{lem:analytic-cluster-Schur-witness} supplies a negative
minor of order at most $C/\sqrt a+1$ and window at most
$C/\sqrt a+p+m$. This proves the upper bounds for $m_*$ and $R_*$.

For the actual unitary character use, as in NN-A027,
\[
       \widetilde\beta_j(a)=\frac1{p_a z_j(-a)^2-1},
                             \qquad p_a=F_{-a}'(0)>0.
\]
The positive endpoint ordering is unchanged. The cluster has the
same Hermite diagonal multiplied by the positive factor
$(1-1/(p_0x_0^2))^{-1}$. Its tail is uniformly $\ell^1$ and separated,
since the normalized original parameters have modulus bounded by
some fixed number less than one near $a=0$. The same lemma thus
gives a negative Toeplitz minor of order at most $C'/\sqrt a+1$
for the translated coefficient sequence. The exact finite-rank
criterion in NN-A021 turns it into failure of positive type at that
rank. The dual partition specifies the actual negative irreducible
coefficient; its entries are bounded by the fixed height $p+m$.

The positive lower bounds of inverse-square-root order are already
proved in NN-A022, including its explicit order/window bound and
its effective unitary-rank bound. Decreasing $c$ and $a_0$ absorbs
the fixed additive constants, proving
\eqref{eq:analytic-any-multiple-detection-scale}.

If every endpoint is simple, NN-A022 instead proves
$\sqrt a\,m_*(-a)\to\infty$ and
$\sqrt a\,d_*(-a)\to\infty$; the elementary inequality
$R_*\ge m_*-1$ gives the window statement. Conversely, the existence
of a multiple endpoint gives the bounded upper quantities just proved,
so none of these three limits can be infinite. A identifies the
positive endpoint zeros with the positive nontrivial zeta zeros,
including their multiplicities. This proves the four-way equivalence.
Rodgers--Tao and A ensure the least witnesses are finite at every
negative time even in the all-simple case; they do not decide which
case occurs.
\end{proof}

\paragraph{Geometric and representation meaning.}
The first nonreduced endpoint local algebra has finite length $m$.
Its universal leading heat unfolding is the Hermite cluster, whose
nonzero traceless Hermitian diagonal has the genuine compact orbit
$U(m)/U(1)^m$. This orbit is a full flag manifold with its usual
coadjoint-orbit symplectic form. The Fourier transform of its Haar
orbital measure is forced to become negative by
\eqref{eq:analytic-orbital-Gaussian}. The exact Weyl--HCIZ identity
then produces negative values of actual integral highest-weight
Schur characters. The finite positive prefix and infinite analytic
tail have been isolated with explicit uniform exponentially small
errors, so the same sign occurs in the actual heat coefficient minor.
The exterior-power and representation-ring operations have the exact
locator SGA~6, Expos\'e~V, Example~2.2(b) and Lemma~2.2.1.
No integrality or prequantization of the Hermite orbit's real symplectic
class is asserted. The result is an exact relation between reducedness
of the endpoint analytic divisor and growth of three concrete finite
certificate families; it is not a proof of reducedness, an arithmetic
purity realization, or a complexity-class separation.


\section{Exact detecting constants for every endpoint multiplicity pattern}
\label{sec:analytic-orbital-detection-constants}

\begin{definition}[The Hermite orbital positivity radius]
For $m\ge2$ let
\[
 D_m=\operatorname{diag}(h_{m,1},\ldots,h_{m,m}),
\]
where the entries are all roots of the probabilists' Hermite polynomial
$\operatorname{He}_m$. Define
\[
 \phi_m(X)=\int_{U(m)}e^{i\operatorname{Tr}(D_mUXU^*)}\,dU,
                 \qquad X=X^*,
\]
and
\begin{equation}\label{eq:analytic-orbital-radius}
 \mathfrak r_m=
   \inf\{\lambda_{\max}(X)-\lambda_{\min}(X):\phi_m(X)<0\}.
\end{equation}
Put $\mathfrak r_1=+\infty$ and $h_m=\max_j|h_{m,j}|$.
The orbital measure is the concrete Haar measure on the compact
Hermite orbit, not a hypothesized arithmetic measure.
\end{definition}

\begin{lemma}[Strictness, bounds and the first two radii]
\label{lem:analytic-orbital-radius-properties}
For $m\ge2$,
\[
                  0<\frac\pi{h_m}\le\mathfrak r_m<\infty.
\]
For each $c<\mathfrak r_m$, $\phi_m$ has a strictly positive minimum
on the compact set of diagonal matrices with entries in $[0,c]$.
Moreover
\begin{equation}\label{eq:analytic-first-orbital-radii}
                   \mathfrak r_2=\pi,\qquad
                   \mathfrak r_3=\frac{2\pi}{\sqrt3}.
\end{equation}
\end{lemma}

\begin{proof}
Finiteness follows from Lemma~\ref{lem:analytic-negative-orbital}.
The trace of $D_m$ is zero, so scalar translation of $X$ does not
change $\phi_m$. For small spectral diameter, subtracting a scalar
makes $\|X\|$ uniformly small. Since $\phi_m(0)=1$, this also proves
that the radius is positive.

We first rule out interior zero values in a nonnegative region.
On the Euclidean space of Hermitian matrices, differentiating under
the compact integral gives the exact elliptic equation
\begin{equation}\label{eq:analytic-orbital-Helmholtz}
                         \Delta\phi_m=-\|D_m\|^2\phi_m.
\end{equation}
The function is real analytic and real valued. If it were nonnegative
near a zero $X_0$, let $P$ be the first nonzero homogeneous term in
its Taylor series there. Such a term exists since $\phi_m$ is not
identically zero. Nonnegativity makes $P$ nonnegative in every real
direction. Equation~\eqref{eq:analytic-orbital-Helmholtz} forces its
lowest-degree part to satisfy $\Delta P=0$. The mean-value property
of harmonic polynomials gives average $P= P(0)=0$ on a sphere, so
nonnegative $P$ must vanish identically, a contradiction. Thus every
zero is approached by negative values. For $c<\mathfrak r_m$, a
neighborhood of every matrix of spectral diameter at most $c$ lies
inside the region with no negative values. No zero can occur there,
and compactness gives the stated strictly positive minimum.

To obtain the explicit lower bound, fix a diagonal spectrum $L$ of
diameter less than $\pi/h_m$, and shift its minimum to zero.
For small $\delta>0$, the finite alphabet
$e^{i\delta h_{m,j}}$ has conjugate pairs of angles at most
$\delta h_m$. The quadratic sector proof in NN-A020 makes it
$PF_r$ whenever $(r+1)\delta h_m<\pi$. Choose the partition
$\mu_i=\lfloor L_i/\delta\rfloor$, in descending order.
Its width satisfies this inequality for sufficiently small $\delta$.
Its Schur value is nonnegative. The exact normalized Weyl--HCIZ
identity \eqref{eq:analytic-Schur-orbital-limit} then gives
$\phi_m(L)\ge0$. Hence no negative value has diameter less than
$\pi/h_m$, proving the bound.

For $m=2$, the HCIZ determinant with Hermite eigenvalues $-1,1$
gives
\[
                \phi_2(X)=\frac{\sin(\lambda_1-\lambda_2)}
                                      {\lambda_1-\lambda_2},
\]
with its continuous value at equality. Its first sign change gives
$\mathfrak r_2=\pi$. For $m=3$, the eigenvalues are
$-\sqrt3,0,\sqrt3$. Factoring the exponential Vandermonde in the
same exact formula gives
\[
 \phi_3(X)=\prod_{1\le i<j\le3}
            \operatorname{sinc}\left(
                          \frac{\sqrt3}{2}(\lambda_i-\lambda_j)\right),
                    \qquad\operatorname{sinc}(u)=\frac{\sin u}{u}.
\]
All factors are positive when the diameter is less than
$2\pi/\sqrt3$. For eigenvalues $0,L/2,L$ with
$2\pi/\sqrt3<L<4\pi/\sqrt3$, the largest-difference factor is
negative and the other two are positive. This proves the second
formula in \eqref{eq:analytic-first-orbital-radii}.
\end{proof}

\begin{lemma}[Uniform finite-cluster positivity]
\label{lem:analytic-uniform-cluster-PF}
Let $m$ be fixed and let a conjugation-stable finite alphabet have
entries $b_\varepsilon e^{\varepsilon a_j(\varepsilon)}$, where
$b_\varepsilon>0$ and
$a_j(\varepsilon)=i\kappa h_{m,j}+O(\varepsilon)$ with fixed
$\kappa>0$. For each $c<\mathfrak r_m/\kappa$, its coefficient
sequence is $PF_{r(\varepsilon)}$ for all sufficiently small
$\varepsilon$, uniformly over integers $r(\varepsilon)\varepsilon\le c$.
The assertion includes every partition width in that range, with no
bound on the degree of the Toeplitz minor.
\end{lemma}

\begin{proof}
The positive common scale does not affect signs. For every partition
$\mu$ of at most $m$ rows and width at most $r$, the exact formula is
\[
 \frac{s_\mu(e^{\varepsilon a})}{s_\mu(1^m)}
       =\frac{\mathcal I(A_\varepsilon,
                       \operatorname{diag}(\varepsilon(\mu+\rho)))}
              {\mathcal I(A_\varepsilon,
                       \operatorname{diag}(\varepsilon\rho))},
            \qquad A_\varepsilon=\operatorname{diag}(a_j),
            \quad \rho=(m-1,\ldots,0).
\]
The second argument in the numerator has entries in
$[0,c+(m-1)\varepsilon]$. Choose $c'$ strictly between $c$ and
$\mathfrak r_m/\kappa$. The preceding lemma gives a positive
minimum for $\phi_m(\kappa X)$ over entries in $[0,c']$.
The Haar integral is uniformly continuous on these compact matrix
sets, while $A_\varepsilon=i\kappa D_m+O(\varepsilon)$.
The numerator is therefore uniformly close to a positive real value,
and the denominator tends to one. The character ratio itself is real,
by conjugation stability. It follows that every such ordinary Schur
value is strictly positive, uniformly in this full partition range.
Partitions with more than $m$ rows have value zero.
Every skew Schur value of width at most $r$ is a nonnegative integer
combination of ordinary Schur values of that width, by the exact
restriction-multiplicity argument in NN-A021. Thus all Toeplitz minors
of order at most $r$ are nonnegative. This proof retains the full
$O(\varepsilon)$ perturbation and the possible coincident limiting
scaled row lengths in the compact set.
\end{proof}

\begin{lemma}[Cluster isolation with earlier multiple endpoints]
\label{lem:analytic-coalescent-prefix-isolation}
The negative-witness conclusion of
Lemma~\ref{lem:analytic-cluster-Schur-witness} remains valid when its
finite prefix is a conjugation-stable multiset of possibly complex
parameters $\zeta_1(\varepsilon),\ldots,\zeta_p(\varepsilon)$ with
positive real limits strictly larger than $b^0$, even when some limits
coincide, provided the parameters are distinct for $\varepsilon>0$ and
\[
                 |\zeta_i-\zeta_j|\ge c\varepsilon
                         \quad(i\ne j)
\]
for a fixed $c>0$. Real prefix entries are positive for small
$\varepsilon$, so $P=\prod_i\zeta_i$ is positive. For any strictly
negative orbital value at spectral diameter $d$, the witness width
can be chosen with
\[
                   \limsup_{\varepsilon\downarrow0}
                               \varepsilon\lambda_1\le d+\delta
\]
for any prescribed $\delta>0$.
\end{lemma}

\begin{proof}
By continuity, first perturb the negative orbital spectrum to have
distinct eigenvalues while preserving strict negativity and increasing
its diameter by less than $\delta/4$. Shift that spectrum to have minimum in $(0,\delta/4)$,
and choose the bounding first-row length constant $C$ less than
$d+\delta$ while keeping a fixed positive gap above its maximum.
Choose the cluster row lengths and the prefix rows of length
$n=\lceil C/\varepsilon\rceil$ exactly as in NN-A028. The gaps
$\mu_m\ge\eta/\varepsilon$ and
$n-\mu_1\ge\eta/\varepsilon$ still hold for some fixed $\eta>0$.
Strict negativity is preserved by the scalar shift because the
limiting Hermite diagonal is traceless.

Here are the changes to the complete determinant estimates.
All $q=p+m$ retained poles now have pairwise separation at least a
constant times $\varepsilon$; the prefix-to-cluster separation remains
bounded below by a positive constant. Thus every partial-fraction
coefficient and the fixed-circle analytic remainder are
$O(\varepsilon^{-(q-1)})$, rather than
$O(\varepsilon^{-(m-1)})$. In a term with remainder columns, the same
sorted-log-modulus argument applies. The entire prefix is assigned
exponent $n$ in the maximal full weight. Put
\[
                  W=P^n b^{|\mu|};
\]
the maximal full weight is $W e^{O(\varepsilon)}$.
The positivity of $P$ makes this the same positive reference scale
as before. Replacing one retained modulus by the smaller tail modulus
still costs a fixed factor to the power $\mu_m$. Consequently the
total tail error is
\[
                     O(W\varepsilon^{-K}q_1^{\mu_m})
\]
for fixed $K$ and $q_1<1$; only the harmless power $K$ has changed.

In the Weyl numerator the leading prefix block is still exactly
$P^{n+m}\Delta(\zeta_1,\ldots,\zeta_p)$. It cancels its entire
prefix Vandermonde, regardless of how small that Vandermonde is.
The leading cross factor is
$\prod_{i,j}(1-\xi_j/\zeta_i)^{-1}$; it is real and tends to a
positive number. Every wrong assignment still exchanges a larger
prefix modulus with a smaller cluster modulus between rows separated
by at least $n-\mu_1$. Its exponential loss is therefore
$q_2^{n-\mu_1}$ with $q_2<1$. Dividing by the full Vandermonde now
costs at most $O(\varepsilon^{-q(q-1)/2})$, by the stated separation.
The wrong-prefix error is again
\[
                      O(W\varepsilon^{-K'}q_2^{n-\mu_1}).
\]
Both errors are smaller than any power of $\varepsilon$ times $W$.
The leading cluster Schur value, divided by its positive dimension
and $b^{|\mu|}$, tends to the fixed strictly negative orbital value.
Its sign is therefore unchanged. The dual-Jacobi--Trudi minor has
order $\lambda_1\le n$, including the empty-prefix case, and fits
in window $0,\ldots,\lambda_1+q-1$. Taking $\varepsilon\to0$
gives the asserted width bound.
\end{proof}

\begin{theorem}[NN-A029: exact orbital formulas for detecting order, window and rank]
\label{thm:analytic-exact-orbital-detection}
Assume A and put $p_0=F_0'(0)>0$. For each positive endpoint zero
$x$ of multiplicity $m_x$, put $\beta_x=1/(p_0x^2)$. Define
\begin{align}
 \mathcal L_F&=\min_{x:m_x\ge2}
                       \frac{x\mathfrak r_{m_x}}{2\sqrt2},
                          \label{eq:analytic-exact-LF}\\
 \mathcal L_U&=\min_{x:m_x\ge2}
                       \frac{x(1-\beta_x)\mathfrak r_{m_x}}{2\sqrt2},
                          \label{eq:analytic-exact-LU}
\end{align}
where a minimum over the empty set means $+\infty$. If there is any
multiple zero, these positive finite minima are attained. Then, in
the extended positive real line,
\begin{equation}\label{eq:analytic-exact-orbital-limits}
 \lim_{a\downarrow0}\sqrt a\,m_*(-a)
 =\lim_{a\downarrow0}\sqrt a\,R_*(-a)=\mathcal L_F,
       \qquad
 \lim_{a\downarrow0}\sqrt a\,d_*(-a)=\mathcal L_U.
\end{equation}
These formulas contain no unstated endpoint multiplicity or simplicity
hypothesis. They express the constants in terms of the actual endpoint
multiplicities and explicit finite-dimensional orbital integrals;
they do not determine those multiplicities.

A double endpoint contributes $\pi x/(2\sqrt2)$ to
\eqref{eq:analytic-exact-LF}; a triple endpoint contributes
$\pi x/\sqrt6$. The corresponding unitary contributions are
multiplied by $1-\beta_x$. Higher-multiplicity contributions are
specified exactly by \eqref{eq:analytic-orbital-radius}, without
asserting a closed formula for every $\mathfrak r_m$.
\end{theorem}

\begin{proof}
The large-radius multiplicity estimate in NN-A022 gives
$h_{m_x}/x\to0$ along positive zeros tending to infinity.
Since $\mathfrak r_m\ge\pi/h_m$, the quantities
$x\mathfrak r_{m_x}$ tend to infinity along multiple zeros escaping
to infinity. Also $1-\beta_x\to1$. Thus each nonempty minimum is
positive, finite and attained among a finite set of endpoint zeros.

Write $\varepsilon=\sqrt a$. The local parameters near an endpoint
$x$ have the expansion, after removing their positive geometric mean,
\[
 z_j(-a)^{-2}=b_x(a)e^{\varepsilon a_{x,j}(\varepsilon)},\qquad
 a_{x,j}(\varepsilon)=-\frac{2i\sqrt2}{x}h_{m_x,j}
                                             +O(\varepsilon).
\]
Changing the sign of the Hermite diagonal only permutes it. At a
simple endpoint the parameter stays positive and real. Uniform
$\ell^1$ tails, compact zero convergence and positivity of the endpoint
products were proved in NN-A027--A028.

Fix any finite $c<\mathcal L_F$. Choose a large circle, avoiding
endpoint zeros, so that the uniform high-radius angular estimate in
NN-A022 gives, for all zeros outside it and all sufficiently small $a$,
\[
                      |\arg(z_j(-a)^{-2})|/\varepsilon
                                      <\frac{\pi}{2(c+1)}.
\]
Every high conjugate quadratic factor is consequently $PF_r$ when
$r=\lfloor c/\varepsilon\rfloor$, for small $\varepsilon$; positive
linear factors are $PF_\infty$. Inside the circle there are finitely
many endpoint clusters. For each multiple cluster,
$c<x\mathfrak r_{m_x}/(2\sqrt2)$, so
Lemma~\ref{lem:analytic-uniform-cluster-PF} gives $PF_r$ uniformly
for its actual perturbed factors. Simple clusters are positive linear
factors. Multiplication and locally uniform convergence now give
$PF_r$ for the full entire coefficient sequence.

In fact every allowed minor of order at most $r$ is strictly positive.
There are infinitely many high factors with constant coefficient one
and positive first coefficient, each $PF_r$. In a Cauchy--Binet path,
move one index at a time one step to its required target without
colliding with another index; each such elementary factor minor is
its positive first coefficient. A finite number of high factors
realizes this path, and all remaining factors may contribute their
positive diagonal minors. Every other term is nonnegative. This is
the explicit strict-path argument used in NN-A020, and it applies
equally to the conjugate quadratic factors here. Thus
$m_*(-a)>r$, and $R_*(-a)\ge m_*(-a)-1$. Since $c<\mathcal L_F$
was arbitrary, this proves the lower limits, also when
$\mathcal L_F=+\infty$.

For the upper limit when a multiple endpoint exists, choose any such
$x$ and a negative orbital value of the scaled Hermite diagonal
$D=(2\sqrt2/x)D_{m_x}$ with spectral diameter less than
$x\mathfrak r_{m_x}/(2\sqrt2)+\delta$.
Its existence follows directly from the infimum defining
$\mathfrak r_{m_x}$. All endpoint parameters below $x$ form a finite
prefix. They need not stay real: earlier multiple endpoints may split
into conjugate clusters. Their moduli, however, remain uniformly
larger than that of the chosen cluster; all limits are positive real.
Within each earlier cluster the Hermite roots give separation at least
a constant times $\varepsilon$, and between distinct endpoint clusters
the separation stays bounded below. Their product is positive.
Hence Lemma~\ref{lem:analytic-coalescent-prefix-isolation} applies.
It gives negative minors of order $n$ and window at most $n+O_x(1)$
with
\[
 \limsup_{a\downarrow0}\sqrt a\,n
                    \le\frac{x\mathfrak r_{m_x}}{2\sqrt2}+2\delta.
\]
Choose a minimizing endpoint and let $\delta\downarrow0$. Together
with the lower limits this proves both limits involving $\mathcal L_F$.

For the unitary statement use the actual translated alphabet
\[
                \widetilde\beta_j(a)=\frac1{p_a z_j(-a)^2-1},
                                \qquad p_a=F_{-a}'(0).
\]
At endpoint $x$ its scaled Hermite diagonal is
\[
                    \frac{2\sqrt2}{x(1-\beta_x)}D_{m_x}.
\]
This follows from the full local expansion in NN-A027; it retains the
$O(\varepsilon)$ perturbation after scaling. The positive endpoint
ordering, the finite coalescent-prefix separation and the uniform
$\ell^1$ tail remain valid. The high-radius angular tail remains
uniformly negligible after division by $\varepsilon$, since
$|\beta_j(a)|\le C/|z_j(-a)|^2$ and
$|\arg(\beta/(1-\beta))|\le|\arg\beta|/(1-|\beta|)$.
The same lower and upper arguments therefore give the constant
\eqref{eq:analytic-exact-LU} for the least negative-minor order of
the translated coefficients. NN-A021 identifies that order exactly
with the least failing unitary rank. Its positive constant
$E_{-a}(-1)$ does not change any minor sign. This proves the remaining
limit. The double and triple values follow from
\eqref{eq:analytic-first-orbital-radii}.
\end{proof}

\paragraph{Geometric and representation meaning.}
For each endpoint local Artin length $m$, the leading heat unfolding
specifies the actual full-flag coadjoint orbit of the Hermite diagonal.
Its Fourier transform has the exact finite-dimensional positivity
radius \eqref{eq:analytic-orbital-radius}. The normalized Weyl--HCIZ
identity controls every integral highest-weight partition in a whole
scaled-width compact set, while its first negative orbital value gives
a negative Schur coefficient after the proved coalescent-prefix and
infinite-tail isolation. These are actual exterior-power tests, with
the representation-ring formalism in SGA~6, Expos\'e~V, Example~2.2(b)
and Lemma~2.2.1. The minima in
\eqref{eq:analytic-exact-LF}--\eqref{eq:analytic-exact-LU} are the exact
asymptotic boundaries of the corresponding analytic positive cones.
The primary HCIZ normalization is Zuber2018 equations (6), (11)--(12),
pp.~311--313, cited in NN-A028. No arithmetic Frobenius object,
integrality of the Hermite symplectic form, or complexity-class
separation is inferred.


\section{A finite Wronskian formula for every Hermite orbital radius}
\label{sec:analytic-orbital-wronskians}

The constants in NN-A029 admit a one-variable description. The
critical-length theorem used below is a classical theorem of extended
Chebyshev spaces; the identification with the Hermite orbital radius
and the resulting detecting constants are the deductions here.

For $m\ge2$ put
\begin{equation}\label{eq:analytic-hermite-ode}
 P_m(z)=i^m\operatorname{He}_m(z/i)
 =\sum_{j=0}^{\lfloor m/2\rfloor}
       \frac{m!}{2^j j!(m-2j)!}\,z^{m-2j}.
\end{equation}
Its distinct roots are $\alpha_j=ih_{m,j}$. Define the real entire
function
\begin{equation}\label{eq:analytic-fundamental-solution}
 S_m(x)=\sum_{j=1}^m\frac{e^{\alpha_jx}}{P_m'(\alpha_j)}
\end{equation}
and the Hankel Wronskians
\begin{equation}\label{eq:analytic-finite-wronskians}
 W_{m,k}(x)=\det\bigl[S_m^{(p+q)}(x)\bigr]_{p,q=0}^k,
 \qquad 0\le k\le\left\lfloor\frac{m-2}{2}\right\rfloor.
\end{equation}
The first positive zero of a function with no positive zero is, by
convention, $+\infty$.

\begin{theorem}[NN-A030: exact Wronskian reduction and computability]
\label{thm:analytic-orbital-wronskians}
The orbital radius of NN-A029 is
\begin{equation}\label{eq:analytic-wronskian-radius}
 \boxed{\displaystyle
 \mathfrak r_m=
 \min_{0\le k\le\lfloor(m-2)/2\rfloor}
        \inf\{x>0:W_{m,k}(x)=0\}.}
\end{equation}
Equivalently, it is the critical interval length for Hermite
interpolation in the real solution space of $P_m(D)f=0$.
In particular,
\begin{equation}\label{eq:analytic-radius-elementary-bounds}
       \frac\pi{h_m}\le\mathfrak r_m
                         \le\frac{(m-1)\pi}{h_m}.
\end{equation}
There is an algorithm, uniform in $m$, which returns rational bounds
$c<\mathfrak r_m<C$ of any prescribed positive rational width.
This is a computability assertion, without a polynomial-time or
uniform practical running-time bound. Consequently every local
constant in NN-A029 is specified by finitely many explicit
one-variable entire functions. Determining the actual endpoint
multiplicities remains a separate problem.
\end{theorem}

\begin{proof}
\emph{The differential equation and its canonical basis.}
The coefficient formula in \eqref{eq:analytic-hermite-ode} follows
directly from the defining formula for the probabilists' Hermite
polynomial. Thus $P_m$ is a real monic polynomial of parity $m$.
Complex conjugation permutes the terms in
\eqref{eq:analytic-fundamental-solution}, so $S_m$ is real on the
real axis. The elementary Lagrange interpolation identities give
\[
 \sum_j\frac{\alpha_j^r}{P_m'(\alpha_j)}=0
       \quad(0\le r<m-1),\qquad
 \sum_j\frac{\alpha_j^{m-1}}{P_m'(\alpha_j)}=1.
\]
Hence $S_m$ is the unique solution of $P_m(D)S_m=0$ with
\[
 S_m^{(r)}(0)=0\quad(0\le r<m-1),\qquad
 S_m^{(m-1)}(0)=1.
\]
Let $E_m$ be the real solution space. It has dimension $m$, and is
closed under translations and reflections. The functions
\[
 u_j=S_m^{(m-1-j)},\qquad 0\le j\le m-1,
\]
form a canonical basis: $u_j^{(r)}(0)=0$ for $r<j$ and
$u_j^{(j)}(0)=1$. Their complex span equals the span of
$e^{\alpha_1x},\ldots,e^{\alpha_mx}$.

\emph{Orbital determinants and arbitrary Hermite nodes.}
Write $Q=m(m-1)/2$ and $C_m=\prod_{q=1}^{m-1}q!$. The normalized
HCIZ identity, with both Vandermondes in the same order, says
\begin{equation}\label{eq:analytic-confluent-hciz}
 \phi_m(\operatorname{diag}(x_1,\ldots,x_m))
 =C_m\frac{\det[e^{\alpha_jx_k}]_{j,k=1}^m}
               {\Delta(\alpha)\Delta(x)},
 \qquad \Delta(\alpha)=i^Q\Delta(h).
\end{equation}
The identity extends continuously to repeated $x_k$. For distinct
nodes $t_1<\cdots<t_s$ with multiplicities $\nu_1+\cdots+\nu_s=m$,
the confluent numerator has columns
\[
 \left(\frac{\alpha_j^r e^{\alpha_jt_\ell}}{r!}\right)_{j=1}^m,
 \qquad 1\le\ell\le s,\quad0\le r<\nu_\ell,
\]
and the remaining denominator is the nonzero constant
$\Delta(\alpha)\prod_{\ell<v}(t_v-t_\ell)^{\nu_\ell\nu_v}$.
Thus the orbital value vanishes exactly when this Hermite
evaluation matrix is singular. Changing from the complex
exponential basis to a real basis of $E_m$ shows that singularity
is equivalent to a nonzero real $f\in E_m$ with zeros of the
prescribed multiplicities. In particular, this includes all
nonreduced node configurations, rather than just distinct nodes.

An $m$-dimensional real space is an extended Chebyshev space on an
interval if every nonzero member has at most $m-1$ zeros there,
counting multiplicities. By NN-A029, $\phi_m$ is strictly positive
for every spectrum of diameter less than $\mathfrak r_m$.
Consequently $E_m$ is an extended Chebyshev space on every
$[0,c]$ with $c<\mathfrak r_m$.

Conversely, if $c>\mathfrak r_m$, choose a negative orbital value
with spectral diameter less than $c$. Subtracting a scalar matrix
does not change it, since $\operatorname{Tr}D_m=0$. Along the
radial path from the zero matrix to this spectrum, the orbital
value starts at $1$ and ends negative. It vanishes at an
intermediate spectrum lying in $[0,c]$. The confluent determinant
criterion therefore makes $E_m$ fail the extended Chebyshev
property there. At the endpoint $c=\mathfrak r_m$, take negative
spectra with diameters decreasing to $\mathfrak r_m$, normalize
their minimum to zero, and extract a convergent subsequence. The
limit has orbital value zero: a negative limit could be radially
contracted to contradict the infimum. The endpoint also fails.
This proves that the critical interval length is exactly
$\mathfrak r_m$.

\emph{The finite Wronskian theorem and its precise indices.}
Carnicer--Mainar--Pe\~na, \emph{Spherical Bessel functions and
critical lengths}, RACSAM \textbf{117} (2023), article~29,
section~3, Proposition~3, p.~6, states that for a
reflection-invariant finite-dimensional space with a canonical
basis $(u_0,\ldots,u_n)$ and finite critical length, that length is
the least positive zero among
\[
                 \det W(u_j,\ldots,u_n),\qquad j>n/2.
\]
All hypotheses hold here; finiteness was proved in NN-A029 and
also follows from the elementary upper bound below. Put $n=m-1$
and $k=m-1-j$. Reversing these Wronskian columns changes only a
constant sign and gives $W_{m,k}$. The condition $j>(m-1)/2$
becomes $0\le k\le\lfloor(m-2)/2\rfloor$. This proves
\eqref{eq:analytic-wronskian-radius}. There are exactly
$\lfloor m/2\rfloor$ functions to test.

The lower bound in \eqref{eq:analytic-radius-elementary-bounds}
is NN-A029. For the upper bound, the nonzero function
$\sin(h_mx)$ belongs to $E_m$ and has the $m$ distinct zeros
\[
                0,\ \pi/h_m,\ \ldots,\ (m-1)\pi/h_m.
\]
The extended Chebyshev property therefore fails on their span.

\emph{Effective two-sided approximation, including tangential
Wronskian zeros.}
All roots of $\operatorname{He}_m$ are distinct real algebraic
numbers; isolating them and evaluating the finite exponential
sums with rigorous errors are effective. At zero,
\begin{equation}\label{eq:analytic-wronskian-leading}
 W_{m,k}(x)=(-1)^{k(k+1)/2}
 \left(\prod_{j=0}^k\frac{j!}{(m-1-j)!}\right)
 x^{(k+1)(m-1-k)}(1+O(x)).
\end{equation}
This follows by inserting the leading monomials
$S_m^{(j)}(x)=x^{m-1-j}/(m-1-j)!+O(x^{m-j})$ and evaluating
their Vandermonde Wronskian. Divide by the displayed power of $x$
and its nonzero signed leading coefficient. The resulting entire
functions have value $1$ at zero, computable Taylor coefficients
and effective uniform approximation bounds on every bounded
interval. Here are explicit bounds. Set $q=k+1$,
$d=q(m-1-k)$, and let $c_{m,k}$ denote the signed nonzero leading
coefficient in \eqref{eq:analytic-wronskian-leading}.
Cauchy--Binet gives the finite exponential sum
\[
 W_{m,k}(x)=\sum_{|J|=q} A_J e^{b_Jx},\qquad
 A_J=\frac{\Delta(\alpha_J)^2}{\prod_{j\in J}P_m'(\alpha_j)},
 \quad b_J=\sum_{j\in J}\alpha_j.
\]
Put $B=\sum_J|A_J|$ and $M=qh_m$. All $b_J$ are purely imaginary
and $|b_J|\le M$. For $T(x)=W_{m,k}(x)/(c_{m,k}x^d)$, its degree-$N$
Taylor polynomial $T_N$ therefore satisfies, on $|x|\le R$,
\[
 |T-T_N|\le\frac{BM^d}{|c_{m,k}|}
      \frac{(MR)^{N+1}e^{MR}}{(N+d+1)!},\qquad
 \sup_{|x|\le R}|T'(x)|
             \le\frac{BM^{d+1}}{|c_{m,k}|}e^{MR}.
\]
Indeed the coefficients of $T$ are
$W_{m,k}^{(n+d)}(0)/(c_{m,k}(n+d)!)$ and their absolute values are at
most $BM^{n+d}/(|c_{m,k}|(n+d)!)$. Summing this majorant gives both
bounds. The algebraic numbers $A_J,b_J$ and the elementary
exponentials allow outward rational bounds for every displayed
quantity.

For any rational $c<\mathfrak r_m$, all these finitely many
normalized Wronskians are strictly positive on $[0,c]$. Their
positive minima can therefore be certified by successive
uniform approximations and a finite mesh with rigorous
derivative bounds. This provides a semidecision procedure for
valid rational lower bounds. A certificate cannot pass at or
beyond the first zero.

Independently enumerate distinct rational spectra and certify
negative orbital values using \eqref{eq:analytic-confluent-hciz}.
Every rational $C>\mathfrak r_m$ admits such a negative spectrum
of diameter less than $C$: a negative spectrum exists by the
definition of the infimum, and negativity persists under a
sufficiently small perturbation to distinct rational entries.
This supplies certified upper bounds. Dovetail the lower and
upper searches until their width is smaller than the requested
error. Both searches terminate for bounds chosen sufficiently
close to the radius. No decision about whether an individual
Wronskian root is multiple, and no unjustified sign-change test
at that root, is used. This proves the algorithmic claim.
\end{proof}

\paragraph{The first examples.}
Here $S_2(x)=\sin x$ and
$S_3(x)=(1-\cos(\sqrt3x))/3$, so the single-Wronskian cases give
$\mathfrak r_2=\pi$ and $\mathfrak r_3=2\pi/\sqrt3$ again.
For $m=4$, put
\[
 a=\sqrt{3-\sqrt6},\qquad b=\sqrt{3+\sqrt6},\qquad
 S_4(x)=\frac{\sin(ax)/a-\sin(bx)/b}{b^2-a^2}.
\]
Then $\mathfrak r_4$ is exactly the smaller of the first positive
zeros of $S_4$ and $S_4S_4''-(S_4')^2$. This is a finite scalar
root problem, without a numerical value being asserted here.

\paragraph{A checked source correction.}
Beccari--Casciola--Mazure, \emph{Critical length: an alternative
approach}, arXiv:1904.09253v1, p.~12, Theorem~4.3 and
Corollary~4.4, equations~(19)--(20), print $W\ne0$ inside the
infima. The original page image was checked: this is a printed
sign error, not an extraction artifact. Those conditions must
read $W=0$, as the preceding sentence's first-zero description
and the correct Proposition~3 above specify. For $P_2(D)=D^2+1$
the printed nonzero condition gives infimum $0$ from $S_2=\sin x$,
whereas the critical length is $\pi$. No claim is made here
about whether later versions corrected that display.

\paragraph{Geometric translation.}
The algebra $\mathbb Q[z]/(P_m)$ is finite etale because $P_m$
has distinct roots; the finite-free discriminant criterion is
SGA~1, Expos\'e~I, section~4, Proposition~10. A companion matrix
$B_m$ of $P_m$ gives the actual algebraic connection $d-B_m\,dx$ on a
rank-$m$ trivial bundle on $\mathbb A^1$. Its analytic horizontal
sections are the solution space above. Evaluation on the finite
nonreduced scheme
\[
              \coprod_{\ell=1}^s
                  \operatorname{Spec}\mathbb C[x]/(x-t_\ell)^{\nu_\ell}
\]
is exactly the Hermite jet map, and its determinant vanishes
precisely when interpolation fails. A zero of $W_{m,k}(x)$ gives
a nonzero solution vanishing to order at least $m-1-k$ at $0$
and at least $k+1$ at $x$: use the span of
$u_{m-1-k},\ldots,u_{m-1}$. Thus the Wronskian formula reduces
this real interval boundary to finitely many two-point jet tests.
It does not turn the exponential connection into an arithmetic
Frobenius object, a pure motive, or a polynomial-time algorithm.


\section{Exact interfaces: arithmetic tests, fluid profiles, and source corrections}

\begin{proposition}[NN-A011: the specified fluid and branch channels preserve the A constraint]
Retain the exact source maps in workbench24 and workbench29i.
Under A, every finite matrix $[W(h_i,h_j)]$ of admissible scalar
Weil tests $h_i\in C_c^\infty(\mathbb R)$ is positive semidefinite,
including every polynomial channel $h_i=L_{P_i}U_i$ constructed in
workbench24. The branch-trace form on its retained two-sheet module
is a different form and can be negative without contradicting this.

In the released-witness arithmetic realization of workbench29i, take
$H=\Xi$ and
\[
 \Psi_{\nu,j}(t,q,\zeta)
 =\Xi(S_j(X_{\nu,t}(q))+\zeta).
\]
For every preterminal time, and almost every terminal label in the
proved terminal construction, all its zeros in the variable $\zeta$
are real under A, with unchanged multiplicity. This holds together
with the source's proved divergence of the specified fixed-label
partial-derivative norm. In this particular map the derivative
singularity cannot yield an off-critical zeta zero.
\end{proposition}
\begin{proof}
For A's real ordinates the exact Weil form is
$W(f,g)=\sum_\rho m_\rho\overline{\widehat f(\gamma_\rho)}
\widehat g(\gamma_\rho)$, absolutely convergent on the specified
test class. Hence
$c^*[W(h_i,h_j)]c=\sum_\rho m_\rho|
\widehat{\sum c_ih_i}(\gamma_\rho)|^2\ge0$.
The polynomial channel preserves that test class and compact support,
as proved by its explicit differential-operator formula in workbench24.
By contrast its branch trace has fibre matrix
$\operatorname{diag}(2,-4\lambda)$ at a real $\lambda>0$; it is
not the componentwise Weil form or its scalar pullback.

The shifts $a=S_j(X_{\nu,t}(q))$ are real at all the stated parameters.
The translation $\zeta\mapsto a+\zeta$ has inverse $z\mapsto z-a$,
and transforms a germ $(z-z_0)^mA(z)$ to
$(\zeta-(z_0-a))^mA(a+\zeta)$. It preserves multiplicity and
imaginary part exactly. A makes the zeros $z_0$ of $\Xi$ real.
The completed arithmetic argument is $s=1/2+i(a+\zeta)$,
so $\Re s=1/2-\Im\zeta=1/2$ at those zeros.
No estimate on $\partial_ta$ enters this identity, so its possible
divergence does not affect the conclusion.
\end{proof}

This closes the particular proposed counterexample route after A; it does
not assert a theorem about all hypothetical fluid-to-arithmetic maps.
The underlying categories are a concrete analytic test-function module,
its Hermitian Gram form, and translated analytic divisors. The branch
trace is an algebra trace in a quadratic algebra. The equality of those
forms was never available. Neither a PDE singularity nor this translation
creates a finite-dimensional pure motive, a negative Weil matrix, or a
complexity-class separation.

The Bost--Connes formulas in source22 also retain their domains. Its
physical Gibbs partition function is $\zeta(\beta)>0$ for real
$\beta>1$. A locates the nontrivial zeros of its meromorphic continuation
on $\Re s=1/2$. The source's strong-limit projection $P_1$ has continued
expectation $1/\zeta(s)$ and therefore poles at those zeros; finite
Euler-projector expectations are entire finite products and the strong
operator limit is not automatically a meromorphic limit through those
poles. A supplies zero location, not a positive Gibbs state at complex
temperature. Any ``Lee--Yang'' terminology must refer to this stated
analytic continuation and coordinate, not silently to a unit-circle
theorem for a finite statistical-mechanical partition polynomial.

\begin{proposition}[NN-A012: two exact source corrections in paper16]
The finite-gerbe path integral printed in
\texttt{thm:finite-gerbe-anomaly} satisfies
$\mathcal A_n(S^4)=n$. Consequently its later assertion
$\mathcal A(q_n)\otimes\mathcal A(-q_n)\simeq\mathbf1$
in \path{prop:anomaly-cancellation} fails for every $n>1$ with
that normalization. Complex conjugation is not the inverse of an
arbitrary non-unit-modulus partition value.

Also, the measure algebra of an arbitrary proper adelic curve is a
Boolean sigma-algebra, but need not be a complete Boolean algebra.
The completeness assertion in \texttt{thm:final-synthesis-v11}(3)
requires an additional measure-space hypothesis such as localizability.
\end{proposition}
\begin{proof}
For $S^4$ and $\Phi_n=\mathbb Z/n\mathbb Z$, one has
$|H^0|=n$, $|H^1|=1$, and $H^2=0$. The sole quadratic action is zero,
so the displayed source formula gives $n$ for both $q_n$ and $-q_n$.
Their product is $n^2$, whereas the tensor unit has value one. On this
example a compensating Euler factor $n^{-\chi(S^4)/2}=n^{-1}$
normalizes each value; a general inverse-theory assertion must specify
the full Euler normalization, not merely replace $q$ by $-q$.

For the second assertion, let $\Omega$ be uncountable and equip it
with the countable-or-cocountable sigma-algebra and counting measure.
Assign the trivial absolute value of any field to every point.
Every principal logarithm is identically zero, so this is a proper
adelic curve under the source's stated definition. The only null set
is empty. The displayed definition at the start of
\path{sec:adelic-support-defect} imposes integrability of principal
logarithms and the product formula, but no sigma-finiteness or
localizability; the counterexample is directed at that exact definition.
Thus the measure algebra
is exactly that sigma-algebra.
Choose an uncountable subset $A$ with uncountable complement. The
family of measurable singleton sets $\{\{a\}:a\in A\}$ has no
least measurable upper bound: any upper bound is cocountable and
contains some $x\notin A$, and deleting $x$ gives a smaller measurable
upper bound. Thus arbitrary suprema need not exist, whereas countable
suprema always do.
\end{proof}

These are finite or explicit set-theoretic counterexamples, not new
consequences of E or A. They are retained as source repairs because
otherwise the proposed geometry would rely on false normalization or
completeness assertions. They supply no computational lower bound.
